% =====================================================================================================
% Chương 3 — bản viết lại, phiên bản 4 (2026-10-12). Dùng như % =====================================================================================================
% Chương 3 — bản viết lại, phiên bản 4 (2026-10-12). Dùng như % =====================================================================================================
% Chương 3 — bản viết lại, phiên bản 4 (2026-10-12). Dùng như % =====================================================================================================
% Chương 3 — bản viết lại, phiên bản 4 (2026-10-12). Dùng như \input{noi-dung/07-chuong3} của main.tex.
% Bản 1 (07-chuong3.tex) và bản trung gian 2 (07-chuong3-v2.tex) được giữ nguyên. Lý do từng thay đổi:
% reviews/CH3-REVISION-LOG.md. Ghi chú v3: 07-chuong3-v3-NOTES.md; thay đổi v4: 07-chuong3-v4-CHANGES-v3.md.
% Mỗi dòng "% EV:" trỏ tới EVIDENCE-REGISTER.md (A01…K52), EVIDENCE-REGISTER-ADDENDUM-1.md (N-01…N-30,
% L1a-01…09), EVIDENCE-REGISTER-ADDENDUM-2-CH3.md (X3-01…X3-41) hoặc các nguồn R26/TESTRUN ghi tại chỗ.
% =====================================================================================================

% Môi trường "Nhận xét" chưa có trong main.tex; định nghĩa có điều kiện ở đây (nên chuyển lên phần mở đầu).
\ifcsname nhanxet\endcsname\else
  \theoremstyle{definition}
  \newtheorem{nhanxet}{Nhận xét}[chapter]
\fi

\chapter{Phương pháp: luật C, lời hứa theo dự báo và sổ theo từng khách}
\label{chuong:phuongphap}

Khi đặt xe đi chung, hành khách được báo giờ xe sẽ trả mình tới nơi; khóa luận gọi đó là lời hứa giờ trả, gọi
tắt là lời hứa. Giờ trả có thể bị dời vì bộ điều phối ghé đón thêm khách, hoặc vì đường đông hơn lúc hứa hay xe
đi nhanh, chậm hơn dự kiến. Luật C, nội dung chính của chương, chỉ giới hạn loại thứ nhất. Ở mỗi lượt quyết định,
hiệu giữa giờ trả khi
chọn phương án mới và khi xe giữ nguyên kế hoạch là phần dời do điều phối; luật C loại phương án làm tổng độ lớn
phần này qua cả chuyến vượt một trần. Nhịp tắc đường quen thuộc trong ngày được cộng trước vào lời hứa đầu nhờ dự
báo; phần còn lại được ghi trong sổ theo từng khách (per-rider ledger). Các mục đi theo thứ tự đó, sau phần phát
biểu bài toán.
% EV: A01; A02; A06; F01; X3-12 (FFA:134-138, ba phần của sổ)

% -----------------------------------------------------------------------------------------------------
\section{Phát biểu bài toán và tổng quan hệ thống}
\label{muc:baitoan}

\subsection{Thuật ngữ và vòng đời của một yêu cầu}
\label{muc:thuatngu}

Bảng~\ref{bang:thuatngu} giải thích các thuật ngữ; tên tiếng Anh chỉ ghi trong bảng. Mọi thực nghiệm chạy trên
bộ mô phỏng FleetPy~\cite{engelhardt2026} với dữ liệu chuẩn Manhattan (mục~\ref{muc:kientruc}).
% EV: B01; J18 (FLEET:1-12, 474)

\begin{table}[htbp]
\centering
\caption[Thuật ngữ chính của phương pháp]{Thuật ngữ dùng trong chương}
\label{bang:thuatngu}
\small
\begin{tabular}{@{}p{4.3cm}p{10.5cm}@{}}
\toprule
\textbf{Thuật ngữ} & \textbf{Nghĩa trong khóa luận}\\
\midrule
Lượt quyết định (decision epoch) & Một lần bộ điều phối nhận một lô sự kiện và trả quyết định cho cả đội xe.\\
Phương án (candidate plan) & Một kế hoạch (dãy điểm đón, điểm trả theo thứ tự) có thể chọn cho một xe.\\
Phương án giữ nguyên kế hoạch (keep-route option) & Xe đi tiếp kế hoạch cũ, không nhận thêm khách.\\
Bộ sinh chỉ chèn thêm (insertion-only generator) & Chèn điểm đón, điểm trả của yêu cầu mới vào kế hoạch của
một xe, không bỏ, không đảo điểm cũ.\\
Bộ chọn (selector); bộ giải (solver) & Chọn mỗi xe một phương án theo thứ tự tiêu chí (lexicographic order):
chỉ khi hòa ở tiêu chí trước mới xét tiêu chí sau; bộ giải của OR-Tools~\cite{ortools} giải bài toán chọn.\\
Lời hứa; lời hứa đầu $p_0$ & Giờ đón, giờ trả báo cho khách; lời hứa đầu công bố khi khách xác nhận đặt xe;
$p_0$ là giờ trả của nó, chưa cộng dự báo.\\
Công bố; bản ghi sửa & Ghi lời hứa mới vào sổ và gửi trong bản tin quyết định; mỗi lần sau lời hứa đầu là một
bản ghi sửa.\\
Khóa (lock); chiều cam kết & Khóa: phần lời hứa không được đổi sau khi xác nhận. Chiều cam kết: thuộc tính của
lời hứa được theo dõi riêng, như giờ trả, việc đổi xe.\\
Trần $B$ (cap); phần đã dùng $L_r$ & Tổng phần dời do điều phối tối đa của một khách cả chuyến; phần khách $r$
đã dùng.\\
Quyết định bị ép (forced decision); luật phục hồi (recovery rule) & Lượt quyết định có ít nhất một xe bị ép:
luật lời hứa đã loại cả phương án giữ nguyên kế hoạch của xe, và xe chọn phương án được giữ nhờ luật phục hồi,
có ít nhất một khách được miễn trừ; luật phục hồi quy định phương án nào được giữ khi đó.\\
Hệ số tắc đường & Số nhân chung cho cả thành phố áp lên mọi thời gian đi, đổi theo khung 15 phút; hệ số 1,2 làm
thời gian đi dài gấp 1,2 lần hệ số 1.\\
Kịch bản & Một ngày, một mẫu nhu cầu, một khung 2 giờ của dữ liệu Manhattan: 108 yêu cầu, 8 xe 4 chỗ.\\
\bottomrule
\end{tabular}
\end{table}
% EV (lượt, phương án, bộ sinh, bộ chọn): X3-01; SPEC:31-34; X3-03; X3-07; X3-22; A14
% EV (lời hứa, công bố, khóa, chiều, trần): A13; X3-25 (fleet_control.py:882-893); A07; A04; A06; K43
% EV (bị ép, phục hồi, hệ số, kịch bản): X3-30; X3-09; B06; STRETCH:3-4; B02; B03; GLOSSARY-VI-v2.md §E (bị ép); C4A-10

Một yêu cầu mới chờ được xét. Khi phương án được chọn của một xe chứa nó, yêu cầu được nhận và khách được mời,
nhưng chưa có lời hứa. Khi khách xác nhận đặt xe, lời hứa đầu được công bố và $L_r = 0$; khách chờ đón, lên xe
rồi được trả. Yêu cầu chưa được nhận được xét lại ở lượt sau nếu còn phương án hợp lệ hay chưa được xét vì giới
hạn số yêu cầu mới, và bị từ chối khi mọi phương án chứa nó đều bị loại. Khóa và luật lời hứa áp từ lúc xác nhận
tới lúc khách xuống xe.
% EV: X3-20 (VALID:272-336; RollingCostPolicy.cs:173-226; fleet_control.py:878-924); A13; X3-32

\subsection{Bài toán ở mỗi lượt quyết định}
\label{muc:baitoan-phatbieu}

Bộ chạy tiến theo đồng hồ sự kiện: yêu cầu đến, xe tới mốc vật lý và mốc cập nhật tắc đường đều có thể gọi một
lượt quyết định. Xác nhận đặt xe và cập nhật trạng thái có thể tạo thêm lượt ở cùng thời điểm mô phỏng; không có
chu kỳ điều phối cố định. Bộ chuyển đổi gom các sự kiện từ lần gọi trước (yêu cầu mới, xác nhận đặt xe, tiến độ xe,
đón, trả, thời gian đi mới) thành một lô.
% EV: R26-CH3-OPEN-QUESTIONS-CODEX.md §Q1; actual_fleetpy_medium_preflight.py:528-543, 569-637; fleet_control.py:395-451; preflight_world_v1_1.py:40-58; preflight_world_v1_2.py:18-20

Khung nhu cầu dài 2 giờ, nhưng lần chạy kết thúc khi hết yêu cầu và hành trình đang hoạt động. Trong 34 lần chạy
luật C trần 120 giây của tập kiểm tra chính, thời lượng mô phỏng từ 118,02 đến 154,93 phút.
% EV: R26-CH3-OPEN-QUESTIONS-CODEX.md §Q2; round26-trace-rider.txt, Q2 C-120 MAIN-PANEL SIMULATED LENGTHS; actual_fleetpy_medium_preflight.py:580-581, 639-683

Bài toán ở mỗi lượt như sau.
\begin{itemize}[leftmargin=1.2cm,itemsep=2pt,topsep=3pt]
\item \textbf{Đầu vào:} trạng thái xe (vị trí, kế hoạch, sức chứa, khách trên xe, khách chờ đón); sổ; yêu cầu
đang chờ với điểm đón, điểm trả, cửa sổ đón (từ giờ đón sớm nhất tới giờ đón muộn nhất), thời gian ngồi xe tối
đa; thời gian đi lúc này, coi là giữ nguyên cho phần đường còn lại.
\item \textbf{Quyết định:} nhận, để lại hay từ chối từng yêu cầu đang chờ; mỗi xe một phương án thỏa các ràng
buộc cứng (Bảng~\ref{bang:rangbuoc}).
\item \textbf{Đầu ra:} lời hứa đầu và các bản ghi sửa (mục~\ref{muc:so}).
\item \textbf{Mục tiêu:} thứ tự tiêu chí của Bảng~\ref{bang:rangbuoc}; các chỉ số đánh giá ở
Chương~\ref{chuong:thucnghiem}.
\end{itemize}
% EV: X3-21 (fleet_control.py:395-429); X3-01; X3-02; A05; X3-20; B03; A06; A07; A14

Bộ chọn giải một quy hoạch nguyên, với $y_{k\pi} \in \{0, 1\}$ bằng 1 khi xe $k$ chọn phương án $\pi$ trong
tập $P_k$ của nó:
\begin{equation}
\sum_{\pi \in P_k} y_{k\pi} = 1 \;\;\text{với mọi xe } k, \qquad
\sum_{k} \sum_{\pi \in P_k,\; u \in \pi} y_{k\pi} \le 1 \;\;\text{với mọi yêu cầu mới } u ,
\label{eq:chon}
\end{equation}
tức mỗi yêu cầu mới thuộc phương án được chọn của nhiều nhất một xe. Các tiêu chí được giải lần lượt, mỗi lần giữ
các tiêu chí trước ở giá trị tốt nhất đã tìm. Bước chọn này theo lối gán xe--chuyến động của Alonso-Mora và cộng
sự~\cite{alonsomora2017} (mục~\ref{muc:c2-baitoan}).
% EV: X3-22 (OrTools:3, 96-125, 180-189, 343-365); X3-34 (HZ:341-345); J17

\begin{table}[htbp]
\centering
\caption[Ràng buộc cứng và thứ tự tiêu chí của bộ chọn]{Ràng buộc cứng của một phương án (nửa trên) và thứ tự tiêu chí của bộ chọn (nửa dưới): so tiêu chí 1
trước, chỉ khi hòa mới xét tiêu chí 2, và cứ thế. Giá trị là cấu hình của các tập thử, không phải chuẩn ngành.}
\label{bang:rangbuoc}
\small
\begin{tabular}{@{}p{2.4cm}p{7.6cm}p{4.8cm}@{}}
\toprule
\textbf{Nhóm} & \textbf{Nội dung} & \textbf{Vai trò}\\
\midrule
Ràng buộc dịch vụ & Đón trong vòng 10 phút kể từ lúc gọi; ngồi xe không quá 1,5 lần thời gian đi thẳng; mỗi
xe tối đa 4 khách & Phương án đổi kế hoạch phải thỏa chặt; phương án giữ nguyên được nới phần vượt thời gian
mà kế hoạch cũ đã có\\
% EV: B03; X3-23 (GEN:161-246; ServiceQualityAllowance.cs:61-94)
Khóa & Không đổi xe, điểm đón, giờ đón; điểm trả không đổi nhờ trần 0. Giờ đón mới phải bằng giờ đón khi giữ
nguyên kế hoạch & So với kế hoạch giữ nguyên, nên riêng tắc đường không làm vi phạm khóa\\
% EV: A07; A08; A10; X3-06 (LOCK:42-60, 103-161, 204-207; CFG C-30.json:66-67, 94-98)
Luật lời hứa & Luật C: $L_r + |p_r - e_r| \le B$ cho mọi khách có lời hứa của xe, hoặc một luật ở
Bảng~\ref{bang:luat} & Lọc trước bộ giải\\
% EV: A06; A09; A11; HVA:62-66
\midrule
Tiêu chí 1 & Ít phương án vi phạm được giữ lại nhất (chỉ khi có xe bị ép) & Chỉ dùng chúng khi không còn cách
khác\\
Tiêu chí 2 & Nhận nhiều yêu cầu mới nhất & Ưu tiên phục vụ\\
Tiêu chí 3 & Mức dùng trần lớn nhất trên đội xe càng nhỏ càng tốt; nếu hòa, tổng phần dời do điều phối nhỏ nhất,
so lần lượt từng chiều cam kết (khi có chiều mang trần) & Ưu tiên ổn định\\
Tiêu chí 4 & Tổng chi phí kế hoạch nhỏ nhất, tức thời lượng còn lại của các tuyến dự kiến & Theo mô hình, không
phải chi phí vận hành thực\\
Tiêu chí 5 & Thứ tự mã phương án & Phá hòa, để chạy lại cho cùng kết quả\\
% EV: A14; X3-07 (SEL:271-330, 365-385, 510-522; HVA:537-636); DIM:19-32; A16; I09
\bottomrule
\end{tabular}
\end{table}

\subsection{Một lượt quyết định}
\label{muc:luotquyetdinh}

Hình~\ref{hinh:vongquyetdinh} và Bảng~\ref{bang:luotquyetdinh} mô tả một lượt. Ngoài sinh phương án và chọn,
mỗi lượt có ba bước cho lời hứa: tính các mốc giờ trả của từng khách, loại phương án vi phạm khóa hay luật lời
hứa trước khi gọi bộ giải, và ghi sổ sau khi chọn. Hình và giả mã dùng một số ký hiệu được định nghĩa ở các mục
sau: với mỗi khách có lời hứa, $v$ là giờ trả đã công bố gần nhất, $e$ là giờ trả nếu xe giữ nguyên kế hoạch, $p$ là
giờ trả theo phương án đang xét, còn $x$, $c$, $z$ là ba hiệu giữa chúng (mục~\ref{muc:bamoc}); $D$ là lời hứa theo
dự báo, tính qua giờ $S$ (mục~\ref{muc:dubao}).
% EV: X3-04 (VALID:201-239, 355-411, 494-520); X3-13; X3-25

\begin{figure}[htbp]
\centering
\begin{tikzpicture}[
  node distance=6mm,
  hop/.style={draw=xamvien,fill=xamnhat,rounded corners=3pt,align=center,
              inner sep=4pt,font=\small,minimum height=12mm,text width=2.75cm},
  ngoai/.style={draw=xamvien,dashed,rounded corners=3pt,align=center,
              inner sep=4pt,font=\small,minimum height=10mm,text width=4.2cm},
  mui/.style={-{Latex[length=2.2mm]},thick}
]
\node[hop] (a) {Nhận lô sự kiện\\mới};
\node[hop,right=of a] (b) {Sinh phương án chỉ chèn thêm, kiểm ràng buộc dịch vụ};
\node[hop,right=of b] (c) {Tính $v$, $e$, $p$ cho từng khách có lời hứa};
\node[hop,right=of c] (d) {Loại phương án vi phạm khóa hay luật lời hứa};
\node[hop,below=9mm of d] (e) {Bộ giải chọn\\một phương án\\mỗi xe};
\node[hop,left=of e] (f) {Kiểm lại cả đội xe; công bố lời hứa $p$; ghi sổ};
\node[ngoai,below=8mm of f] (g) {Sau khi chạy, trên bản ghi đã lưu: lời hứa hiển thị $D = \max(p, S)$};
\draw[mui] (a)--(b); \draw[mui] (b)--(c); \draw[mui] (c)--(d);
\draw[mui] (d)--(e); \draw[mui] (e)--(f);
\draw[mui,dashed] (f.west) -| node[pos=0.3,below,font=\footnotesize]{lượt sau} (a.south);
\draw[-{Latex[length=2.2mm]},dashed] (f.south)--(g.north);
\end{tikzpicture}
\caption[Một lượt quyết định của RideBound]{Một lượt quyết định của RideBound. Bộ điều phối công bố lời hứa không dự báo $p$. Trong các kết quả
chính, lời hứa theo dự báo $D$ được tính sau khi chạy, từ bản ghi đã lưu (ô nét đứt), và không quay lại làm đầu
vào của quyết định nào (mục~\ref{muc:dubao}).}
\label{hinh:vongquyetdinh}
\end{figure}
% EV: X3-03; X3-04; X3-22; X3-25; F03; X3-11 (DISP:1-8); FR:10-11

\begin{table}[htbp]
\centering
\caption[Giả mã một lượt quyết định]{Giả mã một lượt quyết định; dấu $\gets$ nghĩa là gán giá trị bên phải cho biến bên trái}
\label{bang:luotquyetdinh}
\small
\begin{tabular}{@{}r@{\hspace{6pt}}p{10.2cm}@{\hspace{8pt}}p{4.0cm}@{}}
\toprule
\multicolumn{3}{@{}p{14.9cm}@{}}{\textbf{Đầu vào:} trạng thái đội xe; sổ; lô sự kiện mới; luật lời hứa đang
dùng. \textbf{Đầu ra:} kế hoạch mỗi xe; yêu cầu được nhận, để lại hay từ chối; lời hứa mới; bản ghi sổ}\\
\midrule
1: & \textbf{hàm LượtQuyếtĐịnh}(trạng thái, lô sự kiện) & \\
2: & \quad áp lô sự kiện; ghi nhận lần đón muộn hơn giờ đón muộn nhất, nếu có & vị trí xe, đón, trả\\
3: & \quad chọn tối đa 2 yêu cầu đang chờ có giờ đón muộn nhất sớm nhất & yêu cầu khác để lượt sau\\
4: & \quad \textbf{với mỗi} xe $k$ \textbf{thực hiện} & \\
5: & \qquad \textbf{nếu} kế hoạch đang chạy vi phạm ràng buộc cấu trúc \textbf{thì} báo lỗi, dừng lần chạy &
ví dụ vượt sức chứa\\
6: & \qquad $P_k \gets \{$phương án giữ nguyên$\}$, kiểm dịch vụ có nới phần vượt do tắc đường & luôn giữ khi
cắt bớt\\
7: & \qquad thêm vào $P_k$ các phương án chèn yêu cầu ở dòng 3, thỏa chặt ràng buộc dịch vụ; giữ tối đa 100 &
bộ sinh chỉ chèn thêm\\
8: & \qquad \textbf{với mỗi} $\pi \in P_k$ và mỗi khách có lời hứa của xe: tính $v$, $e$, $p$; \textbf{nếu}
$\pi$ vi phạm khóa hay luật lời hứa \textbf{thì} loại $\pi$ & Bảng~\ref{bang:kiemtran}\\
9: & \qquad \textbf{nếu} phương án giữ nguyên bị loại chỉ vì lỗi hạn chót hay lỗi trần \textbf{thì} & xét phục hồi\\
10: & \qquad\quad giữ nó làm phương án được miễn trừ & phục hồi giữ nguyên kế hoạch\\
11: & \qquad\quad giữ thêm phương án đã bị loại không làm vi phạm đó tệ hơn & phục hồi không tệ hơn\\
12: & \qquad \textbf{nếu} phương án giữ nguyên bị loại mà không được miễn \textbf{thì} báo lỗi, dừng lần chạy &
xe kẹt\\
13: & \quad chọn $y$ theo~\eqref{eq:chon} và thứ tự tiêu chí; kiểm lại lựa chọn trên cả đội xe &
Bảng~\ref{bang:rangbuoc}\\
14: & \quad nhận yêu cầu có trong phương án được chọn; để lại hay từ chối yêu cầu khác & \\
15: & \quad khách vừa xác nhận: công bố lời hứa đầu, $L_r \gets 0$; khách có lời hứa mà $x$, $c$ hay $z$ khác 0:
ghi bản ghi sửa, $L_r \gets L_r + |c_r|$ & ghi vi phạm cho khách được miễn\\
16: & \quad \textbf{trả về} quyết định; chờ bộ mô phỏng xác nhận đã áp dụng & \\
17: & \textbf{kết thúc hàm} & \\
\bottomrule
\end{tabular}
\end{table}
% EV (dòng 2–3): X3-01; X3-24 (WP4:8; EventReducer.cs:13-23, 260-282); X3-32 (GEN:54-96)
% EV (dòng 4–7): X3-31 (GEN:171-182); X3-23; X3-27 (CandidatePortfolioRetainer.cs:102-113); X3-03 (GEN:548-575); A24
% EV (dòng 8–12): X3-04; A06; A07; A10; X3-09 (HVA:244-390); X3-31 (POL:336-339; RunnerSession.cs:941-947)
% EV (dòng 13–16): X3-22; X3-07; POL:370-401; X3-20; X3-25 (VALID:272-336, 562-597; RunnerSession.cs:956-987); X3-30; RunnerSession.cs:448-455

Luật lời hứa là bộ lọc cứng: phương án vi phạm bị loại trước khi bộ giải chạy, không được giữ kèm một điểm trừ.
Ngoại lệ là khi xe bị ép (dòng 9--11), chẳng hạn khi riêng việc đường đông hơn đã đẩy giờ trả vượt hạn chót của
luật hạn chót (mục~\ref{muc:doichung}); khi đó tiêu chí 1 chỉ dùng phương án vi phạm nếu xe không còn cách khác.
Bộ sinh không bỏ, không đảo điểm cũ nào, nên khách đã nhận không bị chuyển sang xe khác. Phần nới ở dòng 6 giúp
xe không kẹt vì tắc đường. Mỗi lần vượt như vậy, và mỗi lần khách thực tế lên xe muộn hơn giờ đón muộn nhất, được ghi
lại thay vì làm dừng lần chạy. Vì vậy cửa sổ đón là điều kiện để kiểm phương án, không phải bảo đảm đón đúng hạn.
% EV: A06; HVA:62-66; X3-09 (HVA:255-326); X3-07 (SEL:283-295); A10; X3-03 (GEN:548-575); TEX08:129-131; X3-23 (GEN:161-246; ServiceQualityAllowance.cs:5-18, 61-94; RunnerSession.cs:989-1009); X3-24 (WP4:8; EventReducer.cs:13-23, 260-282)

% -----------------------------------------------------------------------------------------------------
\section{Ba mốc giờ trả và phần dời do điều phối}
\label{muc:bamoc}

Giờ trả có thể đổi vì bộ điều phối ghép thêm khách, hoặc ngay cả khi xe giữ kế hoạch cũ. Để tách hai loại, với
mỗi khách có lời hứa và mỗi phương án, hệ thống so ba mốc giờ trả:
\begin{itemize}[leftmargin=1.2cm,itemsep=2pt,topsep=3pt]
\item $v$: giờ trả đã công bố cho khách ở lần gần nhất;
\item $e$: giờ trả tính lại nếu xe giữ nguyên kế hoạch cũ, từ trạng thái hiện tại với thời gian đi hiện tại;
\item $p$: giờ trả nếu chọn phương án đang xét, với cùng trạng thái và thời gian đi.
\end{itemize}
Ba hiệu sau tính bằng giây (sổ lưu tới mili giây):
\begin{equation}
x = e - v, \qquad c = p - e, \qquad z = p - v = x + c .
\label{eq:bamoc}
\end{equation}
Trong đó, $x$ là phần thay đổi khi giữ nguyên kế hoạch, $c$ là phần dời do điều phối, $z$ là phần khách thấy;
dễ nhớ, $e$ là giờ nếu xe cứ đi như cũ. Mốc $e$ không phải ý của khóa luận: Hörl và Zwick~\cite{horl2022} đã tính
lại giờ trả trên tuyến đang gán theo thời gian đi hiện tại ở mỗi lượt, và khóa luận dùng lại mốc này.
% EV: A01; A02; X3-04 (VALID:355-384); X3-13 (SPEC:39-50); LEDGER:188-207; J01

Mốc $e$ là giờ do mô hình tính, không phải giờ đến thật, nên $x$ gồm mọi thứ đổi khi giữ nguyên kế hoạch: thời
gian đi mới, tiến độ thực của xe, sự kiện vừa xảy ra. Phép tách $z = x + c$ chỉ là phép trừ giữa các giờ của mô
hình, không chứng minh quyết định điều phối thực sự gây ra đúng $c$ giây chậm trễ. Hình~\ref{hinh:bamoc} là một ví
dụ dựng tay: ghép khách M làm A tới muộn hơn lời hứa 70 giây, nhưng chỉ 30 giây là phần dời do điều phối.
% EV: A01 (không phải giờ đến thật, không nhận diện nhân quả); A02; X3-04 (VALID:355-361); K04; ví dụ dựng tay

\begin{figure}[htbp]
\centering
\begin{tikzpicture}[x=1cm,y=1cm,font=\small,
  diem/.style={draw=xamvien,fill=xamnhat,rounded corners=2pt,inner sep=3pt,minimum height=6mm},
  mui/.style={-{Latex[length=2mm]},thick}]
\node[anchor=east] at (2.6,4.0) {Giữ nguyên kế hoạch:};
\node[diem] (x1) at (3.4,4.0) {Xe};
\node[diem] (a1) at (5.6,4.0) {Trả A};
\draw[mui] (x1)--(a1);
\node[anchor=east] at (2.6,3.0) {Ghép thêm khách M:};
\node[diem] (x2) at (3.4,3.0) {Xe};
\node[diem] (b2) at (5.6,3.0) {Đón M};
\node[diem] (a2) at (7.8,3.0) {Trả A};
\node[diem] (c2) at (10.0,3.0) {Trả M};
\draw[mui] (x2)--(b2); \draw[mui] (b2)--(a2); \draw[mui] (a2)--(c2);
\draw[-{Latex[length=2mm]}] (0.8,0)--(12.2,0) node[right]{giờ};
\foreach \vitri/\tenmoc/\giomoc in {2.0/v/9:10:00, 6.8/e/9:10:40, 10.4/p/9:11:10}
  \draw (\vitri,0.12)--(\vitri,-0.12) node[below,align=center]{$\tenmoc$\\\giomoc};
\draw[mui,dashed] (2.0,0.55)--(6.8,0.55) node[midway,above]{$x = e - v = 40$ giây};
\draw[mui] (6.8,0.55)--(10.4,0.55) node[midway,above]{$c = p - e = 30$ giây};
\draw[mui,densely dotted] (2.0,1.45)--(10.4,1.45) node[midway,above]{$z = p - v = 70$ giây};
\end{tikzpicture}
\caption[Ba mốc giờ trả trong một ví dụ dựng tay]{Ví dụ dựng tay: khách A được hứa trả lúc 9:10:00 ($v$); lúc 9:00, xe đang chở A xét ghép thêm khách M.
Tính với thời gian đi hiện tại, giữ nguyên kế hoạch thì A tới lúc 9:10:40 ($e$), ghép M thì tới lúc 9:11:10
($p$). Nét đứt: $x$; nét liền: $c$; nét chấm: $z = x + c$. Luật C chỉ tính $|c|$.}
\label{hinh:bamoc}
\end{figure}

Phần dời do điều phối được tính theo độ lớn: sớm hay muộn hơn $e$ 20 giây đều tính 20 giây, vì luật đo lần sửa
lời hứa chứ không đo lợi hay hại; Fielbaum và Alonso-Mora~\cite{fielbaum2020} thì chỉ tính phần bất lợi. Trên mọi
lần chạy đã ghi, $c \ge 0$ tới mức mili giây; đây là quan sát, chưa phải chứng minh cho trường hợp tổng quát. Trên
tập kiểm tra chính, khóa luận đã kiểm 136 lần chạy của bốn cấu hình: điều phối không trần, luật C trần 30 và 120
giây, luật hạn chót 30 giây. Ở các lần chạy này, mỗi lần kéo giờ trả sớm hơn $e$ chỉ sớm 1 mili giây. Hệ quả: khi
bộ điều phối chỉ chèn thêm, phần tính hai chiều và không hoàn lại của luật C gần như không có việc; nó chỉ có tác
dụng khi kế hoạch của khách đã nhận được phép đổi, ví dụ khi dời khách đang chờ sang xe khác.
% EV: A04; X3-05 (DELTA:137-139, 259-262); J06 (FAM:293-296); N-05 (CORRECTED); ADD1:104-112; ADDENDUM-1 mục 0.1–0.2 (không nêu nguyên nhân của −1 ms)

% -----------------------------------------------------------------------------------------------------
\section{Luật C: trần phần dời giờ trả do điều phối cho từng khách}
\label{muc:luatc}

\begin{dinhnghia}[Luật C]
\label{dn:luatc}
Gọi $B$ là trần, tính bằng giây, và $L_r$ là tổng các $|c|$ đã tính cho khách $r$ kể từ lúc $r$ xác nhận đặt xe,
với $L_r = 0$ lúc xác nhận. Một phương án của một xe chỉ được giữ lại nếu với mọi khách $r$ có lời hứa đang được
gán cho xe đó (chờ đón hay đang trên xe),
\begin{equation}
L_r + |p_r - e_r| \le B .
\label{eq:luatc}
\end{equation}
Sau khi quyết định của lượt được chọn, với mỗi khách $r$ có lời hứa mà phương án được chọn làm $x$, $c$ hay $z$
khác 0, $L_r$ được cộng thêm $|p_r - e_r|$ của phương án đó, kể cả khi giờ trả công bố không đổi ($p_r = v_r$).
\end{dinhnghia}
% EV: A06 (BUDGET:41-105; LEDGER:285-288); A13; X3-05 (BUDGET:107-114; CFG C-30.json:17-23); X3-25 (VALID:272-281, 562-597)

Dấu $\le$ cho phép $L_r + |c|$ đúng bằng $B$. Phương án chỉ được xét không làm đổi $L_r$: mỗi phương án được
kiểm trên bản sao của sổ, chỉ quyết định cuối của lượt mới được ghi vào sổ thật. $B$ giống một ngân sách thời
gian của từng khách: $B - L_r$ là phần trần còn lại, không được cấp lại sau mỗi lần dời. Vì phần dời được cộng
qua cả chuyến, một luật cho mỗi lần dời tới 30 giây cho phép ba lần, tổng 90 giây, còn luật C với $B = 30$ giây thì
không. Ở Hình~\ref{hinh:bamoc}, với $L_A = 0$ và $B = 30$ giây, phương án ghép M ($|c| = 30$ giây) vẫn được giữ và
A dùng hết trần. Khi phép kiểm không được miễn cho $r$ ở lượt nào, tổng $|c|$ của $r$ cả chuyến không vượt $B$.
% EV: X3-05 (BUDGET:88); X3-25 (HVA:161-174; RunnerSession.cs:956-987; VALID:640-644: miễn trừ vẫn cộng vào sổ); A06; K43; E08; J05; CFG C-30.json:3-6 (một chính sách cho mọi khách)

\begin{table}[htbp]
\centering
\caption[Giả mã phép kiểm luật C cho một phương án]{Giả mã phép kiểm luật C cho một phương án (chỉ viết chiều giờ trả)}
\label{bang:kiemtran}
\small
\begin{tabular}{@{}r@{\hspace{6pt}}p{9.6cm}@{\hspace{8pt}}p{4.6cm}@{}}
\toprule
\multicolumn{3}{@{}p{14.9cm}@{}}{\textbf{Đầu vào:} phương án $\pi$ của xe $k$; trần $B$; $L_r$ của mỗi khách
có lời hứa; trạng thái và thời gian đi của lượt. \textbf{Đầu ra:} giữ hay loại $\pi$, kèm bản ghi chặn}\\
\midrule
1: & \textbf{hàm KiểmTrần}($\pi$, $k$, $B$) & \\
2: & \quad \textbf{với mỗi} khách $r$ có lời hứa được gán cho xe $k$, theo thứ tự mã khách \textbf{thực hiện} &
chờ đón hay trên xe\\
3: & \qquad $e_r \gets$ giờ trả của $r$ khi chiếu lại kế hoạch cũ của $k$ từ trạng thái hiện tại & mốc giữ
nguyên kế hoạch\\
4: & \qquad $p_r \gets$ giờ trả của $r$ theo $\pi$, cùng trạng thái và thời gian đi; $c_r \gets p_r - e_r$ &
phần dời do điều phối\\
5: & \qquad \textbf{nếu} $L_r + |c_r| > B$ \textbf{thì trả về} loại $\pi$, kèm ($r$, $L_r$, $|c_r|$,
$L_r + |c_r|$, $B$) & dừng ở khách vi phạm gặp trước\\
6: & \quad \textbf{trả về} giữ $\pi$ & $L_r$ chỉ đổi nếu $\pi$ được chọn\\
7: & \textbf{kết thúc hàm} & \\
\bottomrule
\end{tabular}
\end{table}
% EV (dòng 2–4): X3-04 (VALID:233-239, 355-361; HVA:171); X3-27 (VALID:1048-1064)
% EV (dòng 5–6): X3-05 (BUDGET:6-12, 72-104); A06; A23; X3-25

Cài đặt theo dõi mười chiều cam kết: thay đổi giờ đón, giờ trả; số lần sửa đáng kể; đổi xe; khoảng dời và số lần
đổi của điểm đón, của điểm trả; số lần đảo thứ tự với khách cũ; số điểm dừng mới chèn trước điểm đón. Chỉ bốn
chiều có trần: giờ trả ($B$) và đổi xe, đổi điểm đón, đổi điểm trả (trần 0). Ở tập kiểm tra chính và tập buổi sáng,
phép kiểm dừng ở khách vi phạm gặp trước, nên mỗi bản ghi chặn (bản ghi lý do loại phương án) chỉ nêu lỗi của khách
đó.
% EV: A04 (DIM:3-32); X3-26 (DELTA:137-153); A25; A07; CFG C-30.json:7-89; X3-07 (SEL:319-330, 510-522); A23; X3-04 (VALID:222-239); WP4:1-25; X3-10

\subsection{Phương án giữ nguyên kế hoạch dưới luật C}
\label{muc:giunguyen}

Gọi \textbf{(N)} là giả thiết, áp cho lượt đang xét: (i) kế hoạch đang chạy không vi phạm ràng buộc cấu trúc
(sức chứa, đón trước trả, phần kế hoạch đã cố định), nên bộ sinh dựng được phương án giữ nguyên; (ii) giờ trả
theo phương án giữ nguyên được tính bằng đúng phép chiếu dùng cho $e$. RideBound thỏa (ii) vì dùng cùng một hàm
chiếu với cùng đầu vào, và luôn giữ phương án giữ nguyên khi cắt còn 100 phương án.
% EV: A22; X3-31 (GEN:171-182, 213-229); ServiceQualityAllowance.cs:66-73; X3-27 (VALID:355-360, 1048-1064; HVA:144-165; CandidatePortfolioRetainer.cs:102-113; WP4:15)

\begin{nhanxet}
\label{nx:giunguyen}
Giả sử (N) và mọi khách có lời hứa của xe có $0 \le L_r \le B$. Khi đó luật C không loại phương án giữ nguyên
kế hoạch.
\end{nhanxet}

\noindent\textit{Chứng minh.} Theo (N)(ii), với mọi khách $r$ có lời hứa của xe, $p_r = e_r$, nên
$L_r + |p_r - e_r| = L_r \le B$ và phép kiểm~\eqref{eq:luatc} cho qua. \hfill$\square$
% EV: A22 (theo cấu tạo, giả thiết N); A06; X3-27

Khóa và các trần 0 cũng so với kế hoạch giữ nguyên, nên phương án này qua cả chúng. Nếu (N) đúng ở mọi lượt thì
$L_r \le B$ được giữ: $L_r = 0$ lúc xác nhận, và nếu $L_r \le B$ ở đầu một lượt thì không có miễn trừ nào, phương
án được chọn đã qua~\eqref{eq:luatc}, nên $L_r$ mới vẫn không vượt $B$. Khi đó xe dùng luật C hay điều phối không
trần không bao giờ bị ép; tính chất này vì vậy không riêng của luật C.
% EV: A10 (LOCK:156-161); X3-06; X3-26 (véc-tơ c của phương án giữ nguyên bằng 0 ở mọi chiều); X3-09 (HVA:255-259); X3-25

Nhận xét đúng do cách đặt luật, không phải kết quả thực nghiệm; số liệu chỉ phù hợp với nó: trên tập kiểm tra
chính, 0 trong 20.101 quyết định bị ép với trần 30 giây và 0 trong 20.435 với trần 120 giây (mỗi quyết định ứng với
một lượt của cả đội xe).
% EV: A21; N-21; X3-16 (STG:43-44); N-07; X3-30 (RunnerSession.cs:1013-1027); SPEC:31-34

Nhận xét có ba giới hạn. Thứ nhất, nếu (N)(i) không thỏa, lần chạy dừng vì lỗi cấu trúc, không vì luật C; phần
vượt cửa sổ đón hay thời gian ngồi xe do tắc đường không làm (N) hỏng, nhờ phần nới. Thứ hai, với $v$ = 9:10:00 và
$e = p$ = 9:11:00, phương án giữ nguyên có $c = 0$ nhưng $z = 60$ giây: luật C không bảo đảm giờ đến. Thứ ba, với
$v = p$ = 9:10:00 và $e$ = 9:12:00, kéo giờ trả về đúng giờ đã hứa có $z = 0$ nhưng $|c| = 120$ giây, và $L_r$
tăng 120 giây dù giờ công bố không đổi. Tính chất và giới hạn thứ hai là hai mặt của một lựa chọn: luật C không
bị ép vì không tính phần thay đổi khi giữ nguyên kế hoạch, và vì thế cũng không giới hạn phần đó.
% EV: A22; X3-31 (GEN:171-182; RunnerSession.cs:941-947); X3-23; K05; X3-25 (VALID:562-597: bản ghi sửa khi x ≠ 0 dù p = v); ví dụ dựng tay

% -----------------------------------------------------------------------------------------------------
\section{Các luật đối chứng, bộ chọn và luật phục hồi}
\label{muc:doichung}

Bảng~\ref{bang:luat} định nghĩa bốn cấu hình mà Chương~\ref{chuong:thucnghiem} so với luật C. Trừ ghép tự do,
chúng dùng cùng bộ sinh, bộ kiểm, bộ chọn và luật phục hồi với luật C; tập tắc đường cố định chạy bằng bộ chạy
chưa có luật phục hồi.
% EV: A08; A09; A11; A17; A19; K50; HO:109; B17; NEW:13

\begin{table}[htbp]
\centering
\caption[Các cấu hình được so sánh]{Các cấu hình được so sánh; ký hiệu chỉ dùng trong bảng. $B$ là trần, tính bằng giây; với luật hạn chót,
$B$ là số giây giờ trả được phép muộn hơn lời hứa đầu $p_0$. Các tập thử được mô tả ở mục~\ref{muc:c4-thietlap}.}
\label{bang:luat}
\footnotesize
\setlength{\tabcolsep}{3pt}
\begin{tabular}{@{}p{2.2cm}p{1.1cm}p{3.1cm}p{2.3cm}p{2.6cm}p{3.5cm}@{}}
\toprule
\textbf{Tên trong khóa luận} & \textbf{Ký hiệu} & \textbf{Kiểm điều gì} & \textbf{Khóa} & \textbf{Khi xe bị
ép} & \textbf{Dùng ở tập nào}\\
\midrule
Điều phối không trần & U & Không kiểm giờ trả & Xe, điểm đón, giờ đón; điểm trả có trần 0 & Không bị ép khi (N)
đúng & Tập kiểm tra chính; tập chỉnh trần buổi sáng; tập tắc đường cố định\\
% EV: A07; A08; K50; A10; X3-16 (STG:42: 0 / 20.760); B10; B15; B17; NEW:11
Ghép tự do & FREE & Không kiểm lời hứa; chọn theo chi phí kế hoạch ở từng lượt & Không khóa & Không áp dụng &
Tập tắc đường cố định\\
% EV: A17; K50; HO:109 (b1-rolling-cost, không khóa); B17; X3-38
Luật C & C-$B$ & $\sum |c| \le B$ cả chuyến, cho mọi khách có lời hứa & Như U & Không bị ép khi (N) đúng
(Nhận xét~\ref{nx:giunguyen}) & Tập chính (30, 120); buổi sáng (30, 60, 90, 120, 150, 180, 240, 300); tắc đường
cố định (30, 60, 120); chiều tuyết (30)\\
% EV: A06; A22; B10; L1a-01; X3-28 (run_w07.py:5; run_w07x.py:3-4, 35-36); B17; B20; N-26 (chiều tuyết 0 / 5.067 quyết định bị ép)
Luật hạn chót & M$^{+}$-$B$ & $p - p_0 \le B$ ở mọi lần kiểm; chỉ chặn phía muộn & Như U & Có thể bị ép do tắc
đường; phục hồi không tệ hơn & Tập chính (30, 60); buổi sáng (tám mức như luật C)\\
% EV: A09; A10; A19; B10; X3-28; D04; LOCK:186-197; CFG Mplus-30.json:98
Luật trần độ dời khách thấy & V-$B$ & $\sum |z| \le B$ cả chuyến & Như U & Có thể bị ép do tắc đường; phục hồi
không tệ hơn & Tập chính (30); buổi sáng (30, 90, 120, 150, 180, 240)\\
% EV: A11; A19; B10; X3-28 (run_w07x.py:3-4, 35-36: V-90/120/150/180/240; run_w07.py:5: V-30)
\bottomrule
\end{tabular}
\end{table}

Luật hạn chót chỉ mượn thành phần hạn chót tính từ lời hứa đầu của Schulz và Pfeiffer~\cite{schulz2026}, vốn
dùng thời gian đi cố định và độ chùng mặc định 900 giây. Ở đây mỗi độ chùng bằng một mức trần của luật C trên cùng
tập, và với tắc đường đổi theo khung 15 phút, riêng tắc đường có thể làm luật này bị ép; số quyết định bị ép ở
Chương~\ref{chuong:thucnghiem} nói về một hạn chót chặt như vậy, không về phương pháp của họ.
% EV: J10 (SP:466-490, 568-577); J11; A08; A09; X3-28; X3-41 (D03:214-217); A10

Ba luật lời hứa khác nhau ở đại lượng được cộng. Đánh số các lần cập nhật lời hứa của một khách là
$1, \ldots, n$. Trong sổ, mỗi bản ghi sửa so với lời hứa ngay trước nó ($v_j = p_{j-1}$), nên các giờ ở giữa
triệt tiêu:
\begin{equation}
\sum_{j=1}^{n} (x_j + c_j) \;=\; \sum_{j=1}^{n} z_j \;=\; p_n - p_0 .
\label{eq:dongnhat}
\end{equation}
Luật hạn chót kiểm $p_n - p_0 \le B$, tức cộng có dấu cả hai phần; luật trần độ dời khách thấy cộng $|z_j|$; luật
C chỉ cộng $|c_j|$. Ở Hình~\ref{hinh:bamoc}, nếu 9:10:00 là lời hứa đầu, phương án giữ nguyên có
$e - p_0 = |z| = 40$ giây, vượt trần 30 giây của hai luật kia dù $c = 0$: riêng tắc đường đủ làm chúng bị ép.
% EV: A12; X3-33 (LEDGER:256-269, 298-314; SPEC:62-63); A09; A10 (LOCK:156-161); A11; X3-16 (U 0 / 20.760)

Tiền lệ trực tiếp là Hörl và Zwick~\cite{horl2022} (mục~\ref{muc:c2-rangbuoc}): ở đầu mỗi lượt họ tính mốc $e$,
và công thức 7 đặt giờ trả bắt buộc bằng giá trị lớn hơn giữa giờ cập nhật và giờ trả muộn nhất, như ràng buộc cứng
(ràng buộc 9--10), kèm phạt cho phần vượt giờ gốc (công thức 11). Mục đích của họ, tránh tình huống không còn thứ tự
trả khách khả thi, cũng là mục đích của Nhận xét~\ref{nx:giunguyen}. Đặt giờ trả muộn nhất là $p_0 + B$ cho hạn chót
nới $p \le \max(p_0 + B, e)$, một cách viết lại của khóa luận; nó cho phép $c \le \max(p_0 + B - e, 0)$, về 0 khi
tắc đường đẩy $e$ quá $p_0 + B$, và cũng không loại phương án giữ nguyên. Cấu hình hạn chót M$^{+}$ với luật phục
hồi ``không tệ hơn'' ở chương~4 chính là dạng này. Luật C khác ở ba chỗ:
một trần cho mọi khách có lời hứa, còn phần nới của họ chỉ tác dụng khi giờ cập nhật đã vượt giờ trả muộn nhất;
phần được phép dời, $|c| \le B - L_r$, không phụ thuộc tắc đường; và cộng qua cả chuyến, dù phần cộng dồn hiếm khi là
lý do loại phương án (mục~\ref{muc:c4-congdon}).
% EV: J01; J02 (HZ:481-497); J03; X3-36 (HZ:531-543, 565-606); J05; A22 (suy luận V17 R6); K02; K26; K27; N-15

Từ Fielbaum và Alonso-Mora~\cite{fielbaum2020} (mục~\ref{muc:c2-dolech}), khóa luận kế thừa ý tách phần thay đổi
do quyết định gán xe và kiểm soát nó trước khi quyết định, nhưng mỗi lần sửa được so với mốc $e$ chứ không so điểm
đầu với điểm cuối, nên hai lần dời ngược chiều không triệt tiêu nhau; chiều sớm cũng được tính; đại lượng là giờ
trả chứ không phải thời gian chờ; và tắc đường thay đổi chứ không chọn giờ thấp điểm.
% EV: J06; J07; J08; X3-35 (FAM:90-94, 189-195, 216-220, 293-296, 948-972); K30; K32

Vì đổi xe, đổi điểm đón, đổi điểm trả là các trần 0, phần đầu của tiêu chí 3 (mức dùng trần lớn nhất, tỉ số lớn
nhất giữa phần đã dùng và trần) luôn bằng 100\% ở các lượt đã ghi của tập kiểm tra chính (mục~\ref{muc:c4-hople}),
cả khi không trần; tiêu chí 3 vì vậy chỉ tác dụng qua tổng phần dời theo từng chiều. Điều phối không trần và luật C
cùng cấu hình và bộ chọn, chỉ khác trần giờ trả; chênh lệch giữa chúng là phần đi kèm việc thêm trần khi giữ nguyên
bộ chọn. Khóa luận còn xét thứ tự ưu tiên chi phí: giữ ưu tiên nhận yêu cầu và các ràng buộc, bỏ các mức xếp
hạng theo ổn định rồi so chi phí kế hoạch và tiêu chí phân xử. Phương án vẫn phải qua phép kiểm trần khi trần
được bật. Mục~\ref{muc:c4-vaitro} so hai thứ tự ở trần 30 giây và không trần; đổi thứ tự không làm thay đổi
định nghĩa hay cách cộng trần của luật C.
% EV: S15-03 (costfirst-diff, cấu hình và provenance); S15-01 (mục 0, 2).
% EV: A14; X3-07 (SEL:271-330; HVA:537-636); A15; N-25; A08; H01; H02; K18

Luật phục hồi giữ nguyên kế hoạch giữ phương án giữ nguyên làm phương án được miễn trừ nếu mọi lỗi của nó là lỗi
hạn chót hay lỗi trần. Luật phục hồi không tệ hơn mở rộng luật này: giữ thêm phương án đổi kế hoạch có mỗi lỗi,
cùng khách và cùng chiều, cũng là lỗi của phương án giữ nguyên và không nặng hơn. Lỗi khóa theo giai đoạn (khóa của
khách chờ đón hay trên xe) và lỗi ràng buộc dịch vụ không bao giờ được miễn. Mỗi khách được miễn được ghi một vi
phạm, trừ khi lời hứa, giai đoạn và mã lỗi không đổi so với lần trước. Tập kiểm tra chính, tập buổi sáng và tám
lần chạy chiều tuyết dùng luật thứ hai; cả hai là chi tiết cài đặt, không phải đóng góp của khóa luận.
% EV: A17; A18 (ADR:5066-5092); X3-09 (HVA:79-99, 255-326; VALID:413-420, 475-483, 537-544, 599-638, 791-818; LOCK:42-60); A19 (RUN38:9; WP4:7); X3-10; R26-CH3-OPEN-QUESTIONS-CODEX.md §Q8 (tám lần chạy chiều tuyết; wp4-tier3-v2.json:7)

% -----------------------------------------------------------------------------------------------------
\section{Lời hứa theo dự báo}
\label{muc:dubao}

Vì bộ điều phối dùng thời gian đi hiện tại, lời hứa có thể sớm hơn giờ trả thực tế khi đường đông dần, dù xe giữ
nguyên kế hoạch. Lời hứa theo dự báo, một cách làm đã có (mục~\ref{muc:c2-dubao}), cộng sẵn mức tăng tắc đường các
ngày trước cho thấy vào từng giờ đón, giờ trả được công bố.
% EV: A05; F10; J35; F01; X3-11 (DISP:1-8)

Thời gian đo bằng giây kể từ đầu khung giờ của kịch bản; $i(\tau) = \lfloor \tau / 900 \rfloor$ là khung 15 phút
chứa $\tau$. Gọi $m_i$, $i = 0, \ldots, 8$, là trung bình hệ số tắc đường của khung $i$ trên các ngày nguồn; từ
khung 8 hệ số được giữ ở giá trị cuối, như bộ mô phỏng cũng giữ, nên đặt $\tilde m_i = m_{\min(i, 8)}$. Với lúc
công bố $t$ và giờ công bố $p > t$ (giờ của phương án được chọn, chưa cộng dự báo), $S$ là nghiệm của
\begin{equation}
\int_{t}^{S} \frac{\mathrm{d}\tau}{g(\tau)} = p - t, \qquad
g(\tau) =
\begin{cases}
1 & \text{nếu } i(\tau) = i(t),\\
\tilde m_{i(\tau)} / \tilde m_{i(t)} & \text{nếu không,}
\end{cases}
\label{eq:keodan}
\end{equation}
và lời hứa hiển thị là $D = \max(p, S)$ (giả mã ở Bảng~\ref{bang:dubao}).
% EV: F01; X3-11 (DISP:3-8, 55-64; STRETCH:1-12, 22-52; preflight_world_v1_3.py:51-82)

Theo tốc độ lúc công bố, xe cần thêm $p - t$ giây. Trong khung chậm hơn khung hiện tại $g$ lần, mỗi giây đồng hồ chỉ
đi được phần đường mà ở tốc độ lúc công bố mất $1/g$ giây. Mức tắc đường lúc này đã có trong thời gian đi đang dùng, nên $g = 1$ ở khung
hiện tại; các ngày khác chỉ cho biết khung sau chậm hay nhanh hơn bao nhiêu lần. Phép lấy $\max$ bảo đảm $D \ge p$ ở
cùng lần công bố, không bảo đảm $D$ tăng dần qua các lần cập nhật. Ví dụ dựng tay: trong kịch bản khung
8:00--10:00, lúc $t$ = 9:05:00 hệ thống công bố $p$ = 9:25:00, tức còn 1.200 giây. Nếu các ngày nguồn có hệ số trung
bình 1,25 ở khung 9:00--9:15 và 1,50 ở khung 9:15--9:30, thì sau 9:15 $g = 1{,}50/1{,}25 = 1{,}2$. Tới 9:15 xe đi
được 600 giây quãng đường; 600 giây còn lại mất $600 \times 1{,}2 = 720$ giây, nên $S = D$ = 9:27:00, dài thêm 120
giây.
% EV: X3-11 (STRETCH:1-12; forecast_weight neo ở khung hiện tại); F01; DISP:60-64; ví dụ dựng tay X3-37 (a), tính bằng hàm stretch của dự án; số học ở 07-chuong3-v3-NOTES.md mục 4

\begin{table}[htbp]
\centering
\caption[Giả mã lời hứa theo dự báo]{Giả mã lời hứa theo dự báo cho một giờ đón hay giờ trả đã công bố}
\label{bang:dubao}
\small
\begin{tabular}{@{}r@{\hspace{6pt}}p{9.6cm}@{\hspace{8pt}}p{4.6cm}@{}}
\toprule
\multicolumn{3}{@{}p{14.9cm}@{}}{\textbf{Đầu vào:} lúc công bố $t$; giờ công bố $p$; ngày $d$ và khung giờ của
kịch bản; hệ số tắc đường theo khung của các ngày trước $d$. \textbf{Đầu ra:} lời hứa hiển thị $D$}\\
\midrule
1: & \textbf{hàm LờiHứaTheoDựBáo}($t$, $p$, $d$) & \\
2: & \quad $H \gets$ các ngày cùng loại với $d$ và trước $d$, bỏ ngày 15/11 & chỉ dùng ngày đã qua\\
3: & \quad \textbf{nếu} $H$ có ít hơn 2 ngày \textbf{hoặc} $p \le t$ \textbf{thì trả về} $p$ & không dự báo\\
4: & \quad $m_i \gets$ trung bình hệ số của khung $i$ trên $H$, $i = 0, \ldots, 8$ & khung tính từ đầu khung giờ\\
5: & \quad $W \gets p - t$;\; $\tau \gets t$;\; $i_0 \gets i(t)$ & quãng còn lại theo tốc độ hiện tại\\
6: & \quad \textbf{lặp} & \\
7: & \qquad $g \gets 1$ nếu $i(\tau) = i_0$, ngược lại $g \gets \tilde m_{i(\tau)} / \tilde m_{i_0}$ &
$\tilde m_i = m_{\min(i, 8)}$\\
8: & \qquad $h \gets$ thời gian từ $\tau$ tới hết khung $i(\tau)$ & \\
9: & \qquad \textbf{nếu} $h/g \ge W$ \textbf{thì} $S \gets \tau + W g$ và thoát vòng lặp & \\
10: & \qquad $W \gets W - h/g$;\; $\tau \gets$ đầu khung kế tiếp & \\
11: & \quad \textbf{trả về} $D \gets \max(p, S)$ & $D \ge p$ ở lần công bố này\\
12: & \textbf{kết thúc hàm} & \\
\bottomrule
\end{tabular}
\end{table}
% EV (dòng 2–4): F02; X3-11 (DISP:27-57; preflight_world_v1_3.py:51, 80-82)
% EV (dòng 5–11): F01; X3-11 (STRETCH:25-39, 49-52)

Ngày nguồn là các ngày cùng loại (ngày thường hay cuối tuần) trước ngày của kịch bản, ít nhất 2 ngày, và không bao
giờ là ngày tuyết 15/11/2018 (một lựa chọn của khóa luận; hệ thống thật cần quy tắc tự động nhận ra ngày bất
thường). Vì vậy trong dữ liệu 12--18/11/2018 chỉ ba ngày 14, 15, 16/11 có dự báo: 21 trong 34 kịch bản của tập
kiểm tra chính; 13 kịch bản ngày 17, 18/11 giữ lời hứa không dự báo.
% EV: F02 (DISP:10-16, 27-52); F04; N-16; N-28

Trong các kết quả chính, $D$ được tính lại trên các lần chạy đã ghi, gọi là phát lại (replay): xe nào chở ai giữ
nguyên, chỉ giờ báo cho khách đổi, và $D$ không quay lại sổ hay phép kiểm trần. $D_0 - p_0$, với $D_0$ là lời hứa
đầu có dự báo, là giá của dự báo. Phép kiểm khả thi vẫn dùng giờ không dự báo, nên dưới luật C trần 30 giây, trong
1.471 khách được phục vụ ở 21 kịch bản có dự báo, 34 khách có giờ đón hiển thị ở lời hứa đầu muộn hơn giờ đón
muộn nhất; chỗ này cần sửa trước khi đưa dự báo vào bộ điều phối. Phép thử đưa dự báo vào bộ điều phối
(mục~\ref{muc:c4-dubao}) thì khác: nó đổi thời gian đi mà bộ điều phối nhận, còn FleetPy vẫn chạy theo hệ số đo.
Số khung lấy trung bình là $K$: các cấu hình chính bật dự báo dùng $K = 2$, gồm khung 15 phút hiện tại và khung kế tiếp.
Phép thử giả dùng $K = 1$, cho hệ số bằng 1; khi tắt dự báo không đặt $K$.
% EV: F03 (chỉ phần "phát lại trên các lần chạy đã ghi"); FR:10-11; X3-12 (FFA:84-99); F06; F07 (FR:70); F12; X3-29 (preflight_world_v1_3.py:1-26, 113-133); R26-CH3-OPEN-QUESTIONS-CODEX.md §Q9; preflight_world_v1_3.py:85-89, 99-101; run_forecast.py:67-72

% -----------------------------------------------------------------------------------------------------
\section{Sổ theo từng khách}
\label{muc:so}

Sổ do RideBound ghi trong lúc chạy: mỗi khách có bản ghi lời hứa đầu và một bản ghi sửa cho mỗi lượt mà phương án
được chọn làm $x$, $c$ hay $z$ khác 0 ở bất kỳ chiều nào. Bản ghi sửa lưu các phép chiếu ứng với $v$, $e$, $p$,
độ lớn thay đổi trên mười chiều và $L$ sau lần sửa (Bảng~\ref{bang:vidu-so}); các hiệu có dấu $x$, $c$, $z$ chỉ dùng
cho giờ. Không có ngưỡng: thay đổi 1 mili giây cũng tạo bản ghi sửa, nên số bản ghi sửa không phải số lần một ứng
dụng thật báo lại cho khách.
% EV: A04; X3-26 (DELTA:64-160); X3-25 (VALID:562-597); LEDGER:285-314; X3-13 (SPEC:35-46)

\begin{table}[htbp]
\centering
\caption[Bản ghi sổ của khách A trong ví dụ dựng tay]{Hai bản ghi giờ trả của khách A ở Hình~\ref{hinh:bamoc} (dựng tay), với giả định 9:10:00 là lời hứa đầu,
công bố lúc 8:50:00, và $B = 30$ giây. Bản ghi lời hứa đầu không có lời hứa trước nên sổ quy ước $v = e = p$. Các
độ lớn và $L$ tính bằng giây.}
\label{bang:vidu-so}
\small
\begin{tabular}{@{}lllllrrrr@{}}
\toprule
\textbf{Bản ghi} & \textbf{Lúc ghi} & $v$ & $e$ & $p$ & $|x|$ & $|c|$ & $|z|$ & $L$\\
\midrule
Lời hứa đầu & 8:50:00 & 9:10:00 & 9:10:00 & 9:10:00 & 0 & 0 & 0 & 0\\
Sửa (ghép khách M) & 9:00:00 & 9:10:00 & 9:10:40 & 9:11:10 & 40 & 30 & 70 & 30\\
\bottomrule
\end{tabular}
\end{table}
% EV: ví dụ dựng tay; v = e = p ở bản ghi đầu theo SPEC:46; L = 0 + |c| theo LEDGER:285-288

Trong một bản ghi mô phỏng của luật C trần 120 giây, một khách có $v$ = 17:43:46,044, $e$ = 17:43:46,283 và
$p$ = 17:44:01,203. Phần dời $|p - e| = 14{,}92$ giây làm phần trần còn lại bằng $120 - 14{,}92 = 105{,}08$ giây.
Cả 14 bản ghi sửa sau có $c = 0$ và $|x| < 0{,}25$ giây, nên phần trần đã dùng không tăng.
% EV: round26-trace-rider.txt, Q3 SELECTED RIDER, entries 1-15; R26-CH3-OPEN-QUESTIONS-CODEX.md §Q3; C-120-d20181117-s10-r1-w17/transcript-00.ndjson, final checkpoint ordinal 1716, request rq-c48d442603e9ce8c056b029d6d50c497f4d224b8af469945692690e92a75afb0

Sau khi chạy, độ trễ so với lời hứa đầu được chia ba phần: phần tính cho điều phối, phần được hồ sơ lịch sử giải
thích và phần còn lại,
\begin{equation}
a - p_0 \;=\; q \;+\; f \;+\; (a - D_0 - q), \qquad q = \sum_{j} c_j, \qquad f = D_0 - p_0 .
\label{eq:so}
\end{equation}
Cộng ba phần thì $q$ và $D_0$ triệt tiêu, nên đẳng thức luôn đúng. Trong đó, $a$ là giờ trả thực tế trong mô
phỏng; $q$ cộng các $c$ giữ nguyên dấu (dời muộn dương, dời sớm âm), khác với $\sum |c|$ của luật C, và dưới luật
C, $|q| \le \sum |c| \le B$; $f$ là phần dự báo đã cộng sẵn vào lời hứa đầu, bằng 0 khi không có dự báo.
Theo~\eqref{eq:dongnhat}, phần còn lại bằng $\sum_j x_j + \rho - f$ với $\rho = a - p_n$: thay đổi khi giữ nguyên
kế hoạch mà dự báo chưa tính, cộng độ lệch giữa giờ trả thực tế và lời hứa cuối $p_n$. Với khách A ở
Bảng~\ref{bang:vidu-so}, nếu $D_0$ = 9:10:30 và A tới đúng lời hứa cuối, $a$ = 9:11:10, thì độ trễ 70 giây gồm
$q = 30$, $f = 30$ và 10 giây còn lại, đúng bằng $\sum x + \rho - f = 40 + 0 - 30$.
% EV: X3-12 (FFA:84-97: q = Σ(p − e) có dấu, D0 = lời hứa hiển thị đầu); A12; A06; F09; ví dụ dựng tay X3-37 (b) (số học ở NOTES mục 4)

Chương~\ref{chuong:thucnghiem} cộng ba phần trên các khách có $a - p_0 > 60$ giây, với $p_0$ là lời hứa giờ trả
đầu không dự báo (không dùng $D_0$). Đây là phép tính theo mô hình, không kết luận ai có lỗi: ``phần tính cho điều
phối'' chỉ nói rằng mô hình gán phần này cho quyết định điều phối khi so với kế hoạch giữ nguyên. Phần được hồ sơ
lịch sử giải thích phụ thuộc dự báo đang dùng; phần còn lại là phần dự báo đó chưa giải thích được, không phải phần
``không thể biết trước''.
% EV: X3-12 (FFA:134-138: nhóm A − P0 > 60 s); F09; K04; A01 (lưu ý D03:222-224); K19

% -----------------------------------------------------------------------------------------------------
\section{Cài đặt, kiểm thử và tái lập}
\label{muc:caidat}

\subsection{Kiến trúc và trao đổi với bộ mô phỏng}
\label{muc:kientruc}

RideBound viết bằng C\# trên .NET 10, chia thành các dự án theo tầng (Bảng~\ref{bang:kientruc}). Tầng nghiệp vụ
không tham chiếu dự án hay gói thư viện ngoài nào, OR-Tools chỉ có trong cầu nối bộ giải, và một bộ kiểm thử kiến
trúc kiểm các quy tắc này; nhờ vậy phép kiểm trần, ba mốc và sổ không phụ thuộc bộ mô phỏng hay bộ giải.
% EV: X3-14 (các tệp .csproj: net10.0, ProjectReference, PackageReference; DependencyRuleTests.cs:8-31, 137-166)

\begin{table}[htbp]
\centering
\caption[Các thành phần chính của RideBound và bộ chuyển đổi FleetPy]{Các thành phần chính của RideBound và bộ chuyển đổi FleetPy}
\label{bang:kientruc}
\small
\begin{tabular}{@{}p{4.6cm}p{10.2cm}@{}}
\toprule
\textbf{Thành phần} & \textbf{Vai trò}\\
\midrule
Domain (tầng nghiệp vụ) & Yêu cầu và giai đoạn của khách, lời hứa, khóa, sổ, phép kiểm trần\\
Application (tầng ứng dụng) & Chiếu lịch theo thời gian đi, tính ba mốc giờ trả, kiểm phương án theo khóa và luật
lời hứa\\
Algorithms (thuật toán) & Bộ sinh chỉ chèn thêm, đánh giá phương án theo luật lời hứa và luật phục hồi, bộ chọn\\
Solvers.OrTools (cầu nối bộ giải) & Cầu nối tới bộ giải CP-SAT của OR-Tools 9.15\\
Contracts (định dạng bản tin) & Định dạng các bản tin trao đổi với bộ mô phỏng\\
Runner (bộ chạy) & Nhận bản tin, ghép các thành phần trên và trả quyết định\\
Benchmarking (chuẩn bị dữ liệu) & Chuẩn hóa dữ liệu Manhattan của FleetPy thành kịch bản thử\\
Bộ chuyển đổi FleetPy (Python) & Phần điều khiển đội xe của FleetPy 1.0.2: gửi sự kiện cho bộ chạy, áp quyết định
vào mô phỏng\\
\bottomrule
\end{tabular}
\end{table}
% EV: X3-14 (Google.OrTools 9.15.6755 chỉ ở RideBound.Solvers.OrTools.csproj:9); X3-22 (OrTools:3, CP-SAT); X3-15; NORM; B01

Bộ chuyển đổi viết bằng Python đóng vai phần điều khiển đội xe của FleetPy 1.0.2~\cite{engelhardt2026}, ghim theo
phiên bản; nó chạy bộ chạy RideBound như một chương trình con, hai bên nhắn tin qua đầu vào, đầu ra chuẩn, mỗi bản
tin một dòng JSON (Hình~\ref{hinh:bantin}). Bộ chạy chỉ nhận lô mới sau khi bộ mô phỏng báo đã áp dụng quyết định
trước, nên các quyết định được áp dụng đúng thứ tự. Bản tin của mỗi lần chạy được lưu, và mọi con số ở
Chương~\ref{chuong:thucnghiem} được tính từ đó.
% EV: B01; J18 (FLEET:474); X3-15 (README:1-12; runner_client.py:133-138; Program.cs:103-106; RunnerSession.cs:448-455); X3-13 (SPEC:28-37); A26

\begin{figure}[htbp]
\centering
\begin{tikzpicture}[y=1cm,font=\small,
  ben/.style={draw=xamvien,fill=xamnhat,rounded corners=2pt,align=center,inner sep=4pt,text width=4.3cm},
  tin/.style={-{Latex[length=2mm]},thick},
  nhan/.style={midway,above,font=\footnotesize,align=center}]
\node[ben] (L) at (0,0) {FleetPy 1.0.2\\và bộ chuyển đổi (Python)};
\node[ben] (R) at (9.5,0) {Bộ chạy RideBound\\(C\#, .NET 10)};
\draw[xamvien,thick] (L.south)--(0,-12.5);
\draw[xamvien,thick] (R.south)--(9.5,-12.5);
\draw[tin] (0,-1.8)--(9.5,-1.8) node[nhan]{chào hỏi; khởi tạo lần chạy\\(kịch bản, cấu hình luật)};
\draw[tin] (9.5,-2.9)--(0,-2.9) node[nhan]{đã nhận chào hỏi; đã khởi tạo};
\draw[dashed,xamvien] (-0.6,-3.35) rectangle (12.2,-9.1);
\node[anchor=north west,font=\footnotesize\itshape] at (-0.5,-3.45) {mỗi lượt quyết định};
\draw[tin] (0,-5.1)--(9.5,-5.1) node[nhan]{lô sự kiện: yêu cầu mới, xác nhận đặt xe,\\tiến độ xe, đón, trả,
thời gian đi mới};
\node[ben,text width=3.9cm,font=\footnotesize] at (9.5,-6.15) {sinh phương án, kiểm luật, chọn, ghi sổ
(Bảng~\ref{bang:luotquyetdinh})};
\draw[tin] (9.5,-7.7)--(0,-7.7) node[nhan]{quyết định: kế hoạch mới của các xe, lời hứa};
\draw[tin] (0,-8.7)--(9.5,-8.7) node[nhan]{đã áp dụng quyết định};
\draw[tin] (0,-10.15)--(9.5,-10.15) node[nhan]{yêu cầu trạng thái};
\draw[tin] (9.5,-11.15)--(0,-11.15) node[nhan]{trạng thái, gồm sổ theo từng khách};
\draw[tin] (0,-12.15)--(9.5,-12.15) node[nhan]{kết thúc: bộ chạy dừng, không trả lời};
\end{tikzpicture}
\caption[Các bản tin giữa bộ chuyển đổi FleetPy và bộ chạy RideBound]{Các bản tin chính giữa bộ chuyển đổi FleetPy và bộ chạy RideBound trong một lần chạy (đã giản lược). Mỗi
bản tin là một dòng JSON; phần trong khung nét đứt lặp lại ở mỗi lượt quyết định.}
\label{hinh:bantin}
\end{figure}
% EV: X3-15 (ProtocolPrimitives.cs:104-121; RunnerSession.cs:177-214: chỉ xử lý hello, initializeRun, eventBatch, decisionApplied, checkpoint, restore; shutdown làm bộ chạy dừng, không trả lời; finalizeRun/runSummary chỉ được khai báo; runner_client.py:158-273); X3-13 (SPEC:35-37)

\subsection{Giới hạn tính toán và thời gian chạy}
\label{muc:gioihan}

Ngoài giới hạn 2 yêu cầu mỗi lượt và 100 phương án mỗi xe, các bước sinh, kiểm và giải đều có giới hạn số bước
tính, và bộ giải dùng hạt giống ngẫu nhiên (số khởi tạo phần ngẫu nhiên) cố định bằng 7 để chạy lại cho cùng kết
quả. Vì vậy bộ sinh không xét hết mọi phương án; mục~\ref{muc:c4-hople} báo các bộ đếm của bộ sinh phương án cho biết
giới hạn có cắt bớt phương án hay không. Phép kiểm ở mỗi phương án duyệt mười chiều cho mỗi khách có lời hứa của xe và đếm đảo thứ tự
trên từng cặp điểm dừng, nên không thể coi khối lượng kiểm là tuyến tính theo độ dài kế hoạch; nó chưa được đo
riêng.
% EV: A24 (WP4:9-23); X3-17; X3-32 (GEN:54-79); N-25; X3-05 (BUDGET:76-99); X3-04 (VALID:233-239); X3-26 (DELTA:190-215)

Khóa luận chỉ báo thời gian thực của cả lần chạy: trên tập tắc đường cố định (16 lần chạy mỗi cấu hình), trung vị là
730, 667 và 641 giây với ghép tự do, điều phối không trần và luật C trần 30 giây. Thời gian này gồm mô phỏng FleetPy
và thời gian chờ RideBound trả quyết định, không phải thời gian của một lượt quyết định. Chi phí riêng của phép kiểm trần
không được báo, nên các số này không dùng để so chi phí tính toán.
% EV: H13 (EL17:2-6); X3-19; B17; R26-CH3-OPEN-QUESTIONS-CODEX.md §Q6; wp14r_scientific_protocol.py:801-803, 824-829; wp14r_supervised_process.py:297, 329-330, 979-985

\subsection{Kiểm tra kết quả}
\label{muc:kiemtra}

Kết quả được kiểm ở bốn mức.
\begin{itemize}[leftmargin=1.2cm,itemsep=2pt,topsep=3pt]
\item \textbf{Kiểm thử phần mềm:} chạy lại trên đúng mã nguồn đã dựng bộ chạy, cả 1.025 trong 1.025 kiểm thử tự động
đạt trong một lần chạy đầy đủ ngày 13/10/2026, với kiểu xuống dòng LF và một tệp chưa ghi nhận trong cây làm việc;
chi tiết ở Phụ lục~\ref{phuluc:C}.
\item \textbf{Kiểm tính nhất quán của sổ:} ở cả 382.006 bản ghi sửa của 185 lần chạy (170 lần chạy của năm cấu hình
trên tập kiểm tra chính và 15 lần chạy luật C trần 30 giây trên 15 kịch bản hiệu chỉnh), $|x|$, $|c|$, $|z|$ tính
lại từ $v$, $e$, $p$ trong sổ khớp đúng các trường độ lớn. Phép kiểm này không kiểm $e$ có được tính đúng hay không.
\item \textbf{Tính lại độc lập:} một chương trình viết riêng, độc lập về mã nhưng dùng cùng đặc tả chỉ số, tính lại
13 hạng mục số liệu từ bản ghi thô và một điều kiện luôn phải đúng, không lệch hạng mục nào; các số đếm chính được
ít nhất hai bộ đọc viết độc lập tính lại và khớp trong phạm vi đã kiểm. Hai mức kiểm này được làm bằng các chương
trình đọc riêng, không phải kiểm thủ công; chúng cho biết con số khớp bản ghi, không cho biết mô phỏng giống thực tế
đến đâu.
\item \textbf{Tái lập:} tập kiểm tra chính, tập chỉnh trần buổi sáng, các lần chạy trần 120 giây và tám lần chạy chiều
tuyết dùng cùng một bản bộ chạy đã biên dịch và lưu lại; tập tắc đường cố định dùng một bản cũ hơn. Bản dựng lại, nếu có,
phải đối chiếu với kết quả đã lưu trước khi dùng. Các lượt ưu tiên chi phí dùng bản riêng; bốn cặp kiểm
hành vi không thay thế việc kiểm đồng nhất toàn bộ hai bản chạy (Phụ lục~\ref{phuluc:D}).
\end{itemize}
% EV (kiểm thử): TESTRUN-45f8026-2026-10-12.md (mã nguồn bộ chạy; lượt 3 và 3b; điều kiện LF và tệp chưa ghi nhận); A28
% EV (sổ): A26 (T38:106; T38I:68); B10 (170 + 15); RUN38:1-8 (15 lần chạy C-30 trên 15 kịch bản hiệu chỉnh)
% EV (tính lại): C19 (T38I:79-97: 13 hạng mục + 1 bất biến); ADDENDUM-1:15-19; C04; E01; K41; A26 lưu ý; I01
% EV (tái lập): A28 (RUNW07:40-41); X3-10 (run_w07.py:8; run_w07x.py:7; run_t38c120.py:3-8); B17; NEW:13; I13; R26-CH3-OPEN-QUESTIONS-CODEX.md §Q8 (bản bộ chạy chiều tuyết)

Chương~\ref{chuong:thucnghiem} dùng các thành phần trên để đo luật C được gì và mất gì trên mô phỏng FleetPy, so
với điều phối không trần, luật hạn chót và cách đệm lời hứa, tức cộng thêm một khoảng cố định vào giờ hứa cho mọi
khách.
 của main.tex.
% Bản 1 (07-chuong3.tex) và bản trung gian 2 (07-chuong3-v2.tex) được giữ nguyên. Lý do từng thay đổi:
% reviews/CH3-REVISION-LOG.md. Ghi chú v3: 07-chuong3-v3-NOTES.md; thay đổi v4: 07-chuong3-v4-CHANGES-v3.md.
% Mỗi dòng "% EV:" trỏ tới EVIDENCE-REGISTER.md (A01…K52), EVIDENCE-REGISTER-ADDENDUM-1.md (N-01…N-30,
% L1a-01…09), EVIDENCE-REGISTER-ADDENDUM-2-CH3.md (X3-01…X3-41) hoặc các nguồn R26/TESTRUN ghi tại chỗ.
% =====================================================================================================

% Môi trường "Nhận xét" chưa có trong main.tex; định nghĩa có điều kiện ở đây (nên chuyển lên phần mở đầu).
\ifcsname nhanxet\endcsname\else
  \theoremstyle{definition}
  \newtheorem{nhanxet}{Nhận xét}[chapter]
\fi

\chapter{Phương pháp: luật C, lời hứa theo dự báo và sổ theo từng khách}
\label{chuong:phuongphap}

Khi đặt xe đi chung, hành khách được báo giờ xe sẽ trả mình tới nơi; khóa luận gọi đó là lời hứa giờ trả, gọi
tắt là lời hứa. Giờ trả có thể bị dời vì bộ điều phối ghé đón thêm khách, hoặc vì đường đông hơn lúc hứa hay xe
đi nhanh, chậm hơn dự kiến. Luật C, nội dung chính của chương, chỉ giới hạn loại thứ nhất. Ở mỗi lượt quyết định,
hiệu giữa giờ trả khi
chọn phương án mới và khi xe giữ nguyên kế hoạch là phần dời do điều phối; luật C loại phương án làm tổng độ lớn
phần này qua cả chuyến vượt một trần. Nhịp tắc đường quen thuộc trong ngày được cộng trước vào lời hứa đầu nhờ dự
báo; phần còn lại được ghi trong sổ theo từng khách (per-rider ledger). Các mục đi theo thứ tự đó, sau phần phát
biểu bài toán.
% EV: A01; A02; A06; F01; X3-12 (FFA:134-138, ba phần của sổ)

% -----------------------------------------------------------------------------------------------------
\section{Phát biểu bài toán và tổng quan hệ thống}
\label{muc:baitoan}

\subsection{Thuật ngữ và vòng đời của một yêu cầu}
\label{muc:thuatngu}

Bảng~\ref{bang:thuatngu} giải thích các thuật ngữ; tên tiếng Anh chỉ ghi trong bảng. Mọi thực nghiệm chạy trên
bộ mô phỏng FleetPy~\cite{engelhardt2026} với dữ liệu chuẩn Manhattan (mục~\ref{muc:kientruc}).
% EV: B01; J18 (FLEET:1-12, 474)

\begin{table}[htbp]
\centering
\caption[Thuật ngữ chính của phương pháp]{Thuật ngữ dùng trong chương}
\label{bang:thuatngu}
\small
\begin{tabular}{@{}p{4.3cm}p{10.5cm}@{}}
\toprule
\textbf{Thuật ngữ} & \textbf{Nghĩa trong khóa luận}\\
\midrule
Lượt quyết định (decision epoch) & Một lần bộ điều phối nhận một lô sự kiện và trả quyết định cho cả đội xe.\\
Phương án (candidate plan) & Một kế hoạch (dãy điểm đón, điểm trả theo thứ tự) có thể chọn cho một xe.\\
Phương án giữ nguyên kế hoạch (keep-route option) & Xe đi tiếp kế hoạch cũ, không nhận thêm khách.\\
Bộ sinh chỉ chèn thêm (insertion-only generator) & Chèn điểm đón, điểm trả của yêu cầu mới vào kế hoạch của
một xe, không bỏ, không đảo điểm cũ.\\
Bộ chọn (selector); bộ giải (solver) & Chọn mỗi xe một phương án theo thứ tự tiêu chí (lexicographic order):
chỉ khi hòa ở tiêu chí trước mới xét tiêu chí sau; bộ giải của OR-Tools~\cite{ortools} giải bài toán chọn.\\
Lời hứa; lời hứa đầu $p_0$ & Giờ đón, giờ trả báo cho khách; lời hứa đầu công bố khi khách xác nhận đặt xe;
$p_0$ là giờ trả của nó, chưa cộng dự báo.\\
Công bố; bản ghi sửa & Ghi lời hứa mới vào sổ và gửi trong bản tin quyết định; mỗi lần sau lời hứa đầu là một
bản ghi sửa.\\
Khóa (lock); chiều cam kết & Khóa: phần lời hứa không được đổi sau khi xác nhận. Chiều cam kết: thuộc tính của
lời hứa được theo dõi riêng, như giờ trả, việc đổi xe.\\
Trần $B$ (cap); phần đã dùng $L_r$ & Tổng phần dời do điều phối tối đa của một khách cả chuyến; phần khách $r$
đã dùng.\\
Quyết định bị ép (forced decision); luật phục hồi (recovery rule) & Lượt quyết định có ít nhất một xe bị ép:
luật lời hứa đã loại cả phương án giữ nguyên kế hoạch của xe, và xe chọn phương án được giữ nhờ luật phục hồi,
có ít nhất một khách được miễn trừ; luật phục hồi quy định phương án nào được giữ khi đó.\\
Hệ số tắc đường & Số nhân chung cho cả thành phố áp lên mọi thời gian đi, đổi theo khung 15 phút; hệ số 1,2 làm
thời gian đi dài gấp 1,2 lần hệ số 1.\\
Kịch bản & Một ngày, một mẫu nhu cầu, một khung 2 giờ của dữ liệu Manhattan: 108 yêu cầu, 8 xe 4 chỗ.\\
\bottomrule
\end{tabular}
\end{table}
% EV (lượt, phương án, bộ sinh, bộ chọn): X3-01; SPEC:31-34; X3-03; X3-07; X3-22; A14
% EV (lời hứa, công bố, khóa, chiều, trần): A13; X3-25 (fleet_control.py:882-893); A07; A04; A06; K43
% EV (bị ép, phục hồi, hệ số, kịch bản): X3-30; X3-09; B06; STRETCH:3-4; B02; B03; GLOSSARY-VI-v2.md §E (bị ép); C4A-10

Một yêu cầu mới chờ được xét. Khi phương án được chọn của một xe chứa nó, yêu cầu được nhận và khách được mời,
nhưng chưa có lời hứa. Khi khách xác nhận đặt xe, lời hứa đầu được công bố và $L_r = 0$; khách chờ đón, lên xe
rồi được trả. Yêu cầu chưa được nhận được xét lại ở lượt sau nếu còn phương án hợp lệ hay chưa được xét vì giới
hạn số yêu cầu mới, và bị từ chối khi mọi phương án chứa nó đều bị loại. Khóa và luật lời hứa áp từ lúc xác nhận
tới lúc khách xuống xe.
% EV: X3-20 (VALID:272-336; RollingCostPolicy.cs:173-226; fleet_control.py:878-924); A13; X3-32

\subsection{Bài toán ở mỗi lượt quyết định}
\label{muc:baitoan-phatbieu}

Bộ chạy tiến theo đồng hồ sự kiện: yêu cầu đến, xe tới mốc vật lý và mốc cập nhật tắc đường đều có thể gọi một
lượt quyết định. Xác nhận đặt xe và cập nhật trạng thái có thể tạo thêm lượt ở cùng thời điểm mô phỏng; không có
chu kỳ điều phối cố định. Bộ chuyển đổi gom các sự kiện từ lần gọi trước (yêu cầu mới, xác nhận đặt xe, tiến độ xe,
đón, trả, thời gian đi mới) thành một lô.
% EV: R26-CH3-OPEN-QUESTIONS-CODEX.md §Q1; actual_fleetpy_medium_preflight.py:528-543, 569-637; fleet_control.py:395-451; preflight_world_v1_1.py:40-58; preflight_world_v1_2.py:18-20

Khung nhu cầu dài 2 giờ, nhưng lần chạy kết thúc khi hết yêu cầu và hành trình đang hoạt động. Trong 34 lần chạy
luật C trần 120 giây của tập kiểm tra chính, thời lượng mô phỏng từ 118,02 đến 154,93 phút.
% EV: R26-CH3-OPEN-QUESTIONS-CODEX.md §Q2; round26-trace-rider.txt, Q2 C-120 MAIN-PANEL SIMULATED LENGTHS; actual_fleetpy_medium_preflight.py:580-581, 639-683

Bài toán ở mỗi lượt như sau.
\begin{itemize}[leftmargin=1.2cm,itemsep=2pt,topsep=3pt]
\item \textbf{Đầu vào:} trạng thái xe (vị trí, kế hoạch, sức chứa, khách trên xe, khách chờ đón); sổ; yêu cầu
đang chờ với điểm đón, điểm trả, cửa sổ đón (từ giờ đón sớm nhất tới giờ đón muộn nhất), thời gian ngồi xe tối
đa; thời gian đi lúc này, coi là giữ nguyên cho phần đường còn lại.
\item \textbf{Quyết định:} nhận, để lại hay từ chối từng yêu cầu đang chờ; mỗi xe một phương án thỏa các ràng
buộc cứng (Bảng~\ref{bang:rangbuoc}).
\item \textbf{Đầu ra:} lời hứa đầu và các bản ghi sửa (mục~\ref{muc:so}).
\item \textbf{Mục tiêu:} thứ tự tiêu chí của Bảng~\ref{bang:rangbuoc}; các chỉ số đánh giá ở
Chương~\ref{chuong:thucnghiem}.
\end{itemize}
% EV: X3-21 (fleet_control.py:395-429); X3-01; X3-02; A05; X3-20; B03; A06; A07; A14

Bộ chọn giải một quy hoạch nguyên, với $y_{k\pi} \in \{0, 1\}$ bằng 1 khi xe $k$ chọn phương án $\pi$ trong
tập $P_k$ của nó:
\begin{equation}
\sum_{\pi \in P_k} y_{k\pi} = 1 \;\;\text{với mọi xe } k, \qquad
\sum_{k} \sum_{\pi \in P_k,\; u \in \pi} y_{k\pi} \le 1 \;\;\text{với mọi yêu cầu mới } u ,
\label{eq:chon}
\end{equation}
tức mỗi yêu cầu mới thuộc phương án được chọn của nhiều nhất một xe. Các tiêu chí được giải lần lượt, mỗi lần giữ
các tiêu chí trước ở giá trị tốt nhất đã tìm. Bước chọn này theo lối gán xe--chuyến động của Alonso-Mora và cộng
sự~\cite{alonsomora2017} (mục~\ref{muc:c2-baitoan}).
% EV: X3-22 (OrTools:3, 96-125, 180-189, 343-365); X3-34 (HZ:341-345); J17

\begin{table}[htbp]
\centering
\caption[Ràng buộc cứng và thứ tự tiêu chí của bộ chọn]{Ràng buộc cứng của một phương án (nửa trên) và thứ tự tiêu chí của bộ chọn (nửa dưới): so tiêu chí 1
trước, chỉ khi hòa mới xét tiêu chí 2, và cứ thế. Giá trị là cấu hình của các tập thử, không phải chuẩn ngành.}
\label{bang:rangbuoc}
\small
\begin{tabular}{@{}p{2.4cm}p{7.6cm}p{4.8cm}@{}}
\toprule
\textbf{Nhóm} & \textbf{Nội dung} & \textbf{Vai trò}\\
\midrule
Ràng buộc dịch vụ & Đón trong vòng 10 phút kể từ lúc gọi; ngồi xe không quá 1,5 lần thời gian đi thẳng; mỗi
xe tối đa 4 khách & Phương án đổi kế hoạch phải thỏa chặt; phương án giữ nguyên được nới phần vượt thời gian
mà kế hoạch cũ đã có\\
% EV: B03; X3-23 (GEN:161-246; ServiceQualityAllowance.cs:61-94)
Khóa & Không đổi xe, điểm đón, giờ đón; điểm trả không đổi nhờ trần 0. Giờ đón mới phải bằng giờ đón khi giữ
nguyên kế hoạch & So với kế hoạch giữ nguyên, nên riêng tắc đường không làm vi phạm khóa\\
% EV: A07; A08; A10; X3-06 (LOCK:42-60, 103-161, 204-207; CFG C-30.json:66-67, 94-98)
Luật lời hứa & Luật C: $L_r + |p_r - e_r| \le B$ cho mọi khách có lời hứa của xe, hoặc một luật ở
Bảng~\ref{bang:luat} & Lọc trước bộ giải\\
% EV: A06; A09; A11; HVA:62-66
\midrule
Tiêu chí 1 & Ít phương án vi phạm được giữ lại nhất (chỉ khi có xe bị ép) & Chỉ dùng chúng khi không còn cách
khác\\
Tiêu chí 2 & Nhận nhiều yêu cầu mới nhất & Ưu tiên phục vụ\\
Tiêu chí 3 & Mức dùng trần lớn nhất trên đội xe càng nhỏ càng tốt; nếu hòa, tổng phần dời do điều phối nhỏ nhất,
so lần lượt từng chiều cam kết (khi có chiều mang trần) & Ưu tiên ổn định\\
Tiêu chí 4 & Tổng chi phí kế hoạch nhỏ nhất, tức thời lượng còn lại của các tuyến dự kiến & Theo mô hình, không
phải chi phí vận hành thực\\
Tiêu chí 5 & Thứ tự mã phương án & Phá hòa, để chạy lại cho cùng kết quả\\
% EV: A14; X3-07 (SEL:271-330, 365-385, 510-522; HVA:537-636); DIM:19-32; A16; I09
\bottomrule
\end{tabular}
\end{table}

\subsection{Một lượt quyết định}
\label{muc:luotquyetdinh}

Hình~\ref{hinh:vongquyetdinh} và Bảng~\ref{bang:luotquyetdinh} mô tả một lượt. Ngoài sinh phương án và chọn,
mỗi lượt có ba bước cho lời hứa: tính các mốc giờ trả của từng khách, loại phương án vi phạm khóa hay luật lời
hứa trước khi gọi bộ giải, và ghi sổ sau khi chọn. Hình và giả mã dùng một số ký hiệu được định nghĩa ở các mục
sau: với mỗi khách có lời hứa, $v$ là giờ trả đã công bố gần nhất, $e$ là giờ trả nếu xe giữ nguyên kế hoạch, $p$ là
giờ trả theo phương án đang xét, còn $x$, $c$, $z$ là ba hiệu giữa chúng (mục~\ref{muc:bamoc}); $D$ là lời hứa theo
dự báo, tính qua giờ $S$ (mục~\ref{muc:dubao}).
% EV: X3-04 (VALID:201-239, 355-411, 494-520); X3-13; X3-25

\begin{figure}[htbp]
\centering
\begin{tikzpicture}[
  node distance=6mm,
  hop/.style={draw=xamvien,fill=xamnhat,rounded corners=3pt,align=center,
              inner sep=4pt,font=\small,minimum height=12mm,text width=2.75cm},
  ngoai/.style={draw=xamvien,dashed,rounded corners=3pt,align=center,
              inner sep=4pt,font=\small,minimum height=10mm,text width=4.2cm},
  mui/.style={-{Latex[length=2.2mm]},thick}
]
\node[hop] (a) {Nhận lô sự kiện\\mới};
\node[hop,right=of a] (b) {Sinh phương án chỉ chèn thêm, kiểm ràng buộc dịch vụ};
\node[hop,right=of b] (c) {Tính $v$, $e$, $p$ cho từng khách có lời hứa};
\node[hop,right=of c] (d) {Loại phương án vi phạm khóa hay luật lời hứa};
\node[hop,below=9mm of d] (e) {Bộ giải chọn\\một phương án\\mỗi xe};
\node[hop,left=of e] (f) {Kiểm lại cả đội xe; công bố lời hứa $p$; ghi sổ};
\node[ngoai,below=8mm of f] (g) {Sau khi chạy, trên bản ghi đã lưu: lời hứa hiển thị $D = \max(p, S)$};
\draw[mui] (a)--(b); \draw[mui] (b)--(c); \draw[mui] (c)--(d);
\draw[mui] (d)--(e); \draw[mui] (e)--(f);
\draw[mui,dashed] (f.west) -| node[pos=0.3,below,font=\footnotesize]{lượt sau} (a.south);
\draw[-{Latex[length=2.2mm]},dashed] (f.south)--(g.north);
\end{tikzpicture}
\caption[Một lượt quyết định của RideBound]{Một lượt quyết định của RideBound. Bộ điều phối công bố lời hứa không dự báo $p$. Trong các kết quả
chính, lời hứa theo dự báo $D$ được tính sau khi chạy, từ bản ghi đã lưu (ô nét đứt), và không quay lại làm đầu
vào của quyết định nào (mục~\ref{muc:dubao}).}
\label{hinh:vongquyetdinh}
\end{figure}
% EV: X3-03; X3-04; X3-22; X3-25; F03; X3-11 (DISP:1-8); FR:10-11

\begin{table}[htbp]
\centering
\caption[Giả mã một lượt quyết định]{Giả mã một lượt quyết định; dấu $\gets$ nghĩa là gán giá trị bên phải cho biến bên trái}
\label{bang:luotquyetdinh}
\small
\begin{tabular}{@{}r@{\hspace{6pt}}p{10.2cm}@{\hspace{8pt}}p{4.0cm}@{}}
\toprule
\multicolumn{3}{@{}p{14.9cm}@{}}{\textbf{Đầu vào:} trạng thái đội xe; sổ; lô sự kiện mới; luật lời hứa đang
dùng. \textbf{Đầu ra:} kế hoạch mỗi xe; yêu cầu được nhận, để lại hay từ chối; lời hứa mới; bản ghi sổ}\\
\midrule
1: & \textbf{hàm LượtQuyếtĐịnh}(trạng thái, lô sự kiện) & \\
2: & \quad áp lô sự kiện; ghi nhận lần đón muộn hơn giờ đón muộn nhất, nếu có & vị trí xe, đón, trả\\
3: & \quad chọn tối đa 2 yêu cầu đang chờ có giờ đón muộn nhất sớm nhất & yêu cầu khác để lượt sau\\
4: & \quad \textbf{với mỗi} xe $k$ \textbf{thực hiện} & \\
5: & \qquad \textbf{nếu} kế hoạch đang chạy vi phạm ràng buộc cấu trúc \textbf{thì} báo lỗi, dừng lần chạy &
ví dụ vượt sức chứa\\
6: & \qquad $P_k \gets \{$phương án giữ nguyên$\}$, kiểm dịch vụ có nới phần vượt do tắc đường & luôn giữ khi
cắt bớt\\
7: & \qquad thêm vào $P_k$ các phương án chèn yêu cầu ở dòng 3, thỏa chặt ràng buộc dịch vụ; giữ tối đa 100 &
bộ sinh chỉ chèn thêm\\
8: & \qquad \textbf{với mỗi} $\pi \in P_k$ và mỗi khách có lời hứa của xe: tính $v$, $e$, $p$; \textbf{nếu}
$\pi$ vi phạm khóa hay luật lời hứa \textbf{thì} loại $\pi$ & Bảng~\ref{bang:kiemtran}\\
9: & \qquad \textbf{nếu} phương án giữ nguyên bị loại chỉ vì lỗi hạn chót hay lỗi trần \textbf{thì} & xét phục hồi\\
10: & \qquad\quad giữ nó làm phương án được miễn trừ & phục hồi giữ nguyên kế hoạch\\
11: & \qquad\quad giữ thêm phương án đã bị loại không làm vi phạm đó tệ hơn & phục hồi không tệ hơn\\
12: & \qquad \textbf{nếu} phương án giữ nguyên bị loại mà không được miễn \textbf{thì} báo lỗi, dừng lần chạy &
xe kẹt\\
13: & \quad chọn $y$ theo~\eqref{eq:chon} và thứ tự tiêu chí; kiểm lại lựa chọn trên cả đội xe &
Bảng~\ref{bang:rangbuoc}\\
14: & \quad nhận yêu cầu có trong phương án được chọn; để lại hay từ chối yêu cầu khác & \\
15: & \quad khách vừa xác nhận: công bố lời hứa đầu, $L_r \gets 0$; khách có lời hứa mà $x$, $c$ hay $z$ khác 0:
ghi bản ghi sửa, $L_r \gets L_r + |c_r|$ & ghi vi phạm cho khách được miễn\\
16: & \quad \textbf{trả về} quyết định; chờ bộ mô phỏng xác nhận đã áp dụng & \\
17: & \textbf{kết thúc hàm} & \\
\bottomrule
\end{tabular}
\end{table}
% EV (dòng 2–3): X3-01; X3-24 (WP4:8; EventReducer.cs:13-23, 260-282); X3-32 (GEN:54-96)
% EV (dòng 4–7): X3-31 (GEN:171-182); X3-23; X3-27 (CandidatePortfolioRetainer.cs:102-113); X3-03 (GEN:548-575); A24
% EV (dòng 8–12): X3-04; A06; A07; A10; X3-09 (HVA:244-390); X3-31 (POL:336-339; RunnerSession.cs:941-947)
% EV (dòng 13–16): X3-22; X3-07; POL:370-401; X3-20; X3-25 (VALID:272-336, 562-597; RunnerSession.cs:956-987); X3-30; RunnerSession.cs:448-455

Luật lời hứa là bộ lọc cứng: phương án vi phạm bị loại trước khi bộ giải chạy, không được giữ kèm một điểm trừ.
Ngoại lệ là khi xe bị ép (dòng 9--11), chẳng hạn khi riêng việc đường đông hơn đã đẩy giờ trả vượt hạn chót của
luật hạn chót (mục~\ref{muc:doichung}); khi đó tiêu chí 1 chỉ dùng phương án vi phạm nếu xe không còn cách khác.
Bộ sinh không bỏ, không đảo điểm cũ nào, nên khách đã nhận không bị chuyển sang xe khác. Phần nới ở dòng 6 giúp
xe không kẹt vì tắc đường. Mỗi lần vượt như vậy, và mỗi lần khách thực tế lên xe muộn hơn giờ đón muộn nhất, được ghi
lại thay vì làm dừng lần chạy. Vì vậy cửa sổ đón là điều kiện để kiểm phương án, không phải bảo đảm đón đúng hạn.
% EV: A06; HVA:62-66; X3-09 (HVA:255-326); X3-07 (SEL:283-295); A10; X3-03 (GEN:548-575); TEX08:129-131; X3-23 (GEN:161-246; ServiceQualityAllowance.cs:5-18, 61-94; RunnerSession.cs:989-1009); X3-24 (WP4:8; EventReducer.cs:13-23, 260-282)

% -----------------------------------------------------------------------------------------------------
\section{Ba mốc giờ trả và phần dời do điều phối}
\label{muc:bamoc}

Giờ trả có thể đổi vì bộ điều phối ghép thêm khách, hoặc ngay cả khi xe giữ kế hoạch cũ. Để tách hai loại, với
mỗi khách có lời hứa và mỗi phương án, hệ thống so ba mốc giờ trả:
\begin{itemize}[leftmargin=1.2cm,itemsep=2pt,topsep=3pt]
\item $v$: giờ trả đã công bố cho khách ở lần gần nhất;
\item $e$: giờ trả tính lại nếu xe giữ nguyên kế hoạch cũ, từ trạng thái hiện tại với thời gian đi hiện tại;
\item $p$: giờ trả nếu chọn phương án đang xét, với cùng trạng thái và thời gian đi.
\end{itemize}
Ba hiệu sau tính bằng giây (sổ lưu tới mili giây):
\begin{equation}
x = e - v, \qquad c = p - e, \qquad z = p - v = x + c .
\label{eq:bamoc}
\end{equation}
Trong đó, $x$ là phần thay đổi khi giữ nguyên kế hoạch, $c$ là phần dời do điều phối, $z$ là phần khách thấy;
dễ nhớ, $e$ là giờ nếu xe cứ đi như cũ. Mốc $e$ không phải ý của khóa luận: Hörl và Zwick~\cite{horl2022} đã tính
lại giờ trả trên tuyến đang gán theo thời gian đi hiện tại ở mỗi lượt, và khóa luận dùng lại mốc này.
% EV: A01; A02; X3-04 (VALID:355-384); X3-13 (SPEC:39-50); LEDGER:188-207; J01

Mốc $e$ là giờ do mô hình tính, không phải giờ đến thật, nên $x$ gồm mọi thứ đổi khi giữ nguyên kế hoạch: thời
gian đi mới, tiến độ thực của xe, sự kiện vừa xảy ra. Phép tách $z = x + c$ chỉ là phép trừ giữa các giờ của mô
hình, không chứng minh quyết định điều phối thực sự gây ra đúng $c$ giây chậm trễ. Hình~\ref{hinh:bamoc} là một ví
dụ dựng tay: ghép khách M làm A tới muộn hơn lời hứa 70 giây, nhưng chỉ 30 giây là phần dời do điều phối.
% EV: A01 (không phải giờ đến thật, không nhận diện nhân quả); A02; X3-04 (VALID:355-361); K04; ví dụ dựng tay

\begin{figure}[htbp]
\centering
\begin{tikzpicture}[x=1cm,y=1cm,font=\small,
  diem/.style={draw=xamvien,fill=xamnhat,rounded corners=2pt,inner sep=3pt,minimum height=6mm},
  mui/.style={-{Latex[length=2mm]},thick}]
\node[anchor=east] at (2.6,4.0) {Giữ nguyên kế hoạch:};
\node[diem] (x1) at (3.4,4.0) {Xe};
\node[diem] (a1) at (5.6,4.0) {Trả A};
\draw[mui] (x1)--(a1);
\node[anchor=east] at (2.6,3.0) {Ghép thêm khách M:};
\node[diem] (x2) at (3.4,3.0) {Xe};
\node[diem] (b2) at (5.6,3.0) {Đón M};
\node[diem] (a2) at (7.8,3.0) {Trả A};
\node[diem] (c2) at (10.0,3.0) {Trả M};
\draw[mui] (x2)--(b2); \draw[mui] (b2)--(a2); \draw[mui] (a2)--(c2);
\draw[-{Latex[length=2mm]}] (0.8,0)--(12.2,0) node[right]{giờ};
\foreach \vitri/\tenmoc/\giomoc in {2.0/v/9:10:00, 6.8/e/9:10:40, 10.4/p/9:11:10}
  \draw (\vitri,0.12)--(\vitri,-0.12) node[below,align=center]{$\tenmoc$\\\giomoc};
\draw[mui,dashed] (2.0,0.55)--(6.8,0.55) node[midway,above]{$x = e - v = 40$ giây};
\draw[mui] (6.8,0.55)--(10.4,0.55) node[midway,above]{$c = p - e = 30$ giây};
\draw[mui,densely dotted] (2.0,1.45)--(10.4,1.45) node[midway,above]{$z = p - v = 70$ giây};
\end{tikzpicture}
\caption[Ba mốc giờ trả trong một ví dụ dựng tay]{Ví dụ dựng tay: khách A được hứa trả lúc 9:10:00 ($v$); lúc 9:00, xe đang chở A xét ghép thêm khách M.
Tính với thời gian đi hiện tại, giữ nguyên kế hoạch thì A tới lúc 9:10:40 ($e$), ghép M thì tới lúc 9:11:10
($p$). Nét đứt: $x$; nét liền: $c$; nét chấm: $z = x + c$. Luật C chỉ tính $|c|$.}
\label{hinh:bamoc}
\end{figure}

Phần dời do điều phối được tính theo độ lớn: sớm hay muộn hơn $e$ 20 giây đều tính 20 giây, vì luật đo lần sửa
lời hứa chứ không đo lợi hay hại; Fielbaum và Alonso-Mora~\cite{fielbaum2020} thì chỉ tính phần bất lợi. Trên mọi
lần chạy đã ghi, $c \ge 0$ tới mức mili giây; đây là quan sát, chưa phải chứng minh cho trường hợp tổng quát. Trên
tập kiểm tra chính, khóa luận đã kiểm 136 lần chạy của bốn cấu hình: điều phối không trần, luật C trần 30 và 120
giây, luật hạn chót 30 giây. Ở các lần chạy này, mỗi lần kéo giờ trả sớm hơn $e$ chỉ sớm 1 mili giây. Hệ quả: khi
bộ điều phối chỉ chèn thêm, phần tính hai chiều và không hoàn lại của luật C gần như không có việc; nó chỉ có tác
dụng khi kế hoạch của khách đã nhận được phép đổi, ví dụ khi dời khách đang chờ sang xe khác.
% EV: A04; X3-05 (DELTA:137-139, 259-262); J06 (FAM:293-296); N-05 (CORRECTED); ADD1:104-112; ADDENDUM-1 mục 0.1–0.2 (không nêu nguyên nhân của −1 ms)

% -----------------------------------------------------------------------------------------------------
\section{Luật C: trần phần dời giờ trả do điều phối cho từng khách}
\label{muc:luatc}

\begin{dinhnghia}[Luật C]
\label{dn:luatc}
Gọi $B$ là trần, tính bằng giây, và $L_r$ là tổng các $|c|$ đã tính cho khách $r$ kể từ lúc $r$ xác nhận đặt xe,
với $L_r = 0$ lúc xác nhận. Một phương án của một xe chỉ được giữ lại nếu với mọi khách $r$ có lời hứa đang được
gán cho xe đó (chờ đón hay đang trên xe),
\begin{equation}
L_r + |p_r - e_r| \le B .
\label{eq:luatc}
\end{equation}
Sau khi quyết định của lượt được chọn, với mỗi khách $r$ có lời hứa mà phương án được chọn làm $x$, $c$ hay $z$
khác 0, $L_r$ được cộng thêm $|p_r - e_r|$ của phương án đó, kể cả khi giờ trả công bố không đổi ($p_r = v_r$).
\end{dinhnghia}
% EV: A06 (BUDGET:41-105; LEDGER:285-288); A13; X3-05 (BUDGET:107-114; CFG C-30.json:17-23); X3-25 (VALID:272-281, 562-597)

Dấu $\le$ cho phép $L_r + |c|$ đúng bằng $B$. Phương án chỉ được xét không làm đổi $L_r$: mỗi phương án được
kiểm trên bản sao của sổ, chỉ quyết định cuối của lượt mới được ghi vào sổ thật. $B$ giống một ngân sách thời
gian của từng khách: $B - L_r$ là phần trần còn lại, không được cấp lại sau mỗi lần dời. Vì phần dời được cộng
qua cả chuyến, một luật cho mỗi lần dời tới 30 giây cho phép ba lần, tổng 90 giây, còn luật C với $B = 30$ giây thì
không. Ở Hình~\ref{hinh:bamoc}, với $L_A = 0$ và $B = 30$ giây, phương án ghép M ($|c| = 30$ giây) vẫn được giữ và
A dùng hết trần. Khi phép kiểm không được miễn cho $r$ ở lượt nào, tổng $|c|$ của $r$ cả chuyến không vượt $B$.
% EV: X3-05 (BUDGET:88); X3-25 (HVA:161-174; RunnerSession.cs:956-987; VALID:640-644: miễn trừ vẫn cộng vào sổ); A06; K43; E08; J05; CFG C-30.json:3-6 (một chính sách cho mọi khách)

\begin{table}[htbp]
\centering
\caption[Giả mã phép kiểm luật C cho một phương án]{Giả mã phép kiểm luật C cho một phương án (chỉ viết chiều giờ trả)}
\label{bang:kiemtran}
\small
\begin{tabular}{@{}r@{\hspace{6pt}}p{9.6cm}@{\hspace{8pt}}p{4.6cm}@{}}
\toprule
\multicolumn{3}{@{}p{14.9cm}@{}}{\textbf{Đầu vào:} phương án $\pi$ của xe $k$; trần $B$; $L_r$ của mỗi khách
có lời hứa; trạng thái và thời gian đi của lượt. \textbf{Đầu ra:} giữ hay loại $\pi$, kèm bản ghi chặn}\\
\midrule
1: & \textbf{hàm KiểmTrần}($\pi$, $k$, $B$) & \\
2: & \quad \textbf{với mỗi} khách $r$ có lời hứa được gán cho xe $k$, theo thứ tự mã khách \textbf{thực hiện} &
chờ đón hay trên xe\\
3: & \qquad $e_r \gets$ giờ trả của $r$ khi chiếu lại kế hoạch cũ của $k$ từ trạng thái hiện tại & mốc giữ
nguyên kế hoạch\\
4: & \qquad $p_r \gets$ giờ trả của $r$ theo $\pi$, cùng trạng thái và thời gian đi; $c_r \gets p_r - e_r$ &
phần dời do điều phối\\
5: & \qquad \textbf{nếu} $L_r + |c_r| > B$ \textbf{thì trả về} loại $\pi$, kèm ($r$, $L_r$, $|c_r|$,
$L_r + |c_r|$, $B$) & dừng ở khách vi phạm gặp trước\\
6: & \quad \textbf{trả về} giữ $\pi$ & $L_r$ chỉ đổi nếu $\pi$ được chọn\\
7: & \textbf{kết thúc hàm} & \\
\bottomrule
\end{tabular}
\end{table}
% EV (dòng 2–4): X3-04 (VALID:233-239, 355-361; HVA:171); X3-27 (VALID:1048-1064)
% EV (dòng 5–6): X3-05 (BUDGET:6-12, 72-104); A06; A23; X3-25

Cài đặt theo dõi mười chiều cam kết: thay đổi giờ đón, giờ trả; số lần sửa đáng kể; đổi xe; khoảng dời và số lần
đổi của điểm đón, của điểm trả; số lần đảo thứ tự với khách cũ; số điểm dừng mới chèn trước điểm đón. Chỉ bốn
chiều có trần: giờ trả ($B$) và đổi xe, đổi điểm đón, đổi điểm trả (trần 0). Ở tập kiểm tra chính và tập buổi sáng,
phép kiểm dừng ở khách vi phạm gặp trước, nên mỗi bản ghi chặn (bản ghi lý do loại phương án) chỉ nêu lỗi của khách
đó.
% EV: A04 (DIM:3-32); X3-26 (DELTA:137-153); A25; A07; CFG C-30.json:7-89; X3-07 (SEL:319-330, 510-522); A23; X3-04 (VALID:222-239); WP4:1-25; X3-10

\subsection{Phương án giữ nguyên kế hoạch dưới luật C}
\label{muc:giunguyen}

Gọi \textbf{(N)} là giả thiết, áp cho lượt đang xét: (i) kế hoạch đang chạy không vi phạm ràng buộc cấu trúc
(sức chứa, đón trước trả, phần kế hoạch đã cố định), nên bộ sinh dựng được phương án giữ nguyên; (ii) giờ trả
theo phương án giữ nguyên được tính bằng đúng phép chiếu dùng cho $e$. RideBound thỏa (ii) vì dùng cùng một hàm
chiếu với cùng đầu vào, và luôn giữ phương án giữ nguyên khi cắt còn 100 phương án.
% EV: A22; X3-31 (GEN:171-182, 213-229); ServiceQualityAllowance.cs:66-73; X3-27 (VALID:355-360, 1048-1064; HVA:144-165; CandidatePortfolioRetainer.cs:102-113; WP4:15)

\begin{nhanxet}
\label{nx:giunguyen}
Giả sử (N) và mọi khách có lời hứa của xe có $0 \le L_r \le B$. Khi đó luật C không loại phương án giữ nguyên
kế hoạch.
\end{nhanxet}

\noindent\textit{Chứng minh.} Theo (N)(ii), với mọi khách $r$ có lời hứa của xe, $p_r = e_r$, nên
$L_r + |p_r - e_r| = L_r \le B$ và phép kiểm~\eqref{eq:luatc} cho qua. \hfill$\square$
% EV: A22 (theo cấu tạo, giả thiết N); A06; X3-27

Khóa và các trần 0 cũng so với kế hoạch giữ nguyên, nên phương án này qua cả chúng. Nếu (N) đúng ở mọi lượt thì
$L_r \le B$ được giữ: $L_r = 0$ lúc xác nhận, và nếu $L_r \le B$ ở đầu một lượt thì không có miễn trừ nào, phương
án được chọn đã qua~\eqref{eq:luatc}, nên $L_r$ mới vẫn không vượt $B$. Khi đó xe dùng luật C hay điều phối không
trần không bao giờ bị ép; tính chất này vì vậy không riêng của luật C.
% EV: A10 (LOCK:156-161); X3-06; X3-26 (véc-tơ c của phương án giữ nguyên bằng 0 ở mọi chiều); X3-09 (HVA:255-259); X3-25

Nhận xét đúng do cách đặt luật, không phải kết quả thực nghiệm; số liệu chỉ phù hợp với nó: trên tập kiểm tra
chính, 0 trong 20.101 quyết định bị ép với trần 30 giây và 0 trong 20.435 với trần 120 giây (mỗi quyết định ứng với
một lượt của cả đội xe).
% EV: A21; N-21; X3-16 (STG:43-44); N-07; X3-30 (RunnerSession.cs:1013-1027); SPEC:31-34

Nhận xét có ba giới hạn. Thứ nhất, nếu (N)(i) không thỏa, lần chạy dừng vì lỗi cấu trúc, không vì luật C; phần
vượt cửa sổ đón hay thời gian ngồi xe do tắc đường không làm (N) hỏng, nhờ phần nới. Thứ hai, với $v$ = 9:10:00 và
$e = p$ = 9:11:00, phương án giữ nguyên có $c = 0$ nhưng $z = 60$ giây: luật C không bảo đảm giờ đến. Thứ ba, với
$v = p$ = 9:10:00 và $e$ = 9:12:00, kéo giờ trả về đúng giờ đã hứa có $z = 0$ nhưng $|c| = 120$ giây, và $L_r$
tăng 120 giây dù giờ công bố không đổi. Tính chất và giới hạn thứ hai là hai mặt của một lựa chọn: luật C không
bị ép vì không tính phần thay đổi khi giữ nguyên kế hoạch, và vì thế cũng không giới hạn phần đó.
% EV: A22; X3-31 (GEN:171-182; RunnerSession.cs:941-947); X3-23; K05; X3-25 (VALID:562-597: bản ghi sửa khi x ≠ 0 dù p = v); ví dụ dựng tay

% -----------------------------------------------------------------------------------------------------
\section{Các luật đối chứng, bộ chọn và luật phục hồi}
\label{muc:doichung}

Bảng~\ref{bang:luat} định nghĩa bốn cấu hình mà Chương~\ref{chuong:thucnghiem} so với luật C. Trừ ghép tự do,
chúng dùng cùng bộ sinh, bộ kiểm, bộ chọn và luật phục hồi với luật C; tập tắc đường cố định chạy bằng bộ chạy
chưa có luật phục hồi.
% EV: A08; A09; A11; A17; A19; K50; HO:109; B17; NEW:13

\begin{table}[htbp]
\centering
\caption[Các cấu hình được so sánh]{Các cấu hình được so sánh; ký hiệu chỉ dùng trong bảng. $B$ là trần, tính bằng giây; với luật hạn chót,
$B$ là số giây giờ trả được phép muộn hơn lời hứa đầu $p_0$. Các tập thử được mô tả ở mục~\ref{muc:c4-thietlap}.}
\label{bang:luat}
\footnotesize
\setlength{\tabcolsep}{3pt}
\begin{tabular}{@{}p{2.2cm}p{1.1cm}p{3.1cm}p{2.3cm}p{2.6cm}p{3.5cm}@{}}
\toprule
\textbf{Tên trong khóa luận} & \textbf{Ký hiệu} & \textbf{Kiểm điều gì} & \textbf{Khóa} & \textbf{Khi xe bị
ép} & \textbf{Dùng ở tập nào}\\
\midrule
Điều phối không trần & U & Không kiểm giờ trả & Xe, điểm đón, giờ đón; điểm trả có trần 0 & Không bị ép khi (N)
đúng & Tập kiểm tra chính; tập chỉnh trần buổi sáng; tập tắc đường cố định\\
% EV: A07; A08; K50; A10; X3-16 (STG:42: 0 / 20.760); B10; B15; B17; NEW:11
Ghép tự do & FREE & Không kiểm lời hứa; chọn theo chi phí kế hoạch ở từng lượt & Không khóa & Không áp dụng &
Tập tắc đường cố định\\
% EV: A17; K50; HO:109 (b1-rolling-cost, không khóa); B17; X3-38
Luật C & C-$B$ & $\sum |c| \le B$ cả chuyến, cho mọi khách có lời hứa & Như U & Không bị ép khi (N) đúng
(Nhận xét~\ref{nx:giunguyen}) & Tập chính (30, 120); buổi sáng (30, 60, 90, 120, 150, 180, 240, 300); tắc đường
cố định (30, 60, 120); chiều tuyết (30)\\
% EV: A06; A22; B10; L1a-01; X3-28 (run_w07.py:5; run_w07x.py:3-4, 35-36); B17; B20; N-26 (chiều tuyết 0 / 5.067 quyết định bị ép)
Luật hạn chót & M$^{+}$-$B$ & $p - p_0 \le B$ ở mọi lần kiểm; chỉ chặn phía muộn & Như U & Có thể bị ép do tắc
đường; phục hồi không tệ hơn & Tập chính (30, 60); buổi sáng (tám mức như luật C)\\
% EV: A09; A10; A19; B10; X3-28; D04; LOCK:186-197; CFG Mplus-30.json:98
Luật trần độ dời khách thấy & V-$B$ & $\sum |z| \le B$ cả chuyến & Như U & Có thể bị ép do tắc đường; phục hồi
không tệ hơn & Tập chính (30); buổi sáng (30, 90, 120, 150, 180, 240)\\
% EV: A11; A19; B10; X3-28 (run_w07x.py:3-4, 35-36: V-90/120/150/180/240; run_w07.py:5: V-30)
\bottomrule
\end{tabular}
\end{table}

Luật hạn chót chỉ mượn thành phần hạn chót tính từ lời hứa đầu của Schulz và Pfeiffer~\cite{schulz2026}, vốn
dùng thời gian đi cố định và độ chùng mặc định 900 giây. Ở đây mỗi độ chùng bằng một mức trần của luật C trên cùng
tập, và với tắc đường đổi theo khung 15 phút, riêng tắc đường có thể làm luật này bị ép; số quyết định bị ép ở
Chương~\ref{chuong:thucnghiem} nói về một hạn chót chặt như vậy, không về phương pháp của họ.
% EV: J10 (SP:466-490, 568-577); J11; A08; A09; X3-28; X3-41 (D03:214-217); A10

Ba luật lời hứa khác nhau ở đại lượng được cộng. Đánh số các lần cập nhật lời hứa của một khách là
$1, \ldots, n$. Trong sổ, mỗi bản ghi sửa so với lời hứa ngay trước nó ($v_j = p_{j-1}$), nên các giờ ở giữa
triệt tiêu:
\begin{equation}
\sum_{j=1}^{n} (x_j + c_j) \;=\; \sum_{j=1}^{n} z_j \;=\; p_n - p_0 .
\label{eq:dongnhat}
\end{equation}
Luật hạn chót kiểm $p_n - p_0 \le B$, tức cộng có dấu cả hai phần; luật trần độ dời khách thấy cộng $|z_j|$; luật
C chỉ cộng $|c_j|$. Ở Hình~\ref{hinh:bamoc}, nếu 9:10:00 là lời hứa đầu, phương án giữ nguyên có
$e - p_0 = |z| = 40$ giây, vượt trần 30 giây của hai luật kia dù $c = 0$: riêng tắc đường đủ làm chúng bị ép.
% EV: A12; X3-33 (LEDGER:256-269, 298-314; SPEC:62-63); A09; A10 (LOCK:156-161); A11; X3-16 (U 0 / 20.760)

Tiền lệ trực tiếp là Hörl và Zwick~\cite{horl2022} (mục~\ref{muc:c2-rangbuoc}): ở đầu mỗi lượt họ tính mốc $e$,
và công thức 7 đặt giờ trả bắt buộc bằng giá trị lớn hơn giữa giờ cập nhật và giờ trả muộn nhất, như ràng buộc cứng
(ràng buộc 9--10), kèm phạt cho phần vượt giờ gốc (công thức 11). Mục đích của họ, tránh tình huống không còn thứ tự
trả khách khả thi, cũng là mục đích của Nhận xét~\ref{nx:giunguyen}. Đặt giờ trả muộn nhất là $p_0 + B$ cho hạn chót
nới $p \le \max(p_0 + B, e)$, một cách viết lại của khóa luận; nó cho phép $c \le \max(p_0 + B - e, 0)$, về 0 khi
tắc đường đẩy $e$ quá $p_0 + B$, và cũng không loại phương án giữ nguyên. Cấu hình hạn chót M$^{+}$ với luật phục
hồi ``không tệ hơn'' ở chương~4 chính là dạng này. Luật C khác ở ba chỗ:
một trần cho mọi khách có lời hứa, còn phần nới của họ chỉ tác dụng khi giờ cập nhật đã vượt giờ trả muộn nhất;
phần được phép dời, $|c| \le B - L_r$, không phụ thuộc tắc đường; và cộng qua cả chuyến, dù phần cộng dồn hiếm khi là
lý do loại phương án (mục~\ref{muc:c4-congdon}).
% EV: J01; J02 (HZ:481-497); J03; X3-36 (HZ:531-543, 565-606); J05; A22 (suy luận V17 R6); K02; K26; K27; N-15

Từ Fielbaum và Alonso-Mora~\cite{fielbaum2020} (mục~\ref{muc:c2-dolech}), khóa luận kế thừa ý tách phần thay đổi
do quyết định gán xe và kiểm soát nó trước khi quyết định, nhưng mỗi lần sửa được so với mốc $e$ chứ không so điểm
đầu với điểm cuối, nên hai lần dời ngược chiều không triệt tiêu nhau; chiều sớm cũng được tính; đại lượng là giờ
trả chứ không phải thời gian chờ; và tắc đường thay đổi chứ không chọn giờ thấp điểm.
% EV: J06; J07; J08; X3-35 (FAM:90-94, 189-195, 216-220, 293-296, 948-972); K30; K32

Vì đổi xe, đổi điểm đón, đổi điểm trả là các trần 0, phần đầu của tiêu chí 3 (mức dùng trần lớn nhất, tỉ số lớn
nhất giữa phần đã dùng và trần) luôn bằng 100\% ở các lượt đã ghi của tập kiểm tra chính (mục~\ref{muc:c4-hople}),
cả khi không trần; tiêu chí 3 vì vậy chỉ tác dụng qua tổng phần dời theo từng chiều. Điều phối không trần và luật C
cùng cấu hình và bộ chọn, chỉ khác trần giờ trả; chênh lệch giữa chúng là phần đi kèm việc thêm trần khi giữ nguyên
bộ chọn. Khóa luận còn xét thứ tự ưu tiên chi phí: giữ ưu tiên nhận yêu cầu và các ràng buộc, bỏ các mức xếp
hạng theo ổn định rồi so chi phí kế hoạch và tiêu chí phân xử. Phương án vẫn phải qua phép kiểm trần khi trần
được bật. Mục~\ref{muc:c4-vaitro} so hai thứ tự ở trần 30 giây và không trần; đổi thứ tự không làm thay đổi
định nghĩa hay cách cộng trần của luật C.
% EV: S15-03 (costfirst-diff, cấu hình và provenance); S15-01 (mục 0, 2).
% EV: A14; X3-07 (SEL:271-330; HVA:537-636); A15; N-25; A08; H01; H02; K18

Luật phục hồi giữ nguyên kế hoạch giữ phương án giữ nguyên làm phương án được miễn trừ nếu mọi lỗi của nó là lỗi
hạn chót hay lỗi trần. Luật phục hồi không tệ hơn mở rộng luật này: giữ thêm phương án đổi kế hoạch có mỗi lỗi,
cùng khách và cùng chiều, cũng là lỗi của phương án giữ nguyên và không nặng hơn. Lỗi khóa theo giai đoạn (khóa của
khách chờ đón hay trên xe) và lỗi ràng buộc dịch vụ không bao giờ được miễn. Mỗi khách được miễn được ghi một vi
phạm, trừ khi lời hứa, giai đoạn và mã lỗi không đổi so với lần trước. Tập kiểm tra chính, tập buổi sáng và tám
lần chạy chiều tuyết dùng luật thứ hai; cả hai là chi tiết cài đặt, không phải đóng góp của khóa luận.
% EV: A17; A18 (ADR:5066-5092); X3-09 (HVA:79-99, 255-326; VALID:413-420, 475-483, 537-544, 599-638, 791-818; LOCK:42-60); A19 (RUN38:9; WP4:7); X3-10; R26-CH3-OPEN-QUESTIONS-CODEX.md §Q8 (tám lần chạy chiều tuyết; wp4-tier3-v2.json:7)

% -----------------------------------------------------------------------------------------------------
\section{Lời hứa theo dự báo}
\label{muc:dubao}

Vì bộ điều phối dùng thời gian đi hiện tại, lời hứa có thể sớm hơn giờ trả thực tế khi đường đông dần, dù xe giữ
nguyên kế hoạch. Lời hứa theo dự báo, một cách làm đã có (mục~\ref{muc:c2-dubao}), cộng sẵn mức tăng tắc đường các
ngày trước cho thấy vào từng giờ đón, giờ trả được công bố.
% EV: A05; F10; J35; F01; X3-11 (DISP:1-8)

Thời gian đo bằng giây kể từ đầu khung giờ của kịch bản; $i(\tau) = \lfloor \tau / 900 \rfloor$ là khung 15 phút
chứa $\tau$. Gọi $m_i$, $i = 0, \ldots, 8$, là trung bình hệ số tắc đường của khung $i$ trên các ngày nguồn; từ
khung 8 hệ số được giữ ở giá trị cuối, như bộ mô phỏng cũng giữ, nên đặt $\tilde m_i = m_{\min(i, 8)}$. Với lúc
công bố $t$ và giờ công bố $p > t$ (giờ của phương án được chọn, chưa cộng dự báo), $S$ là nghiệm của
\begin{equation}
\int_{t}^{S} \frac{\mathrm{d}\tau}{g(\tau)} = p - t, \qquad
g(\tau) =
\begin{cases}
1 & \text{nếu } i(\tau) = i(t),\\
\tilde m_{i(\tau)} / \tilde m_{i(t)} & \text{nếu không,}
\end{cases}
\label{eq:keodan}
\end{equation}
và lời hứa hiển thị là $D = \max(p, S)$ (giả mã ở Bảng~\ref{bang:dubao}).
% EV: F01; X3-11 (DISP:3-8, 55-64; STRETCH:1-12, 22-52; preflight_world_v1_3.py:51-82)

Theo tốc độ lúc công bố, xe cần thêm $p - t$ giây. Trong khung chậm hơn khung hiện tại $g$ lần, mỗi giây đồng hồ chỉ
đi được phần đường mà ở tốc độ lúc công bố mất $1/g$ giây. Mức tắc đường lúc này đã có trong thời gian đi đang dùng, nên $g = 1$ ở khung
hiện tại; các ngày khác chỉ cho biết khung sau chậm hay nhanh hơn bao nhiêu lần. Phép lấy $\max$ bảo đảm $D \ge p$ ở
cùng lần công bố, không bảo đảm $D$ tăng dần qua các lần cập nhật. Ví dụ dựng tay: trong kịch bản khung
8:00--10:00, lúc $t$ = 9:05:00 hệ thống công bố $p$ = 9:25:00, tức còn 1.200 giây. Nếu các ngày nguồn có hệ số trung
bình 1,25 ở khung 9:00--9:15 và 1,50 ở khung 9:15--9:30, thì sau 9:15 $g = 1{,}50/1{,}25 = 1{,}2$. Tới 9:15 xe đi
được 600 giây quãng đường; 600 giây còn lại mất $600 \times 1{,}2 = 720$ giây, nên $S = D$ = 9:27:00, dài thêm 120
giây.
% EV: X3-11 (STRETCH:1-12; forecast_weight neo ở khung hiện tại); F01; DISP:60-64; ví dụ dựng tay X3-37 (a), tính bằng hàm stretch của dự án; số học ở 07-chuong3-v3-NOTES.md mục 4

\begin{table}[htbp]
\centering
\caption[Giả mã lời hứa theo dự báo]{Giả mã lời hứa theo dự báo cho một giờ đón hay giờ trả đã công bố}
\label{bang:dubao}
\small
\begin{tabular}{@{}r@{\hspace{6pt}}p{9.6cm}@{\hspace{8pt}}p{4.6cm}@{}}
\toprule
\multicolumn{3}{@{}p{14.9cm}@{}}{\textbf{Đầu vào:} lúc công bố $t$; giờ công bố $p$; ngày $d$ và khung giờ của
kịch bản; hệ số tắc đường theo khung của các ngày trước $d$. \textbf{Đầu ra:} lời hứa hiển thị $D$}\\
\midrule
1: & \textbf{hàm LờiHứaTheoDựBáo}($t$, $p$, $d$) & \\
2: & \quad $H \gets$ các ngày cùng loại với $d$ và trước $d$, bỏ ngày 15/11 & chỉ dùng ngày đã qua\\
3: & \quad \textbf{nếu} $H$ có ít hơn 2 ngày \textbf{hoặc} $p \le t$ \textbf{thì trả về} $p$ & không dự báo\\
4: & \quad $m_i \gets$ trung bình hệ số của khung $i$ trên $H$, $i = 0, \ldots, 8$ & khung tính từ đầu khung giờ\\
5: & \quad $W \gets p - t$;\; $\tau \gets t$;\; $i_0 \gets i(t)$ & quãng còn lại theo tốc độ hiện tại\\
6: & \quad \textbf{lặp} & \\
7: & \qquad $g \gets 1$ nếu $i(\tau) = i_0$, ngược lại $g \gets \tilde m_{i(\tau)} / \tilde m_{i_0}$ &
$\tilde m_i = m_{\min(i, 8)}$\\
8: & \qquad $h \gets$ thời gian từ $\tau$ tới hết khung $i(\tau)$ & \\
9: & \qquad \textbf{nếu} $h/g \ge W$ \textbf{thì} $S \gets \tau + W g$ và thoát vòng lặp & \\
10: & \qquad $W \gets W - h/g$;\; $\tau \gets$ đầu khung kế tiếp & \\
11: & \quad \textbf{trả về} $D \gets \max(p, S)$ & $D \ge p$ ở lần công bố này\\
12: & \textbf{kết thúc hàm} & \\
\bottomrule
\end{tabular}
\end{table}
% EV (dòng 2–4): F02; X3-11 (DISP:27-57; preflight_world_v1_3.py:51, 80-82)
% EV (dòng 5–11): F01; X3-11 (STRETCH:25-39, 49-52)

Ngày nguồn là các ngày cùng loại (ngày thường hay cuối tuần) trước ngày của kịch bản, ít nhất 2 ngày, và không bao
giờ là ngày tuyết 15/11/2018 (một lựa chọn của khóa luận; hệ thống thật cần quy tắc tự động nhận ra ngày bất
thường). Vì vậy trong dữ liệu 12--18/11/2018 chỉ ba ngày 14, 15, 16/11 có dự báo: 21 trong 34 kịch bản của tập
kiểm tra chính; 13 kịch bản ngày 17, 18/11 giữ lời hứa không dự báo.
% EV: F02 (DISP:10-16, 27-52); F04; N-16; N-28

Trong các kết quả chính, $D$ được tính lại trên các lần chạy đã ghi, gọi là phát lại (replay): xe nào chở ai giữ
nguyên, chỉ giờ báo cho khách đổi, và $D$ không quay lại sổ hay phép kiểm trần. $D_0 - p_0$, với $D_0$ là lời hứa
đầu có dự báo, là giá của dự báo. Phép kiểm khả thi vẫn dùng giờ không dự báo, nên dưới luật C trần 30 giây, trong
1.471 khách được phục vụ ở 21 kịch bản có dự báo, 34 khách có giờ đón hiển thị ở lời hứa đầu muộn hơn giờ đón
muộn nhất; chỗ này cần sửa trước khi đưa dự báo vào bộ điều phối. Phép thử đưa dự báo vào bộ điều phối
(mục~\ref{muc:c4-dubao}) thì khác: nó đổi thời gian đi mà bộ điều phối nhận, còn FleetPy vẫn chạy theo hệ số đo.
Số khung lấy trung bình là $K$: các cấu hình chính bật dự báo dùng $K = 2$, gồm khung 15 phút hiện tại và khung kế tiếp.
Phép thử giả dùng $K = 1$, cho hệ số bằng 1; khi tắt dự báo không đặt $K$.
% EV: F03 (chỉ phần "phát lại trên các lần chạy đã ghi"); FR:10-11; X3-12 (FFA:84-99); F06; F07 (FR:70); F12; X3-29 (preflight_world_v1_3.py:1-26, 113-133); R26-CH3-OPEN-QUESTIONS-CODEX.md §Q9; preflight_world_v1_3.py:85-89, 99-101; run_forecast.py:67-72

% -----------------------------------------------------------------------------------------------------
\section{Sổ theo từng khách}
\label{muc:so}

Sổ do RideBound ghi trong lúc chạy: mỗi khách có bản ghi lời hứa đầu và một bản ghi sửa cho mỗi lượt mà phương án
được chọn làm $x$, $c$ hay $z$ khác 0 ở bất kỳ chiều nào. Bản ghi sửa lưu các phép chiếu ứng với $v$, $e$, $p$,
độ lớn thay đổi trên mười chiều và $L$ sau lần sửa (Bảng~\ref{bang:vidu-so}); các hiệu có dấu $x$, $c$, $z$ chỉ dùng
cho giờ. Không có ngưỡng: thay đổi 1 mili giây cũng tạo bản ghi sửa, nên số bản ghi sửa không phải số lần một ứng
dụng thật báo lại cho khách.
% EV: A04; X3-26 (DELTA:64-160); X3-25 (VALID:562-597); LEDGER:285-314; X3-13 (SPEC:35-46)

\begin{table}[htbp]
\centering
\caption[Bản ghi sổ của khách A trong ví dụ dựng tay]{Hai bản ghi giờ trả của khách A ở Hình~\ref{hinh:bamoc} (dựng tay), với giả định 9:10:00 là lời hứa đầu,
công bố lúc 8:50:00, và $B = 30$ giây. Bản ghi lời hứa đầu không có lời hứa trước nên sổ quy ước $v = e = p$. Các
độ lớn và $L$ tính bằng giây.}
\label{bang:vidu-so}
\small
\begin{tabular}{@{}lllllrrrr@{}}
\toprule
\textbf{Bản ghi} & \textbf{Lúc ghi} & $v$ & $e$ & $p$ & $|x|$ & $|c|$ & $|z|$ & $L$\\
\midrule
Lời hứa đầu & 8:50:00 & 9:10:00 & 9:10:00 & 9:10:00 & 0 & 0 & 0 & 0\\
Sửa (ghép khách M) & 9:00:00 & 9:10:00 & 9:10:40 & 9:11:10 & 40 & 30 & 70 & 30\\
\bottomrule
\end{tabular}
\end{table}
% EV: ví dụ dựng tay; v = e = p ở bản ghi đầu theo SPEC:46; L = 0 + |c| theo LEDGER:285-288

Trong một bản ghi mô phỏng của luật C trần 120 giây, một khách có $v$ = 17:43:46,044, $e$ = 17:43:46,283 và
$p$ = 17:44:01,203. Phần dời $|p - e| = 14{,}92$ giây làm phần trần còn lại bằng $120 - 14{,}92 = 105{,}08$ giây.
Cả 14 bản ghi sửa sau có $c = 0$ và $|x| < 0{,}25$ giây, nên phần trần đã dùng không tăng.
% EV: round26-trace-rider.txt, Q3 SELECTED RIDER, entries 1-15; R26-CH3-OPEN-QUESTIONS-CODEX.md §Q3; C-120-d20181117-s10-r1-w17/transcript-00.ndjson, final checkpoint ordinal 1716, request rq-c48d442603e9ce8c056b029d6d50c497f4d224b8af469945692690e92a75afb0

Sau khi chạy, độ trễ so với lời hứa đầu được chia ba phần: phần tính cho điều phối, phần được hồ sơ lịch sử giải
thích và phần còn lại,
\begin{equation}
a - p_0 \;=\; q \;+\; f \;+\; (a - D_0 - q), \qquad q = \sum_{j} c_j, \qquad f = D_0 - p_0 .
\label{eq:so}
\end{equation}
Cộng ba phần thì $q$ và $D_0$ triệt tiêu, nên đẳng thức luôn đúng. Trong đó, $a$ là giờ trả thực tế trong mô
phỏng; $q$ cộng các $c$ giữ nguyên dấu (dời muộn dương, dời sớm âm), khác với $\sum |c|$ của luật C, và dưới luật
C, $|q| \le \sum |c| \le B$; $f$ là phần dự báo đã cộng sẵn vào lời hứa đầu, bằng 0 khi không có dự báo.
Theo~\eqref{eq:dongnhat}, phần còn lại bằng $\sum_j x_j + \rho - f$ với $\rho = a - p_n$: thay đổi khi giữ nguyên
kế hoạch mà dự báo chưa tính, cộng độ lệch giữa giờ trả thực tế và lời hứa cuối $p_n$. Với khách A ở
Bảng~\ref{bang:vidu-so}, nếu $D_0$ = 9:10:30 và A tới đúng lời hứa cuối, $a$ = 9:11:10, thì độ trễ 70 giây gồm
$q = 30$, $f = 30$ và 10 giây còn lại, đúng bằng $\sum x + \rho - f = 40 + 0 - 30$.
% EV: X3-12 (FFA:84-97: q = Σ(p − e) có dấu, D0 = lời hứa hiển thị đầu); A12; A06; F09; ví dụ dựng tay X3-37 (b) (số học ở NOTES mục 4)

Chương~\ref{chuong:thucnghiem} cộng ba phần trên các khách có $a - p_0 > 60$ giây, với $p_0$ là lời hứa giờ trả
đầu không dự báo (không dùng $D_0$). Đây là phép tính theo mô hình, không kết luận ai có lỗi: ``phần tính cho điều
phối'' chỉ nói rằng mô hình gán phần này cho quyết định điều phối khi so với kế hoạch giữ nguyên. Phần được hồ sơ
lịch sử giải thích phụ thuộc dự báo đang dùng; phần còn lại là phần dự báo đó chưa giải thích được, không phải phần
``không thể biết trước''.
% EV: X3-12 (FFA:134-138: nhóm A − P0 > 60 s); F09; K04; A01 (lưu ý D03:222-224); K19

% -----------------------------------------------------------------------------------------------------
\section{Cài đặt, kiểm thử và tái lập}
\label{muc:caidat}

\subsection{Kiến trúc và trao đổi với bộ mô phỏng}
\label{muc:kientruc}

RideBound viết bằng C\# trên .NET 10, chia thành các dự án theo tầng (Bảng~\ref{bang:kientruc}). Tầng nghiệp vụ
không tham chiếu dự án hay gói thư viện ngoài nào, OR-Tools chỉ có trong cầu nối bộ giải, và một bộ kiểm thử kiến
trúc kiểm các quy tắc này; nhờ vậy phép kiểm trần, ba mốc và sổ không phụ thuộc bộ mô phỏng hay bộ giải.
% EV: X3-14 (các tệp .csproj: net10.0, ProjectReference, PackageReference; DependencyRuleTests.cs:8-31, 137-166)

\begin{table}[htbp]
\centering
\caption[Các thành phần chính của RideBound và bộ chuyển đổi FleetPy]{Các thành phần chính của RideBound và bộ chuyển đổi FleetPy}
\label{bang:kientruc}
\small
\begin{tabular}{@{}p{4.6cm}p{10.2cm}@{}}
\toprule
\textbf{Thành phần} & \textbf{Vai trò}\\
\midrule
Domain (tầng nghiệp vụ) & Yêu cầu và giai đoạn của khách, lời hứa, khóa, sổ, phép kiểm trần\\
Application (tầng ứng dụng) & Chiếu lịch theo thời gian đi, tính ba mốc giờ trả, kiểm phương án theo khóa và luật
lời hứa\\
Algorithms (thuật toán) & Bộ sinh chỉ chèn thêm, đánh giá phương án theo luật lời hứa và luật phục hồi, bộ chọn\\
Solvers.OrTools (cầu nối bộ giải) & Cầu nối tới bộ giải CP-SAT của OR-Tools 9.15\\
Contracts (định dạng bản tin) & Định dạng các bản tin trao đổi với bộ mô phỏng\\
Runner (bộ chạy) & Nhận bản tin, ghép các thành phần trên và trả quyết định\\
Benchmarking (chuẩn bị dữ liệu) & Chuẩn hóa dữ liệu Manhattan của FleetPy thành kịch bản thử\\
Bộ chuyển đổi FleetPy (Python) & Phần điều khiển đội xe của FleetPy 1.0.2: gửi sự kiện cho bộ chạy, áp quyết định
vào mô phỏng\\
\bottomrule
\end{tabular}
\end{table}
% EV: X3-14 (Google.OrTools 9.15.6755 chỉ ở RideBound.Solvers.OrTools.csproj:9); X3-22 (OrTools:3, CP-SAT); X3-15; NORM; B01

Bộ chuyển đổi viết bằng Python đóng vai phần điều khiển đội xe của FleetPy 1.0.2~\cite{engelhardt2026}, ghim theo
phiên bản; nó chạy bộ chạy RideBound như một chương trình con, hai bên nhắn tin qua đầu vào, đầu ra chuẩn, mỗi bản
tin một dòng JSON (Hình~\ref{hinh:bantin}). Bộ chạy chỉ nhận lô mới sau khi bộ mô phỏng báo đã áp dụng quyết định
trước, nên các quyết định được áp dụng đúng thứ tự. Bản tin của mỗi lần chạy được lưu, và mọi con số ở
Chương~\ref{chuong:thucnghiem} được tính từ đó.
% EV: B01; J18 (FLEET:474); X3-15 (README:1-12; runner_client.py:133-138; Program.cs:103-106; RunnerSession.cs:448-455); X3-13 (SPEC:28-37); A26

\begin{figure}[htbp]
\centering
\begin{tikzpicture}[y=1cm,font=\small,
  ben/.style={draw=xamvien,fill=xamnhat,rounded corners=2pt,align=center,inner sep=4pt,text width=4.3cm},
  tin/.style={-{Latex[length=2mm]},thick},
  nhan/.style={midway,above,font=\footnotesize,align=center}]
\node[ben] (L) at (0,0) {FleetPy 1.0.2\\và bộ chuyển đổi (Python)};
\node[ben] (R) at (9.5,0) {Bộ chạy RideBound\\(C\#, .NET 10)};
\draw[xamvien,thick] (L.south)--(0,-12.5);
\draw[xamvien,thick] (R.south)--(9.5,-12.5);
\draw[tin] (0,-1.8)--(9.5,-1.8) node[nhan]{chào hỏi; khởi tạo lần chạy\\(kịch bản, cấu hình luật)};
\draw[tin] (9.5,-2.9)--(0,-2.9) node[nhan]{đã nhận chào hỏi; đã khởi tạo};
\draw[dashed,xamvien] (-0.6,-3.35) rectangle (12.2,-9.1);
\node[anchor=north west,font=\footnotesize\itshape] at (-0.5,-3.45) {mỗi lượt quyết định};
\draw[tin] (0,-5.1)--(9.5,-5.1) node[nhan]{lô sự kiện: yêu cầu mới, xác nhận đặt xe,\\tiến độ xe, đón, trả,
thời gian đi mới};
\node[ben,text width=3.9cm,font=\footnotesize] at (9.5,-6.15) {sinh phương án, kiểm luật, chọn, ghi sổ
(Bảng~\ref{bang:luotquyetdinh})};
\draw[tin] (9.5,-7.7)--(0,-7.7) node[nhan]{quyết định: kế hoạch mới của các xe, lời hứa};
\draw[tin] (0,-8.7)--(9.5,-8.7) node[nhan]{đã áp dụng quyết định};
\draw[tin] (0,-10.15)--(9.5,-10.15) node[nhan]{yêu cầu trạng thái};
\draw[tin] (9.5,-11.15)--(0,-11.15) node[nhan]{trạng thái, gồm sổ theo từng khách};
\draw[tin] (0,-12.15)--(9.5,-12.15) node[nhan]{kết thúc: bộ chạy dừng, không trả lời};
\end{tikzpicture}
\caption[Các bản tin giữa bộ chuyển đổi FleetPy và bộ chạy RideBound]{Các bản tin chính giữa bộ chuyển đổi FleetPy và bộ chạy RideBound trong một lần chạy (đã giản lược). Mỗi
bản tin là một dòng JSON; phần trong khung nét đứt lặp lại ở mỗi lượt quyết định.}
\label{hinh:bantin}
\end{figure}
% EV: X3-15 (ProtocolPrimitives.cs:104-121; RunnerSession.cs:177-214: chỉ xử lý hello, initializeRun, eventBatch, decisionApplied, checkpoint, restore; shutdown làm bộ chạy dừng, không trả lời; finalizeRun/runSummary chỉ được khai báo; runner_client.py:158-273); X3-13 (SPEC:35-37)

\subsection{Giới hạn tính toán và thời gian chạy}
\label{muc:gioihan}

Ngoài giới hạn 2 yêu cầu mỗi lượt và 100 phương án mỗi xe, các bước sinh, kiểm và giải đều có giới hạn số bước
tính, và bộ giải dùng hạt giống ngẫu nhiên (số khởi tạo phần ngẫu nhiên) cố định bằng 7 để chạy lại cho cùng kết
quả. Vì vậy bộ sinh không xét hết mọi phương án; mục~\ref{muc:c4-hople} báo các bộ đếm của bộ sinh phương án cho biết
giới hạn có cắt bớt phương án hay không. Phép kiểm ở mỗi phương án duyệt mười chiều cho mỗi khách có lời hứa của xe và đếm đảo thứ tự
trên từng cặp điểm dừng, nên không thể coi khối lượng kiểm là tuyến tính theo độ dài kế hoạch; nó chưa được đo
riêng.
% EV: A24 (WP4:9-23); X3-17; X3-32 (GEN:54-79); N-25; X3-05 (BUDGET:76-99); X3-04 (VALID:233-239); X3-26 (DELTA:190-215)

Khóa luận chỉ báo thời gian thực của cả lần chạy: trên tập tắc đường cố định (16 lần chạy mỗi cấu hình), trung vị là
730, 667 và 641 giây với ghép tự do, điều phối không trần và luật C trần 30 giây. Thời gian này gồm mô phỏng FleetPy
và thời gian chờ RideBound trả quyết định, không phải thời gian của một lượt quyết định. Chi phí riêng của phép kiểm trần
không được báo, nên các số này không dùng để so chi phí tính toán.
% EV: H13 (EL17:2-6); X3-19; B17; R26-CH3-OPEN-QUESTIONS-CODEX.md §Q6; wp14r_scientific_protocol.py:801-803, 824-829; wp14r_supervised_process.py:297, 329-330, 979-985

\subsection{Kiểm tra kết quả}
\label{muc:kiemtra}

Kết quả được kiểm ở bốn mức.
\begin{itemize}[leftmargin=1.2cm,itemsep=2pt,topsep=3pt]
\item \textbf{Kiểm thử phần mềm:} chạy lại trên đúng mã nguồn đã dựng bộ chạy, cả 1.025 trong 1.025 kiểm thử tự động
đạt trong một lần chạy đầy đủ ngày 13/10/2026, với kiểu xuống dòng LF và một tệp chưa ghi nhận trong cây làm việc;
chi tiết ở Phụ lục~\ref{phuluc:C}.
\item \textbf{Kiểm tính nhất quán của sổ:} ở cả 382.006 bản ghi sửa của 185 lần chạy (170 lần chạy của năm cấu hình
trên tập kiểm tra chính và 15 lần chạy luật C trần 30 giây trên 15 kịch bản hiệu chỉnh), $|x|$, $|c|$, $|z|$ tính
lại từ $v$, $e$, $p$ trong sổ khớp đúng các trường độ lớn. Phép kiểm này không kiểm $e$ có được tính đúng hay không.
\item \textbf{Tính lại độc lập:} một chương trình viết riêng, độc lập về mã nhưng dùng cùng đặc tả chỉ số, tính lại
13 hạng mục số liệu từ bản ghi thô và một điều kiện luôn phải đúng, không lệch hạng mục nào; các số đếm chính được
ít nhất hai bộ đọc viết độc lập tính lại và khớp trong phạm vi đã kiểm. Hai mức kiểm này được làm bằng các chương
trình đọc riêng, không phải kiểm thủ công; chúng cho biết con số khớp bản ghi, không cho biết mô phỏng giống thực tế
đến đâu.
\item \textbf{Tái lập:} tập kiểm tra chính, tập chỉnh trần buổi sáng, các lần chạy trần 120 giây và tám lần chạy chiều
tuyết dùng cùng một bản bộ chạy đã biên dịch và lưu lại; tập tắc đường cố định dùng một bản cũ hơn. Bản dựng lại, nếu có,
phải đối chiếu với kết quả đã lưu trước khi dùng. Các lượt ưu tiên chi phí dùng bản riêng; bốn cặp kiểm
hành vi không thay thế việc kiểm đồng nhất toàn bộ hai bản chạy (Phụ lục~\ref{phuluc:D}).
\end{itemize}
% EV (kiểm thử): TESTRUN-45f8026-2026-10-12.md (mã nguồn bộ chạy; lượt 3 và 3b; điều kiện LF và tệp chưa ghi nhận); A28
% EV (sổ): A26 (T38:106; T38I:68); B10 (170 + 15); RUN38:1-8 (15 lần chạy C-30 trên 15 kịch bản hiệu chỉnh)
% EV (tính lại): C19 (T38I:79-97: 13 hạng mục + 1 bất biến); ADDENDUM-1:15-19; C04; E01; K41; A26 lưu ý; I01
% EV (tái lập): A28 (RUNW07:40-41); X3-10 (run_w07.py:8; run_w07x.py:7; run_t38c120.py:3-8); B17; NEW:13; I13; R26-CH3-OPEN-QUESTIONS-CODEX.md §Q8 (bản bộ chạy chiều tuyết)

Chương~\ref{chuong:thucnghiem} dùng các thành phần trên để đo luật C được gì và mất gì trên mô phỏng FleetPy, so
với điều phối không trần, luật hạn chót và cách đệm lời hứa, tức cộng thêm một khoảng cố định vào giờ hứa cho mọi
khách.
 của main.tex.
% Bản 1 (07-chuong3.tex) và bản trung gian 2 (07-chuong3-v2.tex) được giữ nguyên. Lý do từng thay đổi:
% reviews/CH3-REVISION-LOG.md. Ghi chú v3: 07-chuong3-v3-NOTES.md; thay đổi v4: 07-chuong3-v4-CHANGES-v3.md.
% Mỗi dòng "% EV:" trỏ tới EVIDENCE-REGISTER.md (A01…K52), EVIDENCE-REGISTER-ADDENDUM-1.md (N-01…N-30,
% L1a-01…09), EVIDENCE-REGISTER-ADDENDUM-2-CH3.md (X3-01…X3-41) hoặc các nguồn R26/TESTRUN ghi tại chỗ.
% =====================================================================================================

% Môi trường "Nhận xét" chưa có trong main.tex; định nghĩa có điều kiện ở đây (nên chuyển lên phần mở đầu).
\ifcsname nhanxet\endcsname\else
  \theoremstyle{definition}
  \newtheorem{nhanxet}{Nhận xét}[chapter]
\fi

\chapter{Phương pháp: luật C, lời hứa theo dự báo và sổ theo từng khách}
\label{chuong:phuongphap}

Khi đặt xe đi chung, hành khách được báo giờ xe sẽ trả mình tới nơi; khóa luận gọi đó là lời hứa giờ trả, gọi
tắt là lời hứa. Giờ trả có thể bị dời vì bộ điều phối ghé đón thêm khách, hoặc vì đường đông hơn lúc hứa hay xe
đi nhanh, chậm hơn dự kiến. Luật C, nội dung chính của chương, chỉ giới hạn loại thứ nhất. Ở mỗi lượt quyết định,
hiệu giữa giờ trả khi
chọn phương án mới và khi xe giữ nguyên kế hoạch là phần dời do điều phối; luật C loại phương án làm tổng độ lớn
phần này qua cả chuyến vượt một trần. Nhịp tắc đường quen thuộc trong ngày được cộng trước vào lời hứa đầu nhờ dự
báo; phần còn lại được ghi trong sổ theo từng khách (per-rider ledger). Các mục đi theo thứ tự đó, sau phần phát
biểu bài toán.
% EV: A01; A02; A06; F01; X3-12 (FFA:134-138, ba phần của sổ)

% -----------------------------------------------------------------------------------------------------
\section{Phát biểu bài toán và tổng quan hệ thống}
\label{muc:baitoan}

\subsection{Thuật ngữ và vòng đời của một yêu cầu}
\label{muc:thuatngu}

Bảng~\ref{bang:thuatngu} giải thích các thuật ngữ; tên tiếng Anh chỉ ghi trong bảng. Mọi thực nghiệm chạy trên
bộ mô phỏng FleetPy~\cite{engelhardt2026} với dữ liệu chuẩn Manhattan (mục~\ref{muc:kientruc}).
% EV: B01; J18 (FLEET:1-12, 474)

\begin{table}[htbp]
\centering
\caption[Thuật ngữ chính của phương pháp]{Thuật ngữ dùng trong chương}
\label{bang:thuatngu}
\small
\begin{tabular}{@{}p{4.3cm}p{10.5cm}@{}}
\toprule
\textbf{Thuật ngữ} & \textbf{Nghĩa trong khóa luận}\\
\midrule
Lượt quyết định (decision epoch) & Một lần bộ điều phối nhận một lô sự kiện và trả quyết định cho cả đội xe.\\
Phương án (candidate plan) & Một kế hoạch (dãy điểm đón, điểm trả theo thứ tự) có thể chọn cho một xe.\\
Phương án giữ nguyên kế hoạch (keep-route option) & Xe đi tiếp kế hoạch cũ, không nhận thêm khách.\\
Bộ sinh chỉ chèn thêm (insertion-only generator) & Chèn điểm đón, điểm trả của yêu cầu mới vào kế hoạch của
một xe, không bỏ, không đảo điểm cũ.\\
Bộ chọn (selector); bộ giải (solver) & Chọn mỗi xe một phương án theo thứ tự tiêu chí (lexicographic order):
chỉ khi hòa ở tiêu chí trước mới xét tiêu chí sau; bộ giải của OR-Tools~\cite{ortools} giải bài toán chọn.\\
Lời hứa; lời hứa đầu $p_0$ & Giờ đón, giờ trả báo cho khách; lời hứa đầu công bố khi khách xác nhận đặt xe;
$p_0$ là giờ trả của nó, chưa cộng dự báo.\\
Công bố; bản ghi sửa & Ghi lời hứa mới vào sổ và gửi trong bản tin quyết định; mỗi lần sau lời hứa đầu là một
bản ghi sửa.\\
Khóa (lock); chiều cam kết & Khóa: phần lời hứa không được đổi sau khi xác nhận. Chiều cam kết: thuộc tính của
lời hứa được theo dõi riêng, như giờ trả, việc đổi xe.\\
Trần $B$ (cap); phần đã dùng $L_r$ & Tổng phần dời do điều phối tối đa của một khách cả chuyến; phần khách $r$
đã dùng.\\
Quyết định bị ép (forced decision); luật phục hồi (recovery rule) & Lượt quyết định có ít nhất một xe bị ép:
luật lời hứa đã loại cả phương án giữ nguyên kế hoạch của xe, và xe chọn phương án được giữ nhờ luật phục hồi,
có ít nhất một khách được miễn trừ; luật phục hồi quy định phương án nào được giữ khi đó.\\
Hệ số tắc đường & Số nhân chung cho cả thành phố áp lên mọi thời gian đi, đổi theo khung 15 phút; hệ số 1,2 làm
thời gian đi dài gấp 1,2 lần hệ số 1.\\
Kịch bản & Một ngày, một mẫu nhu cầu, một khung 2 giờ của dữ liệu Manhattan: 108 yêu cầu, 8 xe 4 chỗ.\\
\bottomrule
\end{tabular}
\end{table}
% EV (lượt, phương án, bộ sinh, bộ chọn): X3-01; SPEC:31-34; X3-03; X3-07; X3-22; A14
% EV (lời hứa, công bố, khóa, chiều, trần): A13; X3-25 (fleet_control.py:882-893); A07; A04; A06; K43
% EV (bị ép, phục hồi, hệ số, kịch bản): X3-30; X3-09; B06; STRETCH:3-4; B02; B03; GLOSSARY-VI-v2.md §E (bị ép); C4A-10

Một yêu cầu mới chờ được xét. Khi phương án được chọn của một xe chứa nó, yêu cầu được nhận và khách được mời,
nhưng chưa có lời hứa. Khi khách xác nhận đặt xe, lời hứa đầu được công bố và $L_r = 0$; khách chờ đón, lên xe
rồi được trả. Yêu cầu chưa được nhận được xét lại ở lượt sau nếu còn phương án hợp lệ hay chưa được xét vì giới
hạn số yêu cầu mới, và bị từ chối khi mọi phương án chứa nó đều bị loại. Khóa và luật lời hứa áp từ lúc xác nhận
tới lúc khách xuống xe.
% EV: X3-20 (VALID:272-336; RollingCostPolicy.cs:173-226; fleet_control.py:878-924); A13; X3-32

\subsection{Bài toán ở mỗi lượt quyết định}
\label{muc:baitoan-phatbieu}

Bộ chạy tiến theo đồng hồ sự kiện: yêu cầu đến, xe tới mốc vật lý và mốc cập nhật tắc đường đều có thể gọi một
lượt quyết định. Xác nhận đặt xe và cập nhật trạng thái có thể tạo thêm lượt ở cùng thời điểm mô phỏng; không có
chu kỳ điều phối cố định. Bộ chuyển đổi gom các sự kiện từ lần gọi trước (yêu cầu mới, xác nhận đặt xe, tiến độ xe,
đón, trả, thời gian đi mới) thành một lô.
% EV: R26-CH3-OPEN-QUESTIONS-CODEX.md §Q1; actual_fleetpy_medium_preflight.py:528-543, 569-637; fleet_control.py:395-451; preflight_world_v1_1.py:40-58; preflight_world_v1_2.py:18-20

Khung nhu cầu dài 2 giờ, nhưng lần chạy kết thúc khi hết yêu cầu và hành trình đang hoạt động. Trong 34 lần chạy
luật C trần 120 giây của tập kiểm tra chính, thời lượng mô phỏng từ 118,02 đến 154,93 phút.
% EV: R26-CH3-OPEN-QUESTIONS-CODEX.md §Q2; round26-trace-rider.txt, Q2 C-120 MAIN-PANEL SIMULATED LENGTHS; actual_fleetpy_medium_preflight.py:580-581, 639-683

Bài toán ở mỗi lượt như sau.
\begin{itemize}[leftmargin=1.2cm,itemsep=2pt,topsep=3pt]
\item \textbf{Đầu vào:} trạng thái xe (vị trí, kế hoạch, sức chứa, khách trên xe, khách chờ đón); sổ; yêu cầu
đang chờ với điểm đón, điểm trả, cửa sổ đón (từ giờ đón sớm nhất tới giờ đón muộn nhất), thời gian ngồi xe tối
đa; thời gian đi lúc này, coi là giữ nguyên cho phần đường còn lại.
\item \textbf{Quyết định:} nhận, để lại hay từ chối từng yêu cầu đang chờ; mỗi xe một phương án thỏa các ràng
buộc cứng (Bảng~\ref{bang:rangbuoc}).
\item \textbf{Đầu ra:} lời hứa đầu và các bản ghi sửa (mục~\ref{muc:so}).
\item \textbf{Mục tiêu:} thứ tự tiêu chí của Bảng~\ref{bang:rangbuoc}; các chỉ số đánh giá ở
Chương~\ref{chuong:thucnghiem}.
\end{itemize}
% EV: X3-21 (fleet_control.py:395-429); X3-01; X3-02; A05; X3-20; B03; A06; A07; A14

Bộ chọn giải một quy hoạch nguyên, với $y_{k\pi} \in \{0, 1\}$ bằng 1 khi xe $k$ chọn phương án $\pi$ trong
tập $P_k$ của nó:
\begin{equation}
\sum_{\pi \in P_k} y_{k\pi} = 1 \;\;\text{với mọi xe } k, \qquad
\sum_{k} \sum_{\pi \in P_k,\; u \in \pi} y_{k\pi} \le 1 \;\;\text{với mọi yêu cầu mới } u ,
\label{eq:chon}
\end{equation}
tức mỗi yêu cầu mới thuộc phương án được chọn của nhiều nhất một xe. Các tiêu chí được giải lần lượt, mỗi lần giữ
các tiêu chí trước ở giá trị tốt nhất đã tìm. Bước chọn này theo lối gán xe--chuyến động của Alonso-Mora và cộng
sự~\cite{alonsomora2017} (mục~\ref{muc:c2-baitoan}).
% EV: X3-22 (OrTools:3, 96-125, 180-189, 343-365); X3-34 (HZ:341-345); J17

\begin{table}[htbp]
\centering
\caption[Ràng buộc cứng và thứ tự tiêu chí của bộ chọn]{Ràng buộc cứng của một phương án (nửa trên) và thứ tự tiêu chí của bộ chọn (nửa dưới): so tiêu chí 1
trước, chỉ khi hòa mới xét tiêu chí 2, và cứ thế. Giá trị là cấu hình của các tập thử, không phải chuẩn ngành.}
\label{bang:rangbuoc}
\small
\begin{tabular}{@{}p{2.4cm}p{7.6cm}p{4.8cm}@{}}
\toprule
\textbf{Nhóm} & \textbf{Nội dung} & \textbf{Vai trò}\\
\midrule
Ràng buộc dịch vụ & Đón trong vòng 10 phút kể từ lúc gọi; ngồi xe không quá 1,5 lần thời gian đi thẳng; mỗi
xe tối đa 4 khách & Phương án đổi kế hoạch phải thỏa chặt; phương án giữ nguyên được nới phần vượt thời gian
mà kế hoạch cũ đã có\\
% EV: B03; X3-23 (GEN:161-246; ServiceQualityAllowance.cs:61-94)
Khóa & Không đổi xe, điểm đón, giờ đón; điểm trả không đổi nhờ trần 0. Giờ đón mới phải bằng giờ đón khi giữ
nguyên kế hoạch & So với kế hoạch giữ nguyên, nên riêng tắc đường không làm vi phạm khóa\\
% EV: A07; A08; A10; X3-06 (LOCK:42-60, 103-161, 204-207; CFG C-30.json:66-67, 94-98)
Luật lời hứa & Luật C: $L_r + |p_r - e_r| \le B$ cho mọi khách có lời hứa của xe, hoặc một luật ở
Bảng~\ref{bang:luat} & Lọc trước bộ giải\\
% EV: A06; A09; A11; HVA:62-66
\midrule
Tiêu chí 1 & Ít phương án vi phạm được giữ lại nhất (chỉ khi có xe bị ép) & Chỉ dùng chúng khi không còn cách
khác\\
Tiêu chí 2 & Nhận nhiều yêu cầu mới nhất & Ưu tiên phục vụ\\
Tiêu chí 3 & Mức dùng trần lớn nhất trên đội xe càng nhỏ càng tốt; nếu hòa, tổng phần dời do điều phối nhỏ nhất,
so lần lượt từng chiều cam kết (khi có chiều mang trần) & Ưu tiên ổn định\\
Tiêu chí 4 & Tổng chi phí kế hoạch nhỏ nhất, tức thời lượng còn lại của các tuyến dự kiến & Theo mô hình, không
phải chi phí vận hành thực\\
Tiêu chí 5 & Thứ tự mã phương án & Phá hòa, để chạy lại cho cùng kết quả\\
% EV: A14; X3-07 (SEL:271-330, 365-385, 510-522; HVA:537-636); DIM:19-32; A16; I09
\bottomrule
\end{tabular}
\end{table}

\subsection{Một lượt quyết định}
\label{muc:luotquyetdinh}

Hình~\ref{hinh:vongquyetdinh} và Bảng~\ref{bang:luotquyetdinh} mô tả một lượt. Ngoài sinh phương án và chọn,
mỗi lượt có ba bước cho lời hứa: tính các mốc giờ trả của từng khách, loại phương án vi phạm khóa hay luật lời
hứa trước khi gọi bộ giải, và ghi sổ sau khi chọn. Hình và giả mã dùng một số ký hiệu được định nghĩa ở các mục
sau: với mỗi khách có lời hứa, $v$ là giờ trả đã công bố gần nhất, $e$ là giờ trả nếu xe giữ nguyên kế hoạch, $p$ là
giờ trả theo phương án đang xét, còn $x$, $c$, $z$ là ba hiệu giữa chúng (mục~\ref{muc:bamoc}); $D$ là lời hứa theo
dự báo, tính qua giờ $S$ (mục~\ref{muc:dubao}).
% EV: X3-04 (VALID:201-239, 355-411, 494-520); X3-13; X3-25

\begin{figure}[htbp]
\centering
\begin{tikzpicture}[
  node distance=6mm,
  hop/.style={draw=xamvien,fill=xamnhat,rounded corners=3pt,align=center,
              inner sep=4pt,font=\small,minimum height=12mm,text width=2.75cm},
  ngoai/.style={draw=xamvien,dashed,rounded corners=3pt,align=center,
              inner sep=4pt,font=\small,minimum height=10mm,text width=4.2cm},
  mui/.style={-{Latex[length=2.2mm]},thick}
]
\node[hop] (a) {Nhận lô sự kiện\\mới};
\node[hop,right=of a] (b) {Sinh phương án chỉ chèn thêm, kiểm ràng buộc dịch vụ};
\node[hop,right=of b] (c) {Tính $v$, $e$, $p$ cho từng khách có lời hứa};
\node[hop,right=of c] (d) {Loại phương án vi phạm khóa hay luật lời hứa};
\node[hop,below=9mm of d] (e) {Bộ giải chọn\\một phương án\\mỗi xe};
\node[hop,left=of e] (f) {Kiểm lại cả đội xe; công bố lời hứa $p$; ghi sổ};
\node[ngoai,below=8mm of f] (g) {Sau khi chạy, trên bản ghi đã lưu: lời hứa hiển thị $D = \max(p, S)$};
\draw[mui] (a)--(b); \draw[mui] (b)--(c); \draw[mui] (c)--(d);
\draw[mui] (d)--(e); \draw[mui] (e)--(f);
\draw[mui,dashed] (f.west) -| node[pos=0.3,below,font=\footnotesize]{lượt sau} (a.south);
\draw[-{Latex[length=2.2mm]},dashed] (f.south)--(g.north);
\end{tikzpicture}
\caption[Một lượt quyết định của RideBound]{Một lượt quyết định của RideBound. Bộ điều phối công bố lời hứa không dự báo $p$. Trong các kết quả
chính, lời hứa theo dự báo $D$ được tính sau khi chạy, từ bản ghi đã lưu (ô nét đứt), và không quay lại làm đầu
vào của quyết định nào (mục~\ref{muc:dubao}).}
\label{hinh:vongquyetdinh}
\end{figure}
% EV: X3-03; X3-04; X3-22; X3-25; F03; X3-11 (DISP:1-8); FR:10-11

\begin{table}[htbp]
\centering
\caption[Giả mã một lượt quyết định]{Giả mã một lượt quyết định; dấu $\gets$ nghĩa là gán giá trị bên phải cho biến bên trái}
\label{bang:luotquyetdinh}
\small
\begin{tabular}{@{}r@{\hspace{6pt}}p{10.2cm}@{\hspace{8pt}}p{4.0cm}@{}}
\toprule
\multicolumn{3}{@{}p{14.9cm}@{}}{\textbf{Đầu vào:} trạng thái đội xe; sổ; lô sự kiện mới; luật lời hứa đang
dùng. \textbf{Đầu ra:} kế hoạch mỗi xe; yêu cầu được nhận, để lại hay từ chối; lời hứa mới; bản ghi sổ}\\
\midrule
1: & \textbf{hàm LượtQuyếtĐịnh}(trạng thái, lô sự kiện) & \\
2: & \quad áp lô sự kiện; ghi nhận lần đón muộn hơn giờ đón muộn nhất, nếu có & vị trí xe, đón, trả\\
3: & \quad chọn tối đa 2 yêu cầu đang chờ có giờ đón muộn nhất sớm nhất & yêu cầu khác để lượt sau\\
4: & \quad \textbf{với mỗi} xe $k$ \textbf{thực hiện} & \\
5: & \qquad \textbf{nếu} kế hoạch đang chạy vi phạm ràng buộc cấu trúc \textbf{thì} báo lỗi, dừng lần chạy &
ví dụ vượt sức chứa\\
6: & \qquad $P_k \gets \{$phương án giữ nguyên$\}$, kiểm dịch vụ có nới phần vượt do tắc đường & luôn giữ khi
cắt bớt\\
7: & \qquad thêm vào $P_k$ các phương án chèn yêu cầu ở dòng 3, thỏa chặt ràng buộc dịch vụ; giữ tối đa 100 &
bộ sinh chỉ chèn thêm\\
8: & \qquad \textbf{với mỗi} $\pi \in P_k$ và mỗi khách có lời hứa của xe: tính $v$, $e$, $p$; \textbf{nếu}
$\pi$ vi phạm khóa hay luật lời hứa \textbf{thì} loại $\pi$ & Bảng~\ref{bang:kiemtran}\\
9: & \qquad \textbf{nếu} phương án giữ nguyên bị loại chỉ vì lỗi hạn chót hay lỗi trần \textbf{thì} & xét phục hồi\\
10: & \qquad\quad giữ nó làm phương án được miễn trừ & phục hồi giữ nguyên kế hoạch\\
11: & \qquad\quad giữ thêm phương án đã bị loại không làm vi phạm đó tệ hơn & phục hồi không tệ hơn\\
12: & \qquad \textbf{nếu} phương án giữ nguyên bị loại mà không được miễn \textbf{thì} báo lỗi, dừng lần chạy &
xe kẹt\\
13: & \quad chọn $y$ theo~\eqref{eq:chon} và thứ tự tiêu chí; kiểm lại lựa chọn trên cả đội xe &
Bảng~\ref{bang:rangbuoc}\\
14: & \quad nhận yêu cầu có trong phương án được chọn; để lại hay từ chối yêu cầu khác & \\
15: & \quad khách vừa xác nhận: công bố lời hứa đầu, $L_r \gets 0$; khách có lời hứa mà $x$, $c$ hay $z$ khác 0:
ghi bản ghi sửa, $L_r \gets L_r + |c_r|$ & ghi vi phạm cho khách được miễn\\
16: & \quad \textbf{trả về} quyết định; chờ bộ mô phỏng xác nhận đã áp dụng & \\
17: & \textbf{kết thúc hàm} & \\
\bottomrule
\end{tabular}
\end{table}
% EV (dòng 2–3): X3-01; X3-24 (WP4:8; EventReducer.cs:13-23, 260-282); X3-32 (GEN:54-96)
% EV (dòng 4–7): X3-31 (GEN:171-182); X3-23; X3-27 (CandidatePortfolioRetainer.cs:102-113); X3-03 (GEN:548-575); A24
% EV (dòng 8–12): X3-04; A06; A07; A10; X3-09 (HVA:244-390); X3-31 (POL:336-339; RunnerSession.cs:941-947)
% EV (dòng 13–16): X3-22; X3-07; POL:370-401; X3-20; X3-25 (VALID:272-336, 562-597; RunnerSession.cs:956-987); X3-30; RunnerSession.cs:448-455

Luật lời hứa là bộ lọc cứng: phương án vi phạm bị loại trước khi bộ giải chạy, không được giữ kèm một điểm trừ.
Ngoại lệ là khi xe bị ép (dòng 9--11), chẳng hạn khi riêng việc đường đông hơn đã đẩy giờ trả vượt hạn chót của
luật hạn chót (mục~\ref{muc:doichung}); khi đó tiêu chí 1 chỉ dùng phương án vi phạm nếu xe không còn cách khác.
Bộ sinh không bỏ, không đảo điểm cũ nào, nên khách đã nhận không bị chuyển sang xe khác. Phần nới ở dòng 6 giúp
xe không kẹt vì tắc đường. Mỗi lần vượt như vậy, và mỗi lần khách thực tế lên xe muộn hơn giờ đón muộn nhất, được ghi
lại thay vì làm dừng lần chạy. Vì vậy cửa sổ đón là điều kiện để kiểm phương án, không phải bảo đảm đón đúng hạn.
% EV: A06; HVA:62-66; X3-09 (HVA:255-326); X3-07 (SEL:283-295); A10; X3-03 (GEN:548-575); TEX08:129-131; X3-23 (GEN:161-246; ServiceQualityAllowance.cs:5-18, 61-94; RunnerSession.cs:989-1009); X3-24 (WP4:8; EventReducer.cs:13-23, 260-282)

% -----------------------------------------------------------------------------------------------------
\section{Ba mốc giờ trả và phần dời do điều phối}
\label{muc:bamoc}

Giờ trả có thể đổi vì bộ điều phối ghép thêm khách, hoặc ngay cả khi xe giữ kế hoạch cũ. Để tách hai loại, với
mỗi khách có lời hứa và mỗi phương án, hệ thống so ba mốc giờ trả:
\begin{itemize}[leftmargin=1.2cm,itemsep=2pt,topsep=3pt]
\item $v$: giờ trả đã công bố cho khách ở lần gần nhất;
\item $e$: giờ trả tính lại nếu xe giữ nguyên kế hoạch cũ, từ trạng thái hiện tại với thời gian đi hiện tại;
\item $p$: giờ trả nếu chọn phương án đang xét, với cùng trạng thái và thời gian đi.
\end{itemize}
Ba hiệu sau tính bằng giây (sổ lưu tới mili giây):
\begin{equation}
x = e - v, \qquad c = p - e, \qquad z = p - v = x + c .
\label{eq:bamoc}
\end{equation}
Trong đó, $x$ là phần thay đổi khi giữ nguyên kế hoạch, $c$ là phần dời do điều phối, $z$ là phần khách thấy;
dễ nhớ, $e$ là giờ nếu xe cứ đi như cũ. Mốc $e$ không phải ý của khóa luận: Hörl và Zwick~\cite{horl2022} đã tính
lại giờ trả trên tuyến đang gán theo thời gian đi hiện tại ở mỗi lượt, và khóa luận dùng lại mốc này.
% EV: A01; A02; X3-04 (VALID:355-384); X3-13 (SPEC:39-50); LEDGER:188-207; J01

Mốc $e$ là giờ do mô hình tính, không phải giờ đến thật, nên $x$ gồm mọi thứ đổi khi giữ nguyên kế hoạch: thời
gian đi mới, tiến độ thực của xe, sự kiện vừa xảy ra. Phép tách $z = x + c$ chỉ là phép trừ giữa các giờ của mô
hình, không chứng minh quyết định điều phối thực sự gây ra đúng $c$ giây chậm trễ. Hình~\ref{hinh:bamoc} là một ví
dụ dựng tay: ghép khách M làm A tới muộn hơn lời hứa 70 giây, nhưng chỉ 30 giây là phần dời do điều phối.
% EV: A01 (không phải giờ đến thật, không nhận diện nhân quả); A02; X3-04 (VALID:355-361); K04; ví dụ dựng tay

\begin{figure}[htbp]
\centering
\begin{tikzpicture}[x=1cm,y=1cm,font=\small,
  diem/.style={draw=xamvien,fill=xamnhat,rounded corners=2pt,inner sep=3pt,minimum height=6mm},
  mui/.style={-{Latex[length=2mm]},thick}]
\node[anchor=east] at (2.6,4.0) {Giữ nguyên kế hoạch:};
\node[diem] (x1) at (3.4,4.0) {Xe};
\node[diem] (a1) at (5.6,4.0) {Trả A};
\draw[mui] (x1)--(a1);
\node[anchor=east] at (2.6,3.0) {Ghép thêm khách M:};
\node[diem] (x2) at (3.4,3.0) {Xe};
\node[diem] (b2) at (5.6,3.0) {Đón M};
\node[diem] (a2) at (7.8,3.0) {Trả A};
\node[diem] (c2) at (10.0,3.0) {Trả M};
\draw[mui] (x2)--(b2); \draw[mui] (b2)--(a2); \draw[mui] (a2)--(c2);
\draw[-{Latex[length=2mm]}] (0.8,0)--(12.2,0) node[right]{giờ};
\foreach \vitri/\tenmoc/\giomoc in {2.0/v/9:10:00, 6.8/e/9:10:40, 10.4/p/9:11:10}
  \draw (\vitri,0.12)--(\vitri,-0.12) node[below,align=center]{$\tenmoc$\\\giomoc};
\draw[mui,dashed] (2.0,0.55)--(6.8,0.55) node[midway,above]{$x = e - v = 40$ giây};
\draw[mui] (6.8,0.55)--(10.4,0.55) node[midway,above]{$c = p - e = 30$ giây};
\draw[mui,densely dotted] (2.0,1.45)--(10.4,1.45) node[midway,above]{$z = p - v = 70$ giây};
\end{tikzpicture}
\caption[Ba mốc giờ trả trong một ví dụ dựng tay]{Ví dụ dựng tay: khách A được hứa trả lúc 9:10:00 ($v$); lúc 9:00, xe đang chở A xét ghép thêm khách M.
Tính với thời gian đi hiện tại, giữ nguyên kế hoạch thì A tới lúc 9:10:40 ($e$), ghép M thì tới lúc 9:11:10
($p$). Nét đứt: $x$; nét liền: $c$; nét chấm: $z = x + c$. Luật C chỉ tính $|c|$.}
\label{hinh:bamoc}
\end{figure}

Phần dời do điều phối được tính theo độ lớn: sớm hay muộn hơn $e$ 20 giây đều tính 20 giây, vì luật đo lần sửa
lời hứa chứ không đo lợi hay hại; Fielbaum và Alonso-Mora~\cite{fielbaum2020} thì chỉ tính phần bất lợi. Trên mọi
lần chạy đã ghi, $c \ge 0$ tới mức mili giây; đây là quan sát, chưa phải chứng minh cho trường hợp tổng quát. Trên
tập kiểm tra chính, khóa luận đã kiểm 136 lần chạy của bốn cấu hình: điều phối không trần, luật C trần 30 và 120
giây, luật hạn chót 30 giây. Ở các lần chạy này, mỗi lần kéo giờ trả sớm hơn $e$ chỉ sớm 1 mili giây. Hệ quả: khi
bộ điều phối chỉ chèn thêm, phần tính hai chiều và không hoàn lại của luật C gần như không có việc; nó chỉ có tác
dụng khi kế hoạch của khách đã nhận được phép đổi, ví dụ khi dời khách đang chờ sang xe khác.
% EV: A04; X3-05 (DELTA:137-139, 259-262); J06 (FAM:293-296); N-05 (CORRECTED); ADD1:104-112; ADDENDUM-1 mục 0.1–0.2 (không nêu nguyên nhân của −1 ms)

% -----------------------------------------------------------------------------------------------------
\section{Luật C: trần phần dời giờ trả do điều phối cho từng khách}
\label{muc:luatc}

\begin{dinhnghia}[Luật C]
\label{dn:luatc}
Gọi $B$ là trần, tính bằng giây, và $L_r$ là tổng các $|c|$ đã tính cho khách $r$ kể từ lúc $r$ xác nhận đặt xe,
với $L_r = 0$ lúc xác nhận. Một phương án của một xe chỉ được giữ lại nếu với mọi khách $r$ có lời hứa đang được
gán cho xe đó (chờ đón hay đang trên xe),
\begin{equation}
L_r + |p_r - e_r| \le B .
\label{eq:luatc}
\end{equation}
Sau khi quyết định của lượt được chọn, với mỗi khách $r$ có lời hứa mà phương án được chọn làm $x$, $c$ hay $z$
khác 0, $L_r$ được cộng thêm $|p_r - e_r|$ của phương án đó, kể cả khi giờ trả công bố không đổi ($p_r = v_r$).
\end{dinhnghia}
% EV: A06 (BUDGET:41-105; LEDGER:285-288); A13; X3-05 (BUDGET:107-114; CFG C-30.json:17-23); X3-25 (VALID:272-281, 562-597)

Dấu $\le$ cho phép $L_r + |c|$ đúng bằng $B$. Phương án chỉ được xét không làm đổi $L_r$: mỗi phương án được
kiểm trên bản sao của sổ, chỉ quyết định cuối của lượt mới được ghi vào sổ thật. $B$ giống một ngân sách thời
gian của từng khách: $B - L_r$ là phần trần còn lại, không được cấp lại sau mỗi lần dời. Vì phần dời được cộng
qua cả chuyến, một luật cho mỗi lần dời tới 30 giây cho phép ba lần, tổng 90 giây, còn luật C với $B = 30$ giây thì
không. Ở Hình~\ref{hinh:bamoc}, với $L_A = 0$ và $B = 30$ giây, phương án ghép M ($|c| = 30$ giây) vẫn được giữ và
A dùng hết trần. Khi phép kiểm không được miễn cho $r$ ở lượt nào, tổng $|c|$ của $r$ cả chuyến không vượt $B$.
% EV: X3-05 (BUDGET:88); X3-25 (HVA:161-174; RunnerSession.cs:956-987; VALID:640-644: miễn trừ vẫn cộng vào sổ); A06; K43; E08; J05; CFG C-30.json:3-6 (một chính sách cho mọi khách)

\begin{table}[htbp]
\centering
\caption[Giả mã phép kiểm luật C cho một phương án]{Giả mã phép kiểm luật C cho một phương án (chỉ viết chiều giờ trả)}
\label{bang:kiemtran}
\small
\begin{tabular}{@{}r@{\hspace{6pt}}p{9.6cm}@{\hspace{8pt}}p{4.6cm}@{}}
\toprule
\multicolumn{3}{@{}p{14.9cm}@{}}{\textbf{Đầu vào:} phương án $\pi$ của xe $k$; trần $B$; $L_r$ của mỗi khách
có lời hứa; trạng thái và thời gian đi của lượt. \textbf{Đầu ra:} giữ hay loại $\pi$, kèm bản ghi chặn}\\
\midrule
1: & \textbf{hàm KiểmTrần}($\pi$, $k$, $B$) & \\
2: & \quad \textbf{với mỗi} khách $r$ có lời hứa được gán cho xe $k$, theo thứ tự mã khách \textbf{thực hiện} &
chờ đón hay trên xe\\
3: & \qquad $e_r \gets$ giờ trả của $r$ khi chiếu lại kế hoạch cũ của $k$ từ trạng thái hiện tại & mốc giữ
nguyên kế hoạch\\
4: & \qquad $p_r \gets$ giờ trả của $r$ theo $\pi$, cùng trạng thái và thời gian đi; $c_r \gets p_r - e_r$ &
phần dời do điều phối\\
5: & \qquad \textbf{nếu} $L_r + |c_r| > B$ \textbf{thì trả về} loại $\pi$, kèm ($r$, $L_r$, $|c_r|$,
$L_r + |c_r|$, $B$) & dừng ở khách vi phạm gặp trước\\
6: & \quad \textbf{trả về} giữ $\pi$ & $L_r$ chỉ đổi nếu $\pi$ được chọn\\
7: & \textbf{kết thúc hàm} & \\
\bottomrule
\end{tabular}
\end{table}
% EV (dòng 2–4): X3-04 (VALID:233-239, 355-361; HVA:171); X3-27 (VALID:1048-1064)
% EV (dòng 5–6): X3-05 (BUDGET:6-12, 72-104); A06; A23; X3-25

Cài đặt theo dõi mười chiều cam kết: thay đổi giờ đón, giờ trả; số lần sửa đáng kể; đổi xe; khoảng dời và số lần
đổi của điểm đón, của điểm trả; số lần đảo thứ tự với khách cũ; số điểm dừng mới chèn trước điểm đón. Chỉ bốn
chiều có trần: giờ trả ($B$) và đổi xe, đổi điểm đón, đổi điểm trả (trần 0). Ở tập kiểm tra chính và tập buổi sáng,
phép kiểm dừng ở khách vi phạm gặp trước, nên mỗi bản ghi chặn (bản ghi lý do loại phương án) chỉ nêu lỗi của khách
đó.
% EV: A04 (DIM:3-32); X3-26 (DELTA:137-153); A25; A07; CFG C-30.json:7-89; X3-07 (SEL:319-330, 510-522); A23; X3-04 (VALID:222-239); WP4:1-25; X3-10

\subsection{Phương án giữ nguyên kế hoạch dưới luật C}
\label{muc:giunguyen}

Gọi \textbf{(N)} là giả thiết, áp cho lượt đang xét: (i) kế hoạch đang chạy không vi phạm ràng buộc cấu trúc
(sức chứa, đón trước trả, phần kế hoạch đã cố định), nên bộ sinh dựng được phương án giữ nguyên; (ii) giờ trả
theo phương án giữ nguyên được tính bằng đúng phép chiếu dùng cho $e$. RideBound thỏa (ii) vì dùng cùng một hàm
chiếu với cùng đầu vào, và luôn giữ phương án giữ nguyên khi cắt còn 100 phương án.
% EV: A22; X3-31 (GEN:171-182, 213-229); ServiceQualityAllowance.cs:66-73; X3-27 (VALID:355-360, 1048-1064; HVA:144-165; CandidatePortfolioRetainer.cs:102-113; WP4:15)

\begin{nhanxet}
\label{nx:giunguyen}
Giả sử (N) và mọi khách có lời hứa của xe có $0 \le L_r \le B$. Khi đó luật C không loại phương án giữ nguyên
kế hoạch.
\end{nhanxet}

\noindent\textit{Chứng minh.} Theo (N)(ii), với mọi khách $r$ có lời hứa của xe, $p_r = e_r$, nên
$L_r + |p_r - e_r| = L_r \le B$ và phép kiểm~\eqref{eq:luatc} cho qua. \hfill$\square$
% EV: A22 (theo cấu tạo, giả thiết N); A06; X3-27

Khóa và các trần 0 cũng so với kế hoạch giữ nguyên, nên phương án này qua cả chúng. Nếu (N) đúng ở mọi lượt thì
$L_r \le B$ được giữ: $L_r = 0$ lúc xác nhận, và nếu $L_r \le B$ ở đầu một lượt thì không có miễn trừ nào, phương
án được chọn đã qua~\eqref{eq:luatc}, nên $L_r$ mới vẫn không vượt $B$. Khi đó xe dùng luật C hay điều phối không
trần không bao giờ bị ép; tính chất này vì vậy không riêng của luật C.
% EV: A10 (LOCK:156-161); X3-06; X3-26 (véc-tơ c của phương án giữ nguyên bằng 0 ở mọi chiều); X3-09 (HVA:255-259); X3-25

Nhận xét đúng do cách đặt luật, không phải kết quả thực nghiệm; số liệu chỉ phù hợp với nó: trên tập kiểm tra
chính, 0 trong 20.101 quyết định bị ép với trần 30 giây và 0 trong 20.435 với trần 120 giây (mỗi quyết định ứng với
một lượt của cả đội xe).
% EV: A21; N-21; X3-16 (STG:43-44); N-07; X3-30 (RunnerSession.cs:1013-1027); SPEC:31-34

Nhận xét có ba giới hạn. Thứ nhất, nếu (N)(i) không thỏa, lần chạy dừng vì lỗi cấu trúc, không vì luật C; phần
vượt cửa sổ đón hay thời gian ngồi xe do tắc đường không làm (N) hỏng, nhờ phần nới. Thứ hai, với $v$ = 9:10:00 và
$e = p$ = 9:11:00, phương án giữ nguyên có $c = 0$ nhưng $z = 60$ giây: luật C không bảo đảm giờ đến. Thứ ba, với
$v = p$ = 9:10:00 và $e$ = 9:12:00, kéo giờ trả về đúng giờ đã hứa có $z = 0$ nhưng $|c| = 120$ giây, và $L_r$
tăng 120 giây dù giờ công bố không đổi. Tính chất và giới hạn thứ hai là hai mặt của một lựa chọn: luật C không
bị ép vì không tính phần thay đổi khi giữ nguyên kế hoạch, và vì thế cũng không giới hạn phần đó.
% EV: A22; X3-31 (GEN:171-182; RunnerSession.cs:941-947); X3-23; K05; X3-25 (VALID:562-597: bản ghi sửa khi x ≠ 0 dù p = v); ví dụ dựng tay

% -----------------------------------------------------------------------------------------------------
\section{Các luật đối chứng, bộ chọn và luật phục hồi}
\label{muc:doichung}

Bảng~\ref{bang:luat} định nghĩa bốn cấu hình mà Chương~\ref{chuong:thucnghiem} so với luật C. Trừ ghép tự do,
chúng dùng cùng bộ sinh, bộ kiểm, bộ chọn và luật phục hồi với luật C; tập tắc đường cố định chạy bằng bộ chạy
chưa có luật phục hồi.
% EV: A08; A09; A11; A17; A19; K50; HO:109; B17; NEW:13

\begin{table}[htbp]
\centering
\caption[Các cấu hình được so sánh]{Các cấu hình được so sánh; ký hiệu chỉ dùng trong bảng. $B$ là trần, tính bằng giây; với luật hạn chót,
$B$ là số giây giờ trả được phép muộn hơn lời hứa đầu $p_0$. Các tập thử được mô tả ở mục~\ref{muc:c4-thietlap}.}
\label{bang:luat}
\footnotesize
\setlength{\tabcolsep}{3pt}
\begin{tabular}{@{}p{2.2cm}p{1.1cm}p{3.1cm}p{2.3cm}p{2.6cm}p{3.5cm}@{}}
\toprule
\textbf{Tên trong khóa luận} & \textbf{Ký hiệu} & \textbf{Kiểm điều gì} & \textbf{Khóa} & \textbf{Khi xe bị
ép} & \textbf{Dùng ở tập nào}\\
\midrule
Điều phối không trần & U & Không kiểm giờ trả & Xe, điểm đón, giờ đón; điểm trả có trần 0 & Không bị ép khi (N)
đúng & Tập kiểm tra chính; tập chỉnh trần buổi sáng; tập tắc đường cố định\\
% EV: A07; A08; K50; A10; X3-16 (STG:42: 0 / 20.760); B10; B15; B17; NEW:11
Ghép tự do & FREE & Không kiểm lời hứa; chọn theo chi phí kế hoạch ở từng lượt & Không khóa & Không áp dụng &
Tập tắc đường cố định\\
% EV: A17; K50; HO:109 (b1-rolling-cost, không khóa); B17; X3-38
Luật C & C-$B$ & $\sum |c| \le B$ cả chuyến, cho mọi khách có lời hứa & Như U & Không bị ép khi (N) đúng
(Nhận xét~\ref{nx:giunguyen}) & Tập chính (30, 120); buổi sáng (30, 60, 90, 120, 150, 180, 240, 300); tắc đường
cố định (30, 60, 120); chiều tuyết (30)\\
% EV: A06; A22; B10; L1a-01; X3-28 (run_w07.py:5; run_w07x.py:3-4, 35-36); B17; B20; N-26 (chiều tuyết 0 / 5.067 quyết định bị ép)
Luật hạn chót & M$^{+}$-$B$ & $p - p_0 \le B$ ở mọi lần kiểm; chỉ chặn phía muộn & Như U & Có thể bị ép do tắc
đường; phục hồi không tệ hơn & Tập chính (30, 60); buổi sáng (tám mức như luật C)\\
% EV: A09; A10; A19; B10; X3-28; D04; LOCK:186-197; CFG Mplus-30.json:98
Luật trần độ dời khách thấy & V-$B$ & $\sum |z| \le B$ cả chuyến & Như U & Có thể bị ép do tắc đường; phục hồi
không tệ hơn & Tập chính (30); buổi sáng (30, 90, 120, 150, 180, 240)\\
% EV: A11; A19; B10; X3-28 (run_w07x.py:3-4, 35-36: V-90/120/150/180/240; run_w07.py:5: V-30)
\bottomrule
\end{tabular}
\end{table}

Luật hạn chót chỉ mượn thành phần hạn chót tính từ lời hứa đầu của Schulz và Pfeiffer~\cite{schulz2026}, vốn
dùng thời gian đi cố định và độ chùng mặc định 900 giây. Ở đây mỗi độ chùng bằng một mức trần của luật C trên cùng
tập, và với tắc đường đổi theo khung 15 phút, riêng tắc đường có thể làm luật này bị ép; số quyết định bị ép ở
Chương~\ref{chuong:thucnghiem} nói về một hạn chót chặt như vậy, không về phương pháp của họ.
% EV: J10 (SP:466-490, 568-577); J11; A08; A09; X3-28; X3-41 (D03:214-217); A10

Ba luật lời hứa khác nhau ở đại lượng được cộng. Đánh số các lần cập nhật lời hứa của một khách là
$1, \ldots, n$. Trong sổ, mỗi bản ghi sửa so với lời hứa ngay trước nó ($v_j = p_{j-1}$), nên các giờ ở giữa
triệt tiêu:
\begin{equation}
\sum_{j=1}^{n} (x_j + c_j) \;=\; \sum_{j=1}^{n} z_j \;=\; p_n - p_0 .
\label{eq:dongnhat}
\end{equation}
Luật hạn chót kiểm $p_n - p_0 \le B$, tức cộng có dấu cả hai phần; luật trần độ dời khách thấy cộng $|z_j|$; luật
C chỉ cộng $|c_j|$. Ở Hình~\ref{hinh:bamoc}, nếu 9:10:00 là lời hứa đầu, phương án giữ nguyên có
$e - p_0 = |z| = 40$ giây, vượt trần 30 giây của hai luật kia dù $c = 0$: riêng tắc đường đủ làm chúng bị ép.
% EV: A12; X3-33 (LEDGER:256-269, 298-314; SPEC:62-63); A09; A10 (LOCK:156-161); A11; X3-16 (U 0 / 20.760)

Tiền lệ trực tiếp là Hörl và Zwick~\cite{horl2022} (mục~\ref{muc:c2-rangbuoc}): ở đầu mỗi lượt họ tính mốc $e$,
và công thức 7 đặt giờ trả bắt buộc bằng giá trị lớn hơn giữa giờ cập nhật và giờ trả muộn nhất, như ràng buộc cứng
(ràng buộc 9--10), kèm phạt cho phần vượt giờ gốc (công thức 11). Mục đích của họ, tránh tình huống không còn thứ tự
trả khách khả thi, cũng là mục đích của Nhận xét~\ref{nx:giunguyen}. Đặt giờ trả muộn nhất là $p_0 + B$ cho hạn chót
nới $p \le \max(p_0 + B, e)$, một cách viết lại của khóa luận; nó cho phép $c \le \max(p_0 + B - e, 0)$, về 0 khi
tắc đường đẩy $e$ quá $p_0 + B$, và cũng không loại phương án giữ nguyên. Cấu hình hạn chót M$^{+}$ với luật phục
hồi ``không tệ hơn'' ở chương~4 chính là dạng này. Luật C khác ở ba chỗ:
một trần cho mọi khách có lời hứa, còn phần nới của họ chỉ tác dụng khi giờ cập nhật đã vượt giờ trả muộn nhất;
phần được phép dời, $|c| \le B - L_r$, không phụ thuộc tắc đường; và cộng qua cả chuyến, dù phần cộng dồn hiếm khi là
lý do loại phương án (mục~\ref{muc:c4-congdon}).
% EV: J01; J02 (HZ:481-497); J03; X3-36 (HZ:531-543, 565-606); J05; A22 (suy luận V17 R6); K02; K26; K27; N-15

Từ Fielbaum và Alonso-Mora~\cite{fielbaum2020} (mục~\ref{muc:c2-dolech}), khóa luận kế thừa ý tách phần thay đổi
do quyết định gán xe và kiểm soát nó trước khi quyết định, nhưng mỗi lần sửa được so với mốc $e$ chứ không so điểm
đầu với điểm cuối, nên hai lần dời ngược chiều không triệt tiêu nhau; chiều sớm cũng được tính; đại lượng là giờ
trả chứ không phải thời gian chờ; và tắc đường thay đổi chứ không chọn giờ thấp điểm.
% EV: J06; J07; J08; X3-35 (FAM:90-94, 189-195, 216-220, 293-296, 948-972); K30; K32

Vì đổi xe, đổi điểm đón, đổi điểm trả là các trần 0, phần đầu của tiêu chí 3 (mức dùng trần lớn nhất, tỉ số lớn
nhất giữa phần đã dùng và trần) luôn bằng 100\% ở các lượt đã ghi của tập kiểm tra chính (mục~\ref{muc:c4-hople}),
cả khi không trần; tiêu chí 3 vì vậy chỉ tác dụng qua tổng phần dời theo từng chiều. Điều phối không trần và luật C
cùng cấu hình và bộ chọn, chỉ khác trần giờ trả; chênh lệch giữa chúng là phần đi kèm việc thêm trần khi giữ nguyên
bộ chọn. Khóa luận còn xét thứ tự ưu tiên chi phí: giữ ưu tiên nhận yêu cầu và các ràng buộc, bỏ các mức xếp
hạng theo ổn định rồi so chi phí kế hoạch và tiêu chí phân xử. Phương án vẫn phải qua phép kiểm trần khi trần
được bật. Mục~\ref{muc:c4-vaitro} so hai thứ tự ở trần 30 giây và không trần; đổi thứ tự không làm thay đổi
định nghĩa hay cách cộng trần của luật C.
% EV: S15-03 (costfirst-diff, cấu hình và provenance); S15-01 (mục 0, 2).
% EV: A14; X3-07 (SEL:271-330; HVA:537-636); A15; N-25; A08; H01; H02; K18

Luật phục hồi giữ nguyên kế hoạch giữ phương án giữ nguyên làm phương án được miễn trừ nếu mọi lỗi của nó là lỗi
hạn chót hay lỗi trần. Luật phục hồi không tệ hơn mở rộng luật này: giữ thêm phương án đổi kế hoạch có mỗi lỗi,
cùng khách và cùng chiều, cũng là lỗi của phương án giữ nguyên và không nặng hơn. Lỗi khóa theo giai đoạn (khóa của
khách chờ đón hay trên xe) và lỗi ràng buộc dịch vụ không bao giờ được miễn. Mỗi khách được miễn được ghi một vi
phạm, trừ khi lời hứa, giai đoạn và mã lỗi không đổi so với lần trước. Tập kiểm tra chính, tập buổi sáng và tám
lần chạy chiều tuyết dùng luật thứ hai; cả hai là chi tiết cài đặt, không phải đóng góp của khóa luận.
% EV: A17; A18 (ADR:5066-5092); X3-09 (HVA:79-99, 255-326; VALID:413-420, 475-483, 537-544, 599-638, 791-818; LOCK:42-60); A19 (RUN38:9; WP4:7); X3-10; R26-CH3-OPEN-QUESTIONS-CODEX.md §Q8 (tám lần chạy chiều tuyết; wp4-tier3-v2.json:7)

% -----------------------------------------------------------------------------------------------------
\section{Lời hứa theo dự báo}
\label{muc:dubao}

Vì bộ điều phối dùng thời gian đi hiện tại, lời hứa có thể sớm hơn giờ trả thực tế khi đường đông dần, dù xe giữ
nguyên kế hoạch. Lời hứa theo dự báo, một cách làm đã có (mục~\ref{muc:c2-dubao}), cộng sẵn mức tăng tắc đường các
ngày trước cho thấy vào từng giờ đón, giờ trả được công bố.
% EV: A05; F10; J35; F01; X3-11 (DISP:1-8)

Thời gian đo bằng giây kể từ đầu khung giờ của kịch bản; $i(\tau) = \lfloor \tau / 900 \rfloor$ là khung 15 phút
chứa $\tau$. Gọi $m_i$, $i = 0, \ldots, 8$, là trung bình hệ số tắc đường của khung $i$ trên các ngày nguồn; từ
khung 8 hệ số được giữ ở giá trị cuối, như bộ mô phỏng cũng giữ, nên đặt $\tilde m_i = m_{\min(i, 8)}$. Với lúc
công bố $t$ và giờ công bố $p > t$ (giờ của phương án được chọn, chưa cộng dự báo), $S$ là nghiệm của
\begin{equation}
\int_{t}^{S} \frac{\mathrm{d}\tau}{g(\tau)} = p - t, \qquad
g(\tau) =
\begin{cases}
1 & \text{nếu } i(\tau) = i(t),\\
\tilde m_{i(\tau)} / \tilde m_{i(t)} & \text{nếu không,}
\end{cases}
\label{eq:keodan}
\end{equation}
và lời hứa hiển thị là $D = \max(p, S)$ (giả mã ở Bảng~\ref{bang:dubao}).
% EV: F01; X3-11 (DISP:3-8, 55-64; STRETCH:1-12, 22-52; preflight_world_v1_3.py:51-82)

Theo tốc độ lúc công bố, xe cần thêm $p - t$ giây. Trong khung chậm hơn khung hiện tại $g$ lần, mỗi giây đồng hồ chỉ
đi được phần đường mà ở tốc độ lúc công bố mất $1/g$ giây. Mức tắc đường lúc này đã có trong thời gian đi đang dùng, nên $g = 1$ ở khung
hiện tại; các ngày khác chỉ cho biết khung sau chậm hay nhanh hơn bao nhiêu lần. Phép lấy $\max$ bảo đảm $D \ge p$ ở
cùng lần công bố, không bảo đảm $D$ tăng dần qua các lần cập nhật. Ví dụ dựng tay: trong kịch bản khung
8:00--10:00, lúc $t$ = 9:05:00 hệ thống công bố $p$ = 9:25:00, tức còn 1.200 giây. Nếu các ngày nguồn có hệ số trung
bình 1,25 ở khung 9:00--9:15 và 1,50 ở khung 9:15--9:30, thì sau 9:15 $g = 1{,}50/1{,}25 = 1{,}2$. Tới 9:15 xe đi
được 600 giây quãng đường; 600 giây còn lại mất $600 \times 1{,}2 = 720$ giây, nên $S = D$ = 9:27:00, dài thêm 120
giây.
% EV: X3-11 (STRETCH:1-12; forecast_weight neo ở khung hiện tại); F01; DISP:60-64; ví dụ dựng tay X3-37 (a), tính bằng hàm stretch của dự án; số học ở 07-chuong3-v3-NOTES.md mục 4

\begin{table}[htbp]
\centering
\caption[Giả mã lời hứa theo dự báo]{Giả mã lời hứa theo dự báo cho một giờ đón hay giờ trả đã công bố}
\label{bang:dubao}
\small
\begin{tabular}{@{}r@{\hspace{6pt}}p{9.6cm}@{\hspace{8pt}}p{4.6cm}@{}}
\toprule
\multicolumn{3}{@{}p{14.9cm}@{}}{\textbf{Đầu vào:} lúc công bố $t$; giờ công bố $p$; ngày $d$ và khung giờ của
kịch bản; hệ số tắc đường theo khung của các ngày trước $d$. \textbf{Đầu ra:} lời hứa hiển thị $D$}\\
\midrule
1: & \textbf{hàm LờiHứaTheoDựBáo}($t$, $p$, $d$) & \\
2: & \quad $H \gets$ các ngày cùng loại với $d$ và trước $d$, bỏ ngày 15/11 & chỉ dùng ngày đã qua\\
3: & \quad \textbf{nếu} $H$ có ít hơn 2 ngày \textbf{hoặc} $p \le t$ \textbf{thì trả về} $p$ & không dự báo\\
4: & \quad $m_i \gets$ trung bình hệ số của khung $i$ trên $H$, $i = 0, \ldots, 8$ & khung tính từ đầu khung giờ\\
5: & \quad $W \gets p - t$;\; $\tau \gets t$;\; $i_0 \gets i(t)$ & quãng còn lại theo tốc độ hiện tại\\
6: & \quad \textbf{lặp} & \\
7: & \qquad $g \gets 1$ nếu $i(\tau) = i_0$, ngược lại $g \gets \tilde m_{i(\tau)} / \tilde m_{i_0}$ &
$\tilde m_i = m_{\min(i, 8)}$\\
8: & \qquad $h \gets$ thời gian từ $\tau$ tới hết khung $i(\tau)$ & \\
9: & \qquad \textbf{nếu} $h/g \ge W$ \textbf{thì} $S \gets \tau + W g$ và thoát vòng lặp & \\
10: & \qquad $W \gets W - h/g$;\; $\tau \gets$ đầu khung kế tiếp & \\
11: & \quad \textbf{trả về} $D \gets \max(p, S)$ & $D \ge p$ ở lần công bố này\\
12: & \textbf{kết thúc hàm} & \\
\bottomrule
\end{tabular}
\end{table}
% EV (dòng 2–4): F02; X3-11 (DISP:27-57; preflight_world_v1_3.py:51, 80-82)
% EV (dòng 5–11): F01; X3-11 (STRETCH:25-39, 49-52)

Ngày nguồn là các ngày cùng loại (ngày thường hay cuối tuần) trước ngày của kịch bản, ít nhất 2 ngày, và không bao
giờ là ngày tuyết 15/11/2018 (một lựa chọn của khóa luận; hệ thống thật cần quy tắc tự động nhận ra ngày bất
thường). Vì vậy trong dữ liệu 12--18/11/2018 chỉ ba ngày 14, 15, 16/11 có dự báo: 21 trong 34 kịch bản của tập
kiểm tra chính; 13 kịch bản ngày 17, 18/11 giữ lời hứa không dự báo.
% EV: F02 (DISP:10-16, 27-52); F04; N-16; N-28

Trong các kết quả chính, $D$ được tính lại trên các lần chạy đã ghi, gọi là phát lại (replay): xe nào chở ai giữ
nguyên, chỉ giờ báo cho khách đổi, và $D$ không quay lại sổ hay phép kiểm trần. $D_0 - p_0$, với $D_0$ là lời hứa
đầu có dự báo, là giá của dự báo. Phép kiểm khả thi vẫn dùng giờ không dự báo, nên dưới luật C trần 30 giây, trong
1.471 khách được phục vụ ở 21 kịch bản có dự báo, 34 khách có giờ đón hiển thị ở lời hứa đầu muộn hơn giờ đón
muộn nhất; chỗ này cần sửa trước khi đưa dự báo vào bộ điều phối. Phép thử đưa dự báo vào bộ điều phối
(mục~\ref{muc:c4-dubao}) thì khác: nó đổi thời gian đi mà bộ điều phối nhận, còn FleetPy vẫn chạy theo hệ số đo.
Số khung lấy trung bình là $K$: các cấu hình chính bật dự báo dùng $K = 2$, gồm khung 15 phút hiện tại và khung kế tiếp.
Phép thử giả dùng $K = 1$, cho hệ số bằng 1; khi tắt dự báo không đặt $K$.
% EV: F03 (chỉ phần "phát lại trên các lần chạy đã ghi"); FR:10-11; X3-12 (FFA:84-99); F06; F07 (FR:70); F12; X3-29 (preflight_world_v1_3.py:1-26, 113-133); R26-CH3-OPEN-QUESTIONS-CODEX.md §Q9; preflight_world_v1_3.py:85-89, 99-101; run_forecast.py:67-72

% -----------------------------------------------------------------------------------------------------
\section{Sổ theo từng khách}
\label{muc:so}

Sổ do RideBound ghi trong lúc chạy: mỗi khách có bản ghi lời hứa đầu và một bản ghi sửa cho mỗi lượt mà phương án
được chọn làm $x$, $c$ hay $z$ khác 0 ở bất kỳ chiều nào. Bản ghi sửa lưu các phép chiếu ứng với $v$, $e$, $p$,
độ lớn thay đổi trên mười chiều và $L$ sau lần sửa (Bảng~\ref{bang:vidu-so}); các hiệu có dấu $x$, $c$, $z$ chỉ dùng
cho giờ. Không có ngưỡng: thay đổi 1 mili giây cũng tạo bản ghi sửa, nên số bản ghi sửa không phải số lần một ứng
dụng thật báo lại cho khách.
% EV: A04; X3-26 (DELTA:64-160); X3-25 (VALID:562-597); LEDGER:285-314; X3-13 (SPEC:35-46)

\begin{table}[htbp]
\centering
\caption[Bản ghi sổ của khách A trong ví dụ dựng tay]{Hai bản ghi giờ trả của khách A ở Hình~\ref{hinh:bamoc} (dựng tay), với giả định 9:10:00 là lời hứa đầu,
công bố lúc 8:50:00, và $B = 30$ giây. Bản ghi lời hứa đầu không có lời hứa trước nên sổ quy ước $v = e = p$. Các
độ lớn và $L$ tính bằng giây.}
\label{bang:vidu-so}
\small
\begin{tabular}{@{}lllllrrrr@{}}
\toprule
\textbf{Bản ghi} & \textbf{Lúc ghi} & $v$ & $e$ & $p$ & $|x|$ & $|c|$ & $|z|$ & $L$\\
\midrule
Lời hứa đầu & 8:50:00 & 9:10:00 & 9:10:00 & 9:10:00 & 0 & 0 & 0 & 0\\
Sửa (ghép khách M) & 9:00:00 & 9:10:00 & 9:10:40 & 9:11:10 & 40 & 30 & 70 & 30\\
\bottomrule
\end{tabular}
\end{table}
% EV: ví dụ dựng tay; v = e = p ở bản ghi đầu theo SPEC:46; L = 0 + |c| theo LEDGER:285-288

Trong một bản ghi mô phỏng của luật C trần 120 giây, một khách có $v$ = 17:43:46,044, $e$ = 17:43:46,283 và
$p$ = 17:44:01,203. Phần dời $|p - e| = 14{,}92$ giây làm phần trần còn lại bằng $120 - 14{,}92 = 105{,}08$ giây.
Cả 14 bản ghi sửa sau có $c = 0$ và $|x| < 0{,}25$ giây, nên phần trần đã dùng không tăng.
% EV: round26-trace-rider.txt, Q3 SELECTED RIDER, entries 1-15; R26-CH3-OPEN-QUESTIONS-CODEX.md §Q3; C-120-d20181117-s10-r1-w17/transcript-00.ndjson, final checkpoint ordinal 1716, request rq-c48d442603e9ce8c056b029d6d50c497f4d224b8af469945692690e92a75afb0

Sau khi chạy, độ trễ so với lời hứa đầu được chia ba phần: phần tính cho điều phối, phần được hồ sơ lịch sử giải
thích và phần còn lại,
\begin{equation}
a - p_0 \;=\; q \;+\; f \;+\; (a - D_0 - q), \qquad q = \sum_{j} c_j, \qquad f = D_0 - p_0 .
\label{eq:so}
\end{equation}
Cộng ba phần thì $q$ và $D_0$ triệt tiêu, nên đẳng thức luôn đúng. Trong đó, $a$ là giờ trả thực tế trong mô
phỏng; $q$ cộng các $c$ giữ nguyên dấu (dời muộn dương, dời sớm âm), khác với $\sum |c|$ của luật C, và dưới luật
C, $|q| \le \sum |c| \le B$; $f$ là phần dự báo đã cộng sẵn vào lời hứa đầu, bằng 0 khi không có dự báo.
Theo~\eqref{eq:dongnhat}, phần còn lại bằng $\sum_j x_j + \rho - f$ với $\rho = a - p_n$: thay đổi khi giữ nguyên
kế hoạch mà dự báo chưa tính, cộng độ lệch giữa giờ trả thực tế và lời hứa cuối $p_n$. Với khách A ở
Bảng~\ref{bang:vidu-so}, nếu $D_0$ = 9:10:30 và A tới đúng lời hứa cuối, $a$ = 9:11:10, thì độ trễ 70 giây gồm
$q = 30$, $f = 30$ và 10 giây còn lại, đúng bằng $\sum x + \rho - f = 40 + 0 - 30$.
% EV: X3-12 (FFA:84-97: q = Σ(p − e) có dấu, D0 = lời hứa hiển thị đầu); A12; A06; F09; ví dụ dựng tay X3-37 (b) (số học ở NOTES mục 4)

Chương~\ref{chuong:thucnghiem} cộng ba phần trên các khách có $a - p_0 > 60$ giây, với $p_0$ là lời hứa giờ trả
đầu không dự báo (không dùng $D_0$). Đây là phép tính theo mô hình, không kết luận ai có lỗi: ``phần tính cho điều
phối'' chỉ nói rằng mô hình gán phần này cho quyết định điều phối khi so với kế hoạch giữ nguyên. Phần được hồ sơ
lịch sử giải thích phụ thuộc dự báo đang dùng; phần còn lại là phần dự báo đó chưa giải thích được, không phải phần
``không thể biết trước''.
% EV: X3-12 (FFA:134-138: nhóm A − P0 > 60 s); F09; K04; A01 (lưu ý D03:222-224); K19

% -----------------------------------------------------------------------------------------------------
\section{Cài đặt, kiểm thử và tái lập}
\label{muc:caidat}

\subsection{Kiến trúc và trao đổi với bộ mô phỏng}
\label{muc:kientruc}

RideBound viết bằng C\# trên .NET 10, chia thành các dự án theo tầng (Bảng~\ref{bang:kientruc}). Tầng nghiệp vụ
không tham chiếu dự án hay gói thư viện ngoài nào, OR-Tools chỉ có trong cầu nối bộ giải, và một bộ kiểm thử kiến
trúc kiểm các quy tắc này; nhờ vậy phép kiểm trần, ba mốc và sổ không phụ thuộc bộ mô phỏng hay bộ giải.
% EV: X3-14 (các tệp .csproj: net10.0, ProjectReference, PackageReference; DependencyRuleTests.cs:8-31, 137-166)

\begin{table}[htbp]
\centering
\caption[Các thành phần chính của RideBound và bộ chuyển đổi FleetPy]{Các thành phần chính của RideBound và bộ chuyển đổi FleetPy}
\label{bang:kientruc}
\small
\begin{tabular}{@{}p{4.6cm}p{10.2cm}@{}}
\toprule
\textbf{Thành phần} & \textbf{Vai trò}\\
\midrule
Domain (tầng nghiệp vụ) & Yêu cầu và giai đoạn của khách, lời hứa, khóa, sổ, phép kiểm trần\\
Application (tầng ứng dụng) & Chiếu lịch theo thời gian đi, tính ba mốc giờ trả, kiểm phương án theo khóa và luật
lời hứa\\
Algorithms (thuật toán) & Bộ sinh chỉ chèn thêm, đánh giá phương án theo luật lời hứa và luật phục hồi, bộ chọn\\
Solvers.OrTools (cầu nối bộ giải) & Cầu nối tới bộ giải CP-SAT của OR-Tools 9.15\\
Contracts (định dạng bản tin) & Định dạng các bản tin trao đổi với bộ mô phỏng\\
Runner (bộ chạy) & Nhận bản tin, ghép các thành phần trên và trả quyết định\\
Benchmarking (chuẩn bị dữ liệu) & Chuẩn hóa dữ liệu Manhattan của FleetPy thành kịch bản thử\\
Bộ chuyển đổi FleetPy (Python) & Phần điều khiển đội xe của FleetPy 1.0.2: gửi sự kiện cho bộ chạy, áp quyết định
vào mô phỏng\\
\bottomrule
\end{tabular}
\end{table}
% EV: X3-14 (Google.OrTools 9.15.6755 chỉ ở RideBound.Solvers.OrTools.csproj:9); X3-22 (OrTools:3, CP-SAT); X3-15; NORM; B01

Bộ chuyển đổi viết bằng Python đóng vai phần điều khiển đội xe của FleetPy 1.0.2~\cite{engelhardt2026}, ghim theo
phiên bản; nó chạy bộ chạy RideBound như một chương trình con, hai bên nhắn tin qua đầu vào, đầu ra chuẩn, mỗi bản
tin một dòng JSON (Hình~\ref{hinh:bantin}). Bộ chạy chỉ nhận lô mới sau khi bộ mô phỏng báo đã áp dụng quyết định
trước, nên các quyết định được áp dụng đúng thứ tự. Bản tin của mỗi lần chạy được lưu, và mọi con số ở
Chương~\ref{chuong:thucnghiem} được tính từ đó.
% EV: B01; J18 (FLEET:474); X3-15 (README:1-12; runner_client.py:133-138; Program.cs:103-106; RunnerSession.cs:448-455); X3-13 (SPEC:28-37); A26

\begin{figure}[htbp]
\centering
\begin{tikzpicture}[y=1cm,font=\small,
  ben/.style={draw=xamvien,fill=xamnhat,rounded corners=2pt,align=center,inner sep=4pt,text width=4.3cm},
  tin/.style={-{Latex[length=2mm]},thick},
  nhan/.style={midway,above,font=\footnotesize,align=center}]
\node[ben] (L) at (0,0) {FleetPy 1.0.2\\và bộ chuyển đổi (Python)};
\node[ben] (R) at (9.5,0) {Bộ chạy RideBound\\(C\#, .NET 10)};
\draw[xamvien,thick] (L.south)--(0,-12.5);
\draw[xamvien,thick] (R.south)--(9.5,-12.5);
\draw[tin] (0,-1.8)--(9.5,-1.8) node[nhan]{chào hỏi; khởi tạo lần chạy\\(kịch bản, cấu hình luật)};
\draw[tin] (9.5,-2.9)--(0,-2.9) node[nhan]{đã nhận chào hỏi; đã khởi tạo};
\draw[dashed,xamvien] (-0.6,-3.35) rectangle (12.2,-9.1);
\node[anchor=north west,font=\footnotesize\itshape] at (-0.5,-3.45) {mỗi lượt quyết định};
\draw[tin] (0,-5.1)--(9.5,-5.1) node[nhan]{lô sự kiện: yêu cầu mới, xác nhận đặt xe,\\tiến độ xe, đón, trả,
thời gian đi mới};
\node[ben,text width=3.9cm,font=\footnotesize] at (9.5,-6.15) {sinh phương án, kiểm luật, chọn, ghi sổ
(Bảng~\ref{bang:luotquyetdinh})};
\draw[tin] (9.5,-7.7)--(0,-7.7) node[nhan]{quyết định: kế hoạch mới của các xe, lời hứa};
\draw[tin] (0,-8.7)--(9.5,-8.7) node[nhan]{đã áp dụng quyết định};
\draw[tin] (0,-10.15)--(9.5,-10.15) node[nhan]{yêu cầu trạng thái};
\draw[tin] (9.5,-11.15)--(0,-11.15) node[nhan]{trạng thái, gồm sổ theo từng khách};
\draw[tin] (0,-12.15)--(9.5,-12.15) node[nhan]{kết thúc: bộ chạy dừng, không trả lời};
\end{tikzpicture}
\caption[Các bản tin giữa bộ chuyển đổi FleetPy và bộ chạy RideBound]{Các bản tin chính giữa bộ chuyển đổi FleetPy và bộ chạy RideBound trong một lần chạy (đã giản lược). Mỗi
bản tin là một dòng JSON; phần trong khung nét đứt lặp lại ở mỗi lượt quyết định.}
\label{hinh:bantin}
\end{figure}
% EV: X3-15 (ProtocolPrimitives.cs:104-121; RunnerSession.cs:177-214: chỉ xử lý hello, initializeRun, eventBatch, decisionApplied, checkpoint, restore; shutdown làm bộ chạy dừng, không trả lời; finalizeRun/runSummary chỉ được khai báo; runner_client.py:158-273); X3-13 (SPEC:35-37)

\subsection{Giới hạn tính toán và thời gian chạy}
\label{muc:gioihan}

Ngoài giới hạn 2 yêu cầu mỗi lượt và 100 phương án mỗi xe, các bước sinh, kiểm và giải đều có giới hạn số bước
tính, và bộ giải dùng hạt giống ngẫu nhiên (số khởi tạo phần ngẫu nhiên) cố định bằng 7 để chạy lại cho cùng kết
quả. Vì vậy bộ sinh không xét hết mọi phương án; mục~\ref{muc:c4-hople} báo các bộ đếm của bộ sinh phương án cho biết
giới hạn có cắt bớt phương án hay không. Phép kiểm ở mỗi phương án duyệt mười chiều cho mỗi khách có lời hứa của xe và đếm đảo thứ tự
trên từng cặp điểm dừng, nên không thể coi khối lượng kiểm là tuyến tính theo độ dài kế hoạch; nó chưa được đo
riêng.
% EV: A24 (WP4:9-23); X3-17; X3-32 (GEN:54-79); N-25; X3-05 (BUDGET:76-99); X3-04 (VALID:233-239); X3-26 (DELTA:190-215)

Khóa luận chỉ báo thời gian thực của cả lần chạy: trên tập tắc đường cố định (16 lần chạy mỗi cấu hình), trung vị là
730, 667 và 641 giây với ghép tự do, điều phối không trần và luật C trần 30 giây. Thời gian này gồm mô phỏng FleetPy
và thời gian chờ RideBound trả quyết định, không phải thời gian của một lượt quyết định. Chi phí riêng của phép kiểm trần
không được báo, nên các số này không dùng để so chi phí tính toán.
% EV: H13 (EL17:2-6); X3-19; B17; R26-CH3-OPEN-QUESTIONS-CODEX.md §Q6; wp14r_scientific_protocol.py:801-803, 824-829; wp14r_supervised_process.py:297, 329-330, 979-985

\subsection{Kiểm tra kết quả}
\label{muc:kiemtra}

Kết quả được kiểm ở bốn mức.
\begin{itemize}[leftmargin=1.2cm,itemsep=2pt,topsep=3pt]
\item \textbf{Kiểm thử phần mềm:} chạy lại trên đúng mã nguồn đã dựng bộ chạy, cả 1.025 trong 1.025 kiểm thử tự động
đạt trong một lần chạy đầy đủ ngày 13/10/2026, với kiểu xuống dòng LF và một tệp chưa ghi nhận trong cây làm việc;
chi tiết ở Phụ lục~\ref{phuluc:C}.
\item \textbf{Kiểm tính nhất quán của sổ:} ở cả 382.006 bản ghi sửa của 185 lần chạy (170 lần chạy của năm cấu hình
trên tập kiểm tra chính và 15 lần chạy luật C trần 30 giây trên 15 kịch bản hiệu chỉnh), $|x|$, $|c|$, $|z|$ tính
lại từ $v$, $e$, $p$ trong sổ khớp đúng các trường độ lớn. Phép kiểm này không kiểm $e$ có được tính đúng hay không.
\item \textbf{Tính lại độc lập:} một chương trình viết riêng, độc lập về mã nhưng dùng cùng đặc tả chỉ số, tính lại
13 hạng mục số liệu từ bản ghi thô và một điều kiện luôn phải đúng, không lệch hạng mục nào; các số đếm chính được
ít nhất hai bộ đọc viết độc lập tính lại và khớp trong phạm vi đã kiểm. Hai mức kiểm này được làm bằng các chương
trình đọc riêng, không phải kiểm thủ công; chúng cho biết con số khớp bản ghi, không cho biết mô phỏng giống thực tế
đến đâu.
\item \textbf{Tái lập:} tập kiểm tra chính, tập chỉnh trần buổi sáng, các lần chạy trần 120 giây và tám lần chạy chiều
tuyết dùng cùng một bản bộ chạy đã biên dịch và lưu lại; tập tắc đường cố định dùng một bản cũ hơn. Bản dựng lại, nếu có,
phải đối chiếu với kết quả đã lưu trước khi dùng. Các lượt ưu tiên chi phí dùng bản riêng; bốn cặp kiểm
hành vi không thay thế việc kiểm đồng nhất toàn bộ hai bản chạy (Phụ lục~\ref{phuluc:D}).
\end{itemize}
% EV (kiểm thử): TESTRUN-45f8026-2026-10-12.md (mã nguồn bộ chạy; lượt 3 và 3b; điều kiện LF và tệp chưa ghi nhận); A28
% EV (sổ): A26 (T38:106; T38I:68); B10 (170 + 15); RUN38:1-8 (15 lần chạy C-30 trên 15 kịch bản hiệu chỉnh)
% EV (tính lại): C19 (T38I:79-97: 13 hạng mục + 1 bất biến); ADDENDUM-1:15-19; C04; E01; K41; A26 lưu ý; I01
% EV (tái lập): A28 (RUNW07:40-41); X3-10 (run_w07.py:8; run_w07x.py:7; run_t38c120.py:3-8); B17; NEW:13; I13; R26-CH3-OPEN-QUESTIONS-CODEX.md §Q8 (bản bộ chạy chiều tuyết)

Chương~\ref{chuong:thucnghiem} dùng các thành phần trên để đo luật C được gì và mất gì trên mô phỏng FleetPy, so
với điều phối không trần, luật hạn chót và cách đệm lời hứa, tức cộng thêm một khoảng cố định vào giờ hứa cho mọi
khách.
 của main.tex.
% Bản 1 (07-chuong3.tex) và bản trung gian 2 (07-chuong3-v2.tex) được giữ nguyên. Lý do từng thay đổi:
% reviews/CH3-REVISION-LOG.md. Ghi chú v3: 07-chuong3-v3-NOTES.md; thay đổi v4: 07-chuong3-v4-CHANGES-v3.md.
% Mỗi dòng "% EV:" trỏ tới EVIDENCE-REGISTER.md (A01…K52), EVIDENCE-REGISTER-ADDENDUM-1.md (N-01…N-30,
% L1a-01…09), EVIDENCE-REGISTER-ADDENDUM-2-CH3.md (X3-01…X3-41) hoặc các nguồn R26/TESTRUN ghi tại chỗ.
% =====================================================================================================

% Môi trường "Nhận xét" chưa có trong main.tex; định nghĩa có điều kiện ở đây (nên chuyển lên phần mở đầu).
\ifcsname nhanxet\endcsname\else
  \theoremstyle{definition}
  \newtheorem{nhanxet}{Nhận xét}[chapter]
\fi

\chapter{Phương pháp: luật C, lời hứa theo dự báo và sổ theo từng khách}
\label{chuong:phuongphap}

Khi đặt xe đi chung, hành khách được báo giờ xe sẽ trả mình tới nơi; khóa luận gọi đó là lời hứa giờ trả, gọi
tắt là lời hứa. Giờ trả có thể bị dời vì bộ điều phối ghé đón thêm khách, hoặc vì đường đông hơn lúc hứa hay xe
đi nhanh, chậm hơn dự kiến. Luật C, nội dung chính của chương, chỉ giới hạn loại thứ nhất. Ở mỗi lượt quyết định,
hiệu giữa giờ trả khi
chọn phương án mới và khi xe giữ nguyên kế hoạch là phần dời do điều phối; luật C loại phương án làm tổng độ lớn
phần này qua cả chuyến vượt một trần. Nhịp tắc đường quen thuộc trong ngày được cộng trước vào lời hứa đầu nhờ dự
báo; phần còn lại được ghi trong sổ theo từng khách (per-rider ledger). Các mục đi theo thứ tự đó, sau phần phát
biểu bài toán.
% EV: A01; A02; A06; F01; X3-12 (FFA:134-138, ba phần của sổ)

% -----------------------------------------------------------------------------------------------------
\section{Phát biểu bài toán và tổng quan hệ thống}
\label{muc:baitoan}

\subsection{Thuật ngữ và vòng đời của một yêu cầu}
\label{muc:thuatngu}

Bảng~\ref{bang:thuatngu} giải thích các thuật ngữ; tên tiếng Anh chỉ ghi trong bảng. Mọi thực nghiệm chạy trên
bộ mô phỏng FleetPy~\cite{engelhardt2026} với dữ liệu chuẩn Manhattan (mục~\ref{muc:kientruc}).
% EV: B01; J18 (FLEET:1-12, 474)

\begin{table}[htbp]
\centering
\caption[Thuật ngữ chính của phương pháp]{Thuật ngữ dùng trong chương}
\label{bang:thuatngu}
\small
\begin{tabular}{@{}p{4.3cm}p{10.5cm}@{}}
\toprule
\textbf{Thuật ngữ} & \textbf{Nghĩa trong khóa luận}\\
\midrule
Lượt quyết định (decision epoch) & Một lần bộ điều phối nhận một lô sự kiện và trả quyết định cho cả đội xe.\\
Phương án (candidate plan) & Một kế hoạch (dãy điểm đón, điểm trả theo thứ tự) có thể chọn cho một xe.\\
Phương án giữ nguyên kế hoạch (keep-route option) & Xe đi tiếp kế hoạch cũ, không nhận thêm khách.\\
Bộ sinh chỉ chèn thêm (insertion-only generator) & Chèn điểm đón, điểm trả của yêu cầu mới vào kế hoạch của
một xe, không bỏ, không đảo điểm cũ.\\
Bộ chọn (selector); bộ giải (solver) & Chọn mỗi xe một phương án theo thứ tự tiêu chí (lexicographic order):
chỉ khi hòa ở tiêu chí trước mới xét tiêu chí sau; bộ giải của OR-Tools~\cite{ortools} giải bài toán chọn.\\
Lời hứa; lời hứa đầu $p_0$ & Giờ đón, giờ trả báo cho khách; lời hứa đầu công bố khi khách xác nhận đặt xe;
$p_0$ là giờ trả của nó, chưa cộng dự báo.\\
Công bố; bản ghi sửa & Ghi lời hứa mới vào sổ và gửi trong bản tin quyết định; mỗi lần sau lời hứa đầu là một
bản ghi sửa.\\
Khóa (lock); chiều cam kết & Khóa: phần lời hứa không được đổi sau khi xác nhận. Chiều cam kết: thuộc tính của
lời hứa được theo dõi riêng, như giờ trả, việc đổi xe.\\
Trần $B$ (cap); phần đã dùng $L_r$ & Tổng phần dời do điều phối tối đa của một khách cả chuyến; phần khách $r$
đã dùng.\\
Quyết định bị ép (forced decision); luật phục hồi (recovery rule) & Lượt quyết định có ít nhất một xe bị ép:
luật lời hứa đã loại cả phương án giữ nguyên kế hoạch của xe, và xe chọn phương án được giữ nhờ luật phục hồi,
có ít nhất một khách được miễn trừ; luật phục hồi quy định phương án nào được giữ khi đó.\\
Hệ số tắc đường & Số nhân chung cho cả thành phố áp lên mọi thời gian đi, đổi theo khung 15 phút; hệ số 1,2 làm
thời gian đi dài gấp 1,2 lần hệ số 1.\\
Kịch bản & Một ngày, một mẫu nhu cầu, một khung 2 giờ của dữ liệu Manhattan: 108 yêu cầu, 8 xe 4 chỗ.\\
\bottomrule
\end{tabular}
\end{table}
% EV (lượt, phương án, bộ sinh, bộ chọn): X3-01; SPEC:31-34; X3-03; X3-07; X3-22; A14
% EV (lời hứa, công bố, khóa, chiều, trần): A13; X3-25 (fleet_control.py:882-893); A07; A04; A06; K43
% EV (bị ép, phục hồi, hệ số, kịch bản): X3-30; X3-09; B06; STRETCH:3-4; B02; B03; GLOSSARY-VI-v2.md §E (bị ép); C4A-10

Một yêu cầu mới chờ được xét. Khi phương án được chọn của một xe chứa nó, yêu cầu được nhận và khách được mời,
nhưng chưa có lời hứa. Khi khách xác nhận đặt xe, lời hứa đầu được công bố và $L_r = 0$; khách chờ đón, lên xe
rồi được trả. Yêu cầu chưa được nhận được xét lại ở lượt sau nếu còn phương án hợp lệ hay chưa được xét vì giới
hạn số yêu cầu mới, và bị từ chối khi mọi phương án chứa nó đều bị loại. Khóa và luật lời hứa áp từ lúc xác nhận
tới lúc khách xuống xe.
% EV: X3-20 (VALID:272-336; RollingCostPolicy.cs:173-226; fleet_control.py:878-924); A13; X3-32

\subsection{Bài toán ở mỗi lượt quyết định}
\label{muc:baitoan-phatbieu}

Bộ chạy tiến theo đồng hồ sự kiện: yêu cầu đến, xe tới mốc vật lý và mốc cập nhật tắc đường đều có thể gọi một
lượt quyết định. Xác nhận đặt xe và cập nhật trạng thái có thể tạo thêm lượt ở cùng thời điểm mô phỏng; không có
chu kỳ điều phối cố định. Bộ chuyển đổi gom các sự kiện từ lần gọi trước (yêu cầu mới, xác nhận đặt xe, tiến độ xe,
đón, trả, thời gian đi mới) thành một lô.
% EV: R26-CH3-OPEN-QUESTIONS-CODEX.md §Q1; actual_fleetpy_medium_preflight.py:528-543, 569-637; fleet_control.py:395-451; preflight_world_v1_1.py:40-58; preflight_world_v1_2.py:18-20

Khung nhu cầu dài 2 giờ, nhưng lần chạy kết thúc khi hết yêu cầu và hành trình đang hoạt động. Trong 34 lần chạy
luật C trần 120 giây của tập kiểm tra chính, thời lượng mô phỏng từ 118,02 đến 154,93 phút.
% EV: R26-CH3-OPEN-QUESTIONS-CODEX.md §Q2; round26-trace-rider.txt, Q2 C-120 MAIN-PANEL SIMULATED LENGTHS; actual_fleetpy_medium_preflight.py:580-581, 639-683

Bài toán ở mỗi lượt như sau.
\begin{itemize}[leftmargin=1.2cm,itemsep=2pt,topsep=3pt]
\item \textbf{Đầu vào:} trạng thái xe (vị trí, kế hoạch, sức chứa, khách trên xe, khách chờ đón); sổ; yêu cầu
đang chờ với điểm đón, điểm trả, cửa sổ đón (từ giờ đón sớm nhất tới giờ đón muộn nhất), thời gian ngồi xe tối
đa; thời gian đi lúc này, coi là giữ nguyên cho phần đường còn lại.
\item \textbf{Quyết định:} nhận, để lại hay từ chối từng yêu cầu đang chờ; mỗi xe một phương án thỏa các ràng
buộc cứng (Bảng~\ref{bang:rangbuoc}).
\item \textbf{Đầu ra:} lời hứa đầu và các bản ghi sửa (mục~\ref{muc:so}).
\item \textbf{Mục tiêu:} thứ tự tiêu chí của Bảng~\ref{bang:rangbuoc}; các chỉ số đánh giá ở
Chương~\ref{chuong:thucnghiem}.
\end{itemize}
% EV: X3-21 (fleet_control.py:395-429); X3-01; X3-02; A05; X3-20; B03; A06; A07; A14

Bộ chọn giải một quy hoạch nguyên, với $y_{k\pi} \in \{0, 1\}$ bằng 1 khi xe $k$ chọn phương án $\pi$ trong
tập $P_k$ của nó:
\begin{equation}
\sum_{\pi \in P_k} y_{k\pi} = 1 \;\;\text{với mọi xe } k, \qquad
\sum_{k} \sum_{\pi \in P_k,\; u \in \pi} y_{k\pi} \le 1 \;\;\text{với mọi yêu cầu mới } u ,
\label{eq:chon}
\end{equation}
tức mỗi yêu cầu mới thuộc phương án được chọn của nhiều nhất một xe. Các tiêu chí được giải lần lượt, mỗi lần giữ
các tiêu chí trước ở giá trị tốt nhất đã tìm. Bước chọn này theo lối gán xe--chuyến động của Alonso-Mora và cộng
sự~\cite{alonsomora2017} (mục~\ref{muc:c2-baitoan}).
% EV: X3-22 (OrTools:3, 96-125, 180-189, 343-365); X3-34 (HZ:341-345); J17

\begin{table}[htbp]
\centering
\caption[Ràng buộc cứng và thứ tự tiêu chí của bộ chọn]{Ràng buộc cứng của một phương án (nửa trên) và thứ tự tiêu chí của bộ chọn (nửa dưới): so tiêu chí 1
trước, chỉ khi hòa mới xét tiêu chí 2, và cứ thế. Giá trị là cấu hình của các tập thử, không phải chuẩn ngành.}
\label{bang:rangbuoc}
\small
\begin{tabular}{@{}p{2.4cm}p{7.6cm}p{4.8cm}@{}}
\toprule
\textbf{Nhóm} & \textbf{Nội dung} & \textbf{Vai trò}\\
\midrule
Ràng buộc dịch vụ & Đón trong vòng 10 phút kể từ lúc gọi; ngồi xe không quá 1,5 lần thời gian đi thẳng; mỗi
xe tối đa 4 khách & Phương án đổi kế hoạch phải thỏa chặt; phương án giữ nguyên được nới phần vượt thời gian
mà kế hoạch cũ đã có\\
% EV: B03; X3-23 (GEN:161-246; ServiceQualityAllowance.cs:61-94)
Khóa & Không đổi xe, điểm đón, giờ đón; điểm trả không đổi nhờ trần 0. Giờ đón mới phải bằng giờ đón khi giữ
nguyên kế hoạch & So với kế hoạch giữ nguyên, nên riêng tắc đường không làm vi phạm khóa\\
% EV: A07; A08; A10; X3-06 (LOCK:42-60, 103-161, 204-207; CFG C-30.json:66-67, 94-98)
Luật lời hứa & Luật C: $L_r + |p_r - e_r| \le B$ cho mọi khách có lời hứa của xe, hoặc một luật ở
Bảng~\ref{bang:luat} & Lọc trước bộ giải\\
% EV: A06; A09; A11; HVA:62-66
\midrule
Tiêu chí 1 & Ít phương án vi phạm được giữ lại nhất (chỉ khi có xe bị ép) & Chỉ dùng chúng khi không còn cách
khác\\
Tiêu chí 2 & Nhận nhiều yêu cầu mới nhất & Ưu tiên phục vụ\\
Tiêu chí 3 & Mức dùng trần lớn nhất trên đội xe càng nhỏ càng tốt; nếu hòa, tổng phần dời do điều phối nhỏ nhất,
so lần lượt từng chiều cam kết (khi có chiều mang trần) & Ưu tiên ổn định\\
Tiêu chí 4 & Tổng chi phí kế hoạch nhỏ nhất, tức thời lượng còn lại của các tuyến dự kiến & Theo mô hình, không
phải chi phí vận hành thực\\
Tiêu chí 5 & Thứ tự mã phương án & Phá hòa, để chạy lại cho cùng kết quả\\
% EV: A14; X3-07 (SEL:271-330, 365-385, 510-522; HVA:537-636); DIM:19-32; A16; I09
\bottomrule
\end{tabular}
\end{table}

\subsection{Một lượt quyết định}
\label{muc:luotquyetdinh}

Hình~\ref{hinh:vongquyetdinh} và Bảng~\ref{bang:luotquyetdinh} mô tả một lượt. Ngoài sinh phương án và chọn,
mỗi lượt có ba bước cho lời hứa: tính các mốc giờ trả của từng khách, loại phương án vi phạm khóa hay luật lời
hứa trước khi gọi bộ giải, và ghi sổ sau khi chọn. Hình và giả mã dùng một số ký hiệu được định nghĩa ở các mục
sau: với mỗi khách có lời hứa, $v$ là giờ trả đã công bố gần nhất, $e$ là giờ trả nếu xe giữ nguyên kế hoạch, $p$ là
giờ trả theo phương án đang xét, còn $x$, $c$, $z$ là ba hiệu giữa chúng (mục~\ref{muc:bamoc}); $D$ là lời hứa theo
dự báo, tính qua giờ $S$ (mục~\ref{muc:dubao}).
% EV: X3-04 (VALID:201-239, 355-411, 494-520); X3-13; X3-25

\begin{figure}[htbp]
\centering
\begin{tikzpicture}[
  node distance=6mm,
  hop/.style={draw=xamvien,fill=xamnhat,rounded corners=3pt,align=center,
              inner sep=4pt,font=\small,minimum height=12mm,text width=2.75cm},
  ngoai/.style={draw=xamvien,dashed,rounded corners=3pt,align=center,
              inner sep=4pt,font=\small,minimum height=10mm,text width=4.2cm},
  mui/.style={-{Latex[length=2.2mm]},thick}
]
\node[hop] (a) {Nhận lô sự kiện\\mới};
\node[hop,right=of a] (b) {Sinh phương án chỉ chèn thêm, kiểm ràng buộc dịch vụ};
\node[hop,right=of b] (c) {Tính $v$, $e$, $p$ cho từng khách có lời hứa};
\node[hop,right=of c] (d) {Loại phương án vi phạm khóa hay luật lời hứa};
\node[hop,below=9mm of d] (e) {Bộ giải chọn\\một phương án\\mỗi xe};
\node[hop,left=of e] (f) {Kiểm lại cả đội xe; công bố lời hứa $p$; ghi sổ};
\node[ngoai,below=8mm of f] (g) {Sau khi chạy, trên bản ghi đã lưu: lời hứa hiển thị $D = \max(p, S)$};
\draw[mui] (a)--(b); \draw[mui] (b)--(c); \draw[mui] (c)--(d);
\draw[mui] (d)--(e); \draw[mui] (e)--(f);
\draw[mui,dashed] (f.west) -| node[pos=0.3,below,font=\footnotesize]{lượt sau} (a.south);
\draw[-{Latex[length=2.2mm]},dashed] (f.south)--(g.north);
\end{tikzpicture}
\caption[Một lượt quyết định của RideBound]{Một lượt quyết định của RideBound. Bộ điều phối công bố lời hứa không dự báo $p$. Trong các kết quả
chính, lời hứa theo dự báo $D$ được tính sau khi chạy, từ bản ghi đã lưu (ô nét đứt), và không quay lại làm đầu
vào của quyết định nào (mục~\ref{muc:dubao}).}
\label{hinh:vongquyetdinh}
\end{figure}
% EV: X3-03; X3-04; X3-22; X3-25; F03; X3-11 (DISP:1-8); FR:10-11

\begin{table}[htbp]
\centering
\caption[Giả mã một lượt quyết định]{Giả mã một lượt quyết định; dấu $\gets$ nghĩa là gán giá trị bên phải cho biến bên trái}
\label{bang:luotquyetdinh}
\small
\begin{tabular}{@{}r@{\hspace{6pt}}p{10.2cm}@{\hspace{8pt}}p{4.0cm}@{}}
\toprule
\multicolumn{3}{@{}p{14.9cm}@{}}{\textbf{Đầu vào:} trạng thái đội xe; sổ; lô sự kiện mới; luật lời hứa đang
dùng. \textbf{Đầu ra:} kế hoạch mỗi xe; yêu cầu được nhận, để lại hay từ chối; lời hứa mới; bản ghi sổ}\\
\midrule
1: & \textbf{hàm LượtQuyếtĐịnh}(trạng thái, lô sự kiện) & \\
2: & \quad áp lô sự kiện; ghi nhận lần đón muộn hơn giờ đón muộn nhất, nếu có & vị trí xe, đón, trả\\
3: & \quad chọn tối đa 2 yêu cầu đang chờ có giờ đón muộn nhất sớm nhất & yêu cầu khác để lượt sau\\
4: & \quad \textbf{với mỗi} xe $k$ \textbf{thực hiện} & \\
5: & \qquad \textbf{nếu} kế hoạch đang chạy vi phạm ràng buộc cấu trúc \textbf{thì} báo lỗi, dừng lần chạy &
ví dụ vượt sức chứa\\
6: & \qquad $P_k \gets \{$phương án giữ nguyên$\}$, kiểm dịch vụ có nới phần vượt do tắc đường & luôn giữ khi
cắt bớt\\
7: & \qquad thêm vào $P_k$ các phương án chèn yêu cầu ở dòng 3, thỏa chặt ràng buộc dịch vụ; giữ tối đa 100 &
bộ sinh chỉ chèn thêm\\
8: & \qquad \textbf{với mỗi} $\pi \in P_k$ và mỗi khách có lời hứa của xe: tính $v$, $e$, $p$; \textbf{nếu}
$\pi$ vi phạm khóa hay luật lời hứa \textbf{thì} loại $\pi$ & Bảng~\ref{bang:kiemtran}\\
9: & \qquad \textbf{nếu} phương án giữ nguyên bị loại chỉ vì lỗi hạn chót hay lỗi trần \textbf{thì} & xét phục hồi\\
10: & \qquad\quad giữ nó làm phương án được miễn trừ & phục hồi giữ nguyên kế hoạch\\
11: & \qquad\quad giữ thêm phương án đã bị loại không làm vi phạm đó tệ hơn & phục hồi không tệ hơn\\
12: & \qquad \textbf{nếu} phương án giữ nguyên bị loại mà không được miễn \textbf{thì} báo lỗi, dừng lần chạy &
xe kẹt\\
13: & \quad chọn $y$ theo~\eqref{eq:chon} và thứ tự tiêu chí; kiểm lại lựa chọn trên cả đội xe &
Bảng~\ref{bang:rangbuoc}\\
14: & \quad nhận yêu cầu có trong phương án được chọn; để lại hay từ chối yêu cầu khác & \\
15: & \quad khách vừa xác nhận: công bố lời hứa đầu, $L_r \gets 0$; khách có lời hứa mà $x$, $c$ hay $z$ khác 0:
ghi bản ghi sửa, $L_r \gets L_r + |c_r|$ & ghi vi phạm cho khách được miễn\\
16: & \quad \textbf{trả về} quyết định; chờ bộ mô phỏng xác nhận đã áp dụng & \\
17: & \textbf{kết thúc hàm} & \\
\bottomrule
\end{tabular}
\end{table}
% EV (dòng 2–3): X3-01; X3-24 (WP4:8; EventReducer.cs:13-23, 260-282); X3-32 (GEN:54-96)
% EV (dòng 4–7): X3-31 (GEN:171-182); X3-23; X3-27 (CandidatePortfolioRetainer.cs:102-113); X3-03 (GEN:548-575); A24
% EV (dòng 8–12): X3-04; A06; A07; A10; X3-09 (HVA:244-390); X3-31 (POL:336-339; RunnerSession.cs:941-947)
% EV (dòng 13–16): X3-22; X3-07; POL:370-401; X3-20; X3-25 (VALID:272-336, 562-597; RunnerSession.cs:956-987); X3-30; RunnerSession.cs:448-455

Luật lời hứa là bộ lọc cứng: phương án vi phạm bị loại trước khi bộ giải chạy, không được giữ kèm một điểm trừ.
Ngoại lệ là khi xe bị ép (dòng 9--11), chẳng hạn khi riêng việc đường đông hơn đã đẩy giờ trả vượt hạn chót của
luật hạn chót (mục~\ref{muc:doichung}); khi đó tiêu chí 1 chỉ dùng phương án vi phạm nếu xe không còn cách khác.
Bộ sinh không bỏ, không đảo điểm cũ nào, nên khách đã nhận không bị chuyển sang xe khác. Phần nới ở dòng 6 giúp
xe không kẹt vì tắc đường. Mỗi lần vượt như vậy, và mỗi lần khách thực tế lên xe muộn hơn giờ đón muộn nhất, được ghi
lại thay vì làm dừng lần chạy. Vì vậy cửa sổ đón là điều kiện để kiểm phương án, không phải bảo đảm đón đúng hạn.
% EV: A06; HVA:62-66; X3-09 (HVA:255-326); X3-07 (SEL:283-295); A10; X3-03 (GEN:548-575); TEX08:129-131; X3-23 (GEN:161-246; ServiceQualityAllowance.cs:5-18, 61-94; RunnerSession.cs:989-1009); X3-24 (WP4:8; EventReducer.cs:13-23, 260-282)

% -----------------------------------------------------------------------------------------------------
\section{Ba mốc giờ trả và phần dời do điều phối}
\label{muc:bamoc}

Giờ trả có thể đổi vì bộ điều phối ghép thêm khách, hoặc ngay cả khi xe giữ kế hoạch cũ. Để tách hai loại, với
mỗi khách có lời hứa và mỗi phương án, hệ thống so ba mốc giờ trả:
\begin{itemize}[leftmargin=1.2cm,itemsep=2pt,topsep=3pt]
\item $v$: giờ trả đã công bố cho khách ở lần gần nhất;
\item $e$: giờ trả tính lại nếu xe giữ nguyên kế hoạch cũ, từ trạng thái hiện tại với thời gian đi hiện tại;
\item $p$: giờ trả nếu chọn phương án đang xét, với cùng trạng thái và thời gian đi.
\end{itemize}
Ba hiệu sau tính bằng giây (sổ lưu tới mili giây):
\begin{equation}
x = e - v, \qquad c = p - e, \qquad z = p - v = x + c .
\label{eq:bamoc}
\end{equation}
Trong đó, $x$ là phần thay đổi khi giữ nguyên kế hoạch, $c$ là phần dời do điều phối, $z$ là phần khách thấy;
dễ nhớ, $e$ là giờ nếu xe cứ đi như cũ. Mốc $e$ không phải ý của khóa luận: Hörl và Zwick~\cite{horl2022} đã tính
lại giờ trả trên tuyến đang gán theo thời gian đi hiện tại ở mỗi lượt, và khóa luận dùng lại mốc này.
% EV: A01; A02; X3-04 (VALID:355-384); X3-13 (SPEC:39-50); LEDGER:188-207; J01

Mốc $e$ là giờ do mô hình tính, không phải giờ đến thật, nên $x$ gồm mọi thứ đổi khi giữ nguyên kế hoạch: thời
gian đi mới, tiến độ thực của xe, sự kiện vừa xảy ra. Phép tách $z = x + c$ chỉ là phép trừ giữa các giờ của mô
hình, không chứng minh quyết định điều phối thực sự gây ra đúng $c$ giây chậm trễ. Hình~\ref{hinh:bamoc} là một ví
dụ dựng tay: ghép khách M làm A tới muộn hơn lời hứa 70 giây, nhưng chỉ 30 giây là phần dời do điều phối.
% EV: A01 (không phải giờ đến thật, không nhận diện nhân quả); A02; X3-04 (VALID:355-361); K04; ví dụ dựng tay

\begin{figure}[htbp]
\centering
\begin{tikzpicture}[x=1cm,y=1cm,font=\small,
  diem/.style={draw=xamvien,fill=xamnhat,rounded corners=2pt,inner sep=3pt,minimum height=6mm},
  mui/.style={-{Latex[length=2mm]},thick}]
\node[anchor=east] at (2.6,4.0) {Giữ nguyên kế hoạch:};
\node[diem] (x1) at (3.4,4.0) {Xe};
\node[diem] (a1) at (5.6,4.0) {Trả A};
\draw[mui] (x1)--(a1);
\node[anchor=east] at (2.6,3.0) {Ghép thêm khách M:};
\node[diem] (x2) at (3.4,3.0) {Xe};
\node[diem] (b2) at (5.6,3.0) {Đón M};
\node[diem] (a2) at (7.8,3.0) {Trả A};
\node[diem] (c2) at (10.0,3.0) {Trả M};
\draw[mui] (x2)--(b2); \draw[mui] (b2)--(a2); \draw[mui] (a2)--(c2);
\draw[-{Latex[length=2mm]}] (0.8,0)--(12.2,0) node[right]{giờ};
\foreach \vitri/\tenmoc/\giomoc in {2.0/v/9:10:00, 6.8/e/9:10:40, 10.4/p/9:11:10}
  \draw (\vitri,0.12)--(\vitri,-0.12) node[below,align=center]{$\tenmoc$\\\giomoc};
\draw[mui,dashed] (2.0,0.55)--(6.8,0.55) node[midway,above]{$x = e - v = 40$ giây};
\draw[mui] (6.8,0.55)--(10.4,0.55) node[midway,above]{$c = p - e = 30$ giây};
\draw[mui,densely dotted] (2.0,1.45)--(10.4,1.45) node[midway,above]{$z = p - v = 70$ giây};
\end{tikzpicture}
\caption[Ba mốc giờ trả trong một ví dụ dựng tay]{Ví dụ dựng tay: khách A được hứa trả lúc 9:10:00 ($v$); lúc 9:00, xe đang chở A xét ghép thêm khách M.
Tính với thời gian đi hiện tại, giữ nguyên kế hoạch thì A tới lúc 9:10:40 ($e$), ghép M thì tới lúc 9:11:10
($p$). Nét đứt: $x$; nét liền: $c$; nét chấm: $z = x + c$. Luật C chỉ tính $|c|$.}
\label{hinh:bamoc}
\end{figure}

Phần dời do điều phối được tính theo độ lớn: sớm hay muộn hơn $e$ 20 giây đều tính 20 giây, vì luật đo lần sửa
lời hứa chứ không đo lợi hay hại; Fielbaum và Alonso-Mora~\cite{fielbaum2020} thì chỉ tính phần bất lợi. Trên mọi
lần chạy đã ghi, $c \ge 0$ tới mức mili giây; đây là quan sát, chưa phải chứng minh cho trường hợp tổng quát. Trên
tập kiểm tra chính, khóa luận đã kiểm 136 lần chạy của bốn cấu hình: điều phối không trần, luật C trần 30 và 120
giây, luật hạn chót 30 giây. Ở các lần chạy này, mỗi lần kéo giờ trả sớm hơn $e$ chỉ sớm 1 mili giây. Hệ quả: khi
bộ điều phối chỉ chèn thêm, phần tính hai chiều và không hoàn lại của luật C gần như không có việc; nó chỉ có tác
dụng khi kế hoạch của khách đã nhận được phép đổi, ví dụ khi dời khách đang chờ sang xe khác.
% EV: A04; X3-05 (DELTA:137-139, 259-262); J06 (FAM:293-296); N-05 (CORRECTED); ADD1:104-112; ADDENDUM-1 mục 0.1–0.2 (không nêu nguyên nhân của −1 ms)

% -----------------------------------------------------------------------------------------------------
\section{Luật C: trần phần dời giờ trả do điều phối cho từng khách}
\label{muc:luatc}

\begin{dinhnghia}[Luật C]
\label{dn:luatc}
Gọi $B$ là trần, tính bằng giây, và $L_r$ là tổng các $|c|$ đã tính cho khách $r$ kể từ lúc $r$ xác nhận đặt xe,
với $L_r = 0$ lúc xác nhận. Một phương án của một xe chỉ được giữ lại nếu với mọi khách $r$ có lời hứa đang được
gán cho xe đó (chờ đón hay đang trên xe),
\begin{equation}
L_r + |p_r - e_r| \le B .
\label{eq:luatc}
\end{equation}
Sau khi quyết định của lượt được chọn, với mỗi khách $r$ có lời hứa mà phương án được chọn làm $x$, $c$ hay $z$
khác 0, $L_r$ được cộng thêm $|p_r - e_r|$ của phương án đó, kể cả khi giờ trả công bố không đổi ($p_r = v_r$).
\end{dinhnghia}
% EV: A06 (BUDGET:41-105; LEDGER:285-288); A13; X3-05 (BUDGET:107-114; CFG C-30.json:17-23); X3-25 (VALID:272-281, 562-597)

Dấu $\le$ cho phép $L_r + |c|$ đúng bằng $B$. Phương án chỉ được xét không làm đổi $L_r$: mỗi phương án được
kiểm trên bản sao của sổ, chỉ quyết định cuối của lượt mới được ghi vào sổ thật. $B$ giống một ngân sách thời
gian của từng khách: $B - L_r$ là phần trần còn lại, không được cấp lại sau mỗi lần dời. Vì phần dời được cộng
qua cả chuyến, một luật cho mỗi lần dời tới 30 giây cho phép ba lần, tổng 90 giây, còn luật C với $B = 30$ giây thì
không. Ở Hình~\ref{hinh:bamoc}, với $L_A = 0$ và $B = 30$ giây, phương án ghép M ($|c| = 30$ giây) vẫn được giữ và
A dùng hết trần. Khi phép kiểm không được miễn cho $r$ ở lượt nào, tổng $|c|$ của $r$ cả chuyến không vượt $B$.
% EV: X3-05 (BUDGET:88); X3-25 (HVA:161-174; RunnerSession.cs:956-987; VALID:640-644: miễn trừ vẫn cộng vào sổ); A06; K43; E08; J05; CFG C-30.json:3-6 (một chính sách cho mọi khách)

\begin{table}[htbp]
\centering
\caption[Giả mã phép kiểm luật C cho một phương án]{Giả mã phép kiểm luật C cho một phương án (chỉ viết chiều giờ trả)}
\label{bang:kiemtran}
\small
\begin{tabular}{@{}r@{\hspace{6pt}}p{9.6cm}@{\hspace{8pt}}p{4.6cm}@{}}
\toprule
\multicolumn{3}{@{}p{14.9cm}@{}}{\textbf{Đầu vào:} phương án $\pi$ của xe $k$; trần $B$; $L_r$ của mỗi khách
có lời hứa; trạng thái và thời gian đi của lượt. \textbf{Đầu ra:} giữ hay loại $\pi$, kèm bản ghi chặn}\\
\midrule
1: & \textbf{hàm KiểmTrần}($\pi$, $k$, $B$) & \\
2: & \quad \textbf{với mỗi} khách $r$ có lời hứa được gán cho xe $k$, theo thứ tự mã khách \textbf{thực hiện} &
chờ đón hay trên xe\\
3: & \qquad $e_r \gets$ giờ trả của $r$ khi chiếu lại kế hoạch cũ của $k$ từ trạng thái hiện tại & mốc giữ
nguyên kế hoạch\\
4: & \qquad $p_r \gets$ giờ trả của $r$ theo $\pi$, cùng trạng thái và thời gian đi; $c_r \gets p_r - e_r$ &
phần dời do điều phối\\
5: & \qquad \textbf{nếu} $L_r + |c_r| > B$ \textbf{thì trả về} loại $\pi$, kèm ($r$, $L_r$, $|c_r|$,
$L_r + |c_r|$, $B$) & dừng ở khách vi phạm gặp trước\\
6: & \quad \textbf{trả về} giữ $\pi$ & $L_r$ chỉ đổi nếu $\pi$ được chọn\\
7: & \textbf{kết thúc hàm} & \\
\bottomrule
\end{tabular}
\end{table}
% EV (dòng 2–4): X3-04 (VALID:233-239, 355-361; HVA:171); X3-27 (VALID:1048-1064)
% EV (dòng 5–6): X3-05 (BUDGET:6-12, 72-104); A06; A23; X3-25

Cài đặt theo dõi mười chiều cam kết: thay đổi giờ đón, giờ trả; số lần sửa đáng kể; đổi xe; khoảng dời và số lần
đổi của điểm đón, của điểm trả; số lần đảo thứ tự với khách cũ; số điểm dừng mới chèn trước điểm đón. Chỉ bốn
chiều có trần: giờ trả ($B$) và đổi xe, đổi điểm đón, đổi điểm trả (trần 0). Ở tập kiểm tra chính và tập buổi sáng,
phép kiểm dừng ở khách vi phạm gặp trước, nên mỗi bản ghi chặn (bản ghi lý do loại phương án) chỉ nêu lỗi của khách
đó.
% EV: A04 (DIM:3-32); X3-26 (DELTA:137-153); A25; A07; CFG C-30.json:7-89; X3-07 (SEL:319-330, 510-522); A23; X3-04 (VALID:222-239); WP4:1-25; X3-10

\subsection{Phương án giữ nguyên kế hoạch dưới luật C}
\label{muc:giunguyen}

Gọi \textbf{(N)} là giả thiết, áp cho lượt đang xét: (i) kế hoạch đang chạy không vi phạm ràng buộc cấu trúc
(sức chứa, đón trước trả, phần kế hoạch đã cố định), nên bộ sinh dựng được phương án giữ nguyên; (ii) giờ trả
theo phương án giữ nguyên được tính bằng đúng phép chiếu dùng cho $e$. RideBound thỏa (ii) vì dùng cùng một hàm
chiếu với cùng đầu vào, và luôn giữ phương án giữ nguyên khi cắt còn 100 phương án.
% EV: A22; X3-31 (GEN:171-182, 213-229); ServiceQualityAllowance.cs:66-73; X3-27 (VALID:355-360, 1048-1064; HVA:144-165; CandidatePortfolioRetainer.cs:102-113; WP4:15)

\begin{nhanxet}
\label{nx:giunguyen}
Giả sử (N) và mọi khách có lời hứa của xe có $0 \le L_r \le B$. Khi đó luật C không loại phương án giữ nguyên
kế hoạch.
\end{nhanxet}

\noindent\textit{Chứng minh.} Theo (N)(ii), với mọi khách $r$ có lời hứa của xe, $p_r = e_r$, nên
$L_r + |p_r - e_r| = L_r \le B$ và phép kiểm~\eqref{eq:luatc} cho qua. \hfill$\square$
% EV: A22 (theo cấu tạo, giả thiết N); A06; X3-27

Khóa và các trần 0 cũng so với kế hoạch giữ nguyên, nên phương án này qua cả chúng. Nếu (N) đúng ở mọi lượt thì
$L_r \le B$ được giữ: $L_r = 0$ lúc xác nhận, và nếu $L_r \le B$ ở đầu một lượt thì không có miễn trừ nào, phương
án được chọn đã qua~\eqref{eq:luatc}, nên $L_r$ mới vẫn không vượt $B$. Khi đó xe dùng luật C hay điều phối không
trần không bao giờ bị ép; tính chất này vì vậy không riêng của luật C.
% EV: A10 (LOCK:156-161); X3-06; X3-26 (véc-tơ c của phương án giữ nguyên bằng 0 ở mọi chiều); X3-09 (HVA:255-259); X3-25

Nhận xét đúng do cách đặt luật, không phải kết quả thực nghiệm; số liệu chỉ phù hợp với nó: trên tập kiểm tra
chính, 0 trong 20.101 quyết định bị ép với trần 30 giây và 0 trong 20.435 với trần 120 giây (mỗi quyết định ứng với
một lượt của cả đội xe).
% EV: A21; N-21; X3-16 (STG:43-44); N-07; X3-30 (RunnerSession.cs:1013-1027); SPEC:31-34

Nhận xét có ba giới hạn. Thứ nhất, nếu (N)(i) không thỏa, lần chạy dừng vì lỗi cấu trúc, không vì luật C; phần
vượt cửa sổ đón hay thời gian ngồi xe do tắc đường không làm (N) hỏng, nhờ phần nới. Thứ hai, với $v$ = 9:10:00 và
$e = p$ = 9:11:00, phương án giữ nguyên có $c = 0$ nhưng $z = 60$ giây: luật C không bảo đảm giờ đến. Thứ ba, với
$v = p$ = 9:10:00 và $e$ = 9:12:00, kéo giờ trả về đúng giờ đã hứa có $z = 0$ nhưng $|c| = 120$ giây, và $L_r$
tăng 120 giây dù giờ công bố không đổi. Tính chất và giới hạn thứ hai là hai mặt của một lựa chọn: luật C không
bị ép vì không tính phần thay đổi khi giữ nguyên kế hoạch, và vì thế cũng không giới hạn phần đó.
% EV: A22; X3-31 (GEN:171-182; RunnerSession.cs:941-947); X3-23; K05; X3-25 (VALID:562-597: bản ghi sửa khi x ≠ 0 dù p = v); ví dụ dựng tay

% -----------------------------------------------------------------------------------------------------
\section{Các luật đối chứng, bộ chọn và luật phục hồi}
\label{muc:doichung}

Bảng~\ref{bang:luat} định nghĩa bốn cấu hình mà Chương~\ref{chuong:thucnghiem} so với luật C. Trừ ghép tự do,
chúng dùng cùng bộ sinh, bộ kiểm, bộ chọn và luật phục hồi với luật C; tập tắc đường cố định chạy bằng bộ chạy
chưa có luật phục hồi.
% EV: A08; A09; A11; A17; A19; K50; HO:109; B17; NEW:13

\begin{table}[htbp]
\centering
\caption[Các cấu hình được so sánh]{Các cấu hình được so sánh; ký hiệu chỉ dùng trong bảng. $B$ là trần, tính bằng giây; với luật hạn chót,
$B$ là số giây giờ trả được phép muộn hơn lời hứa đầu $p_0$. Các tập thử được mô tả ở mục~\ref{muc:c4-thietlap}.}
\label{bang:luat}
\footnotesize
\setlength{\tabcolsep}{3pt}
\begin{tabular}{@{}p{2.2cm}p{1.1cm}p{3.1cm}p{2.3cm}p{2.6cm}p{3.5cm}@{}}
\toprule
\textbf{Tên trong khóa luận} & \textbf{Ký hiệu} & \textbf{Kiểm điều gì} & \textbf{Khóa} & \textbf{Khi xe bị
ép} & \textbf{Dùng ở tập nào}\\
\midrule
Điều phối không trần & U & Không kiểm giờ trả & Xe, điểm đón, giờ đón; điểm trả có trần 0 & Không bị ép khi (N)
đúng & Tập kiểm tra chính; tập chỉnh trần buổi sáng; tập tắc đường cố định\\
% EV: A07; A08; K50; A10; X3-16 (STG:42: 0 / 20.760); B10; B15; B17; NEW:11
Ghép tự do & FREE & Không kiểm lời hứa; chọn theo chi phí kế hoạch ở từng lượt & Không khóa & Không áp dụng &
Tập tắc đường cố định\\
% EV: A17; K50; HO:109 (b1-rolling-cost, không khóa); B17; X3-38
Luật C & C-$B$ & $\sum |c| \le B$ cả chuyến, cho mọi khách có lời hứa & Như U & Không bị ép khi (N) đúng
(Nhận xét~\ref{nx:giunguyen}) & Tập chính (30, 120); buổi sáng (30, 60, 90, 120, 150, 180, 240, 300); tắc đường
cố định (30, 60, 120); chiều tuyết (30)\\
% EV: A06; A22; B10; L1a-01; X3-28 (run_w07.py:5; run_w07x.py:3-4, 35-36); B17; B20; N-26 (chiều tuyết 0 / 5.067 quyết định bị ép)
Luật hạn chót & M$^{+}$-$B$ & $p - p_0 \le B$ ở mọi lần kiểm; chỉ chặn phía muộn & Như U & Có thể bị ép do tắc
đường; phục hồi không tệ hơn & Tập chính (30, 60); buổi sáng (tám mức như luật C)\\
% EV: A09; A10; A19; B10; X3-28; D04; LOCK:186-197; CFG Mplus-30.json:98
Luật trần độ dời khách thấy & V-$B$ & $\sum |z| \le B$ cả chuyến & Như U & Có thể bị ép do tắc đường; phục hồi
không tệ hơn & Tập chính (30); buổi sáng (30, 90, 120, 150, 180, 240)\\
% EV: A11; A19; B10; X3-28 (run_w07x.py:3-4, 35-36: V-90/120/150/180/240; run_w07.py:5: V-30)
\bottomrule
\end{tabular}
\end{table}

Luật hạn chót chỉ mượn thành phần hạn chót tính từ lời hứa đầu của Schulz và Pfeiffer~\cite{schulz2026}, vốn
dùng thời gian đi cố định và độ chùng mặc định 900 giây. Ở đây mỗi độ chùng bằng một mức trần của luật C trên cùng
tập, và với tắc đường đổi theo khung 15 phút, riêng tắc đường có thể làm luật này bị ép; số quyết định bị ép ở
Chương~\ref{chuong:thucnghiem} nói về một hạn chót chặt như vậy, không về phương pháp của họ.
% EV: J10 (SP:466-490, 568-577); J11; A08; A09; X3-28; X3-41 (D03:214-217); A10

Ba luật lời hứa khác nhau ở đại lượng được cộng. Đánh số các lần cập nhật lời hứa của một khách là
$1, \ldots, n$. Trong sổ, mỗi bản ghi sửa so với lời hứa ngay trước nó ($v_j = p_{j-1}$), nên các giờ ở giữa
triệt tiêu:
\begin{equation}
\sum_{j=1}^{n} (x_j + c_j) \;=\; \sum_{j=1}^{n} z_j \;=\; p_n - p_0 .
\label{eq:dongnhat}
\end{equation}
Luật hạn chót kiểm $p_n - p_0 \le B$, tức cộng có dấu cả hai phần; luật trần độ dời khách thấy cộng $|z_j|$; luật
C chỉ cộng $|c_j|$. Ở Hình~\ref{hinh:bamoc}, nếu 9:10:00 là lời hứa đầu, phương án giữ nguyên có
$e - p_0 = |z| = 40$ giây, vượt trần 30 giây của hai luật kia dù $c = 0$: riêng tắc đường đủ làm chúng bị ép.
% EV: A12; X3-33 (LEDGER:256-269, 298-314; SPEC:62-63); A09; A10 (LOCK:156-161); A11; X3-16 (U 0 / 20.760)

Tiền lệ trực tiếp là Hörl và Zwick~\cite{horl2022} (mục~\ref{muc:c2-rangbuoc}): ở đầu mỗi lượt họ tính mốc $e$,
và công thức 7 đặt giờ trả bắt buộc bằng giá trị lớn hơn giữa giờ cập nhật và giờ trả muộn nhất, như ràng buộc cứng
(ràng buộc 9--10), kèm phạt cho phần vượt giờ gốc (công thức 11). Mục đích của họ, tránh tình huống không còn thứ tự
trả khách khả thi, cũng là mục đích của Nhận xét~\ref{nx:giunguyen}. Đặt giờ trả muộn nhất là $p_0 + B$ cho hạn chót
nới $p \le \max(p_0 + B, e)$, một cách viết lại của khóa luận; nó cho phép $c \le \max(p_0 + B - e, 0)$, về 0 khi
tắc đường đẩy $e$ quá $p_0 + B$, và cũng không loại phương án giữ nguyên. Cấu hình hạn chót M$^{+}$ với luật phục
hồi ``không tệ hơn'' ở chương~4 chính là dạng này. Luật C khác ở ba chỗ:
một trần cho mọi khách có lời hứa, còn phần nới của họ chỉ tác dụng khi giờ cập nhật đã vượt giờ trả muộn nhất;
phần được phép dời, $|c| \le B - L_r$, không phụ thuộc tắc đường; và cộng qua cả chuyến, dù phần cộng dồn hiếm khi là
lý do loại phương án (mục~\ref{muc:c4-congdon}).
% EV: J01; J02 (HZ:481-497); J03; X3-36 (HZ:531-543, 565-606); J05; A22 (suy luận V17 R6); K02; K26; K27; N-15

Từ Fielbaum và Alonso-Mora~\cite{fielbaum2020} (mục~\ref{muc:c2-dolech}), khóa luận kế thừa ý tách phần thay đổi
do quyết định gán xe và kiểm soát nó trước khi quyết định, nhưng mỗi lần sửa được so với mốc $e$ chứ không so điểm
đầu với điểm cuối, nên hai lần dời ngược chiều không triệt tiêu nhau; chiều sớm cũng được tính; đại lượng là giờ
trả chứ không phải thời gian chờ; và tắc đường thay đổi chứ không chọn giờ thấp điểm.
% EV: J06; J07; J08; X3-35 (FAM:90-94, 189-195, 216-220, 293-296, 948-972); K30; K32

Vì đổi xe, đổi điểm đón, đổi điểm trả là các trần 0, phần đầu của tiêu chí 3 (mức dùng trần lớn nhất, tỉ số lớn
nhất giữa phần đã dùng và trần) luôn bằng 100\% ở các lượt đã ghi của tập kiểm tra chính (mục~\ref{muc:c4-hople}),
cả khi không trần; tiêu chí 3 vì vậy chỉ tác dụng qua tổng phần dời theo từng chiều. Điều phối không trần và luật C
cùng cấu hình và bộ chọn, chỉ khác trần giờ trả; chênh lệch giữa chúng là phần đi kèm việc thêm trần khi giữ nguyên
bộ chọn. Khóa luận còn xét thứ tự ưu tiên chi phí: giữ ưu tiên nhận yêu cầu và các ràng buộc, bỏ các mức xếp
hạng theo ổn định rồi so chi phí kế hoạch và tiêu chí phân xử. Phương án vẫn phải qua phép kiểm trần khi trần
được bật. Mục~\ref{muc:c4-vaitro} so hai thứ tự ở trần 30 giây và không trần; đổi thứ tự không làm thay đổi
định nghĩa hay cách cộng trần của luật C.
% EV: S15-03 (costfirst-diff, cấu hình và provenance); S15-01 (mục 0, 2).
% EV: A14; X3-07 (SEL:271-330; HVA:537-636); A15; N-25; A08; H01; H02; K18

Luật phục hồi giữ nguyên kế hoạch giữ phương án giữ nguyên làm phương án được miễn trừ nếu mọi lỗi của nó là lỗi
hạn chót hay lỗi trần. Luật phục hồi không tệ hơn mở rộng luật này: giữ thêm phương án đổi kế hoạch có mỗi lỗi,
cùng khách và cùng chiều, cũng là lỗi của phương án giữ nguyên và không nặng hơn. Lỗi khóa theo giai đoạn (khóa của
khách chờ đón hay trên xe) và lỗi ràng buộc dịch vụ không bao giờ được miễn. Mỗi khách được miễn được ghi một vi
phạm, trừ khi lời hứa, giai đoạn và mã lỗi không đổi so với lần trước. Tập kiểm tra chính, tập buổi sáng và tám
lần chạy chiều tuyết dùng luật thứ hai; cả hai là chi tiết cài đặt, không phải đóng góp của khóa luận.
% EV: A17; A18 (ADR:5066-5092); X3-09 (HVA:79-99, 255-326; VALID:413-420, 475-483, 537-544, 599-638, 791-818; LOCK:42-60); A19 (RUN38:9; WP4:7); X3-10; R26-CH3-OPEN-QUESTIONS-CODEX.md §Q8 (tám lần chạy chiều tuyết; wp4-tier3-v2.json:7)

% -----------------------------------------------------------------------------------------------------
\section{Lời hứa theo dự báo}
\label{muc:dubao}

Vì bộ điều phối dùng thời gian đi hiện tại, lời hứa có thể sớm hơn giờ trả thực tế khi đường đông dần, dù xe giữ
nguyên kế hoạch. Lời hứa theo dự báo, một cách làm đã có (mục~\ref{muc:c2-dubao}), cộng sẵn mức tăng tắc đường các
ngày trước cho thấy vào từng giờ đón, giờ trả được công bố.
% EV: A05; F10; J35; F01; X3-11 (DISP:1-8)

Thời gian đo bằng giây kể từ đầu khung giờ của kịch bản; $i(\tau) = \lfloor \tau / 900 \rfloor$ là khung 15 phút
chứa $\tau$. Gọi $m_i$, $i = 0, \ldots, 8$, là trung bình hệ số tắc đường của khung $i$ trên các ngày nguồn; từ
khung 8 hệ số được giữ ở giá trị cuối, như bộ mô phỏng cũng giữ, nên đặt $\tilde m_i = m_{\min(i, 8)}$. Với lúc
công bố $t$ và giờ công bố $p > t$ (giờ của phương án được chọn, chưa cộng dự báo), $S$ là nghiệm của
\begin{equation}
\int_{t}^{S} \frac{\mathrm{d}\tau}{g(\tau)} = p - t, \qquad
g(\tau) =
\begin{cases}
1 & \text{nếu } i(\tau) = i(t),\\
\tilde m_{i(\tau)} / \tilde m_{i(t)} & \text{nếu không,}
\end{cases}
\label{eq:keodan}
\end{equation}
và lời hứa hiển thị là $D = \max(p, S)$ (giả mã ở Bảng~\ref{bang:dubao}).
% EV: F01; X3-11 (DISP:3-8, 55-64; STRETCH:1-12, 22-52; preflight_world_v1_3.py:51-82)

Theo tốc độ lúc công bố, xe cần thêm $p - t$ giây. Trong khung chậm hơn khung hiện tại $g$ lần, mỗi giây đồng hồ chỉ
đi được phần đường mà ở tốc độ lúc công bố mất $1/g$ giây. Mức tắc đường lúc này đã có trong thời gian đi đang dùng, nên $g = 1$ ở khung
hiện tại; các ngày khác chỉ cho biết khung sau chậm hay nhanh hơn bao nhiêu lần. Phép lấy $\max$ bảo đảm $D \ge p$ ở
cùng lần công bố, không bảo đảm $D$ tăng dần qua các lần cập nhật. Ví dụ dựng tay: trong kịch bản khung
8:00--10:00, lúc $t$ = 9:05:00 hệ thống công bố $p$ = 9:25:00, tức còn 1.200 giây. Nếu các ngày nguồn có hệ số trung
bình 1,25 ở khung 9:00--9:15 và 1,50 ở khung 9:15--9:30, thì sau 9:15 $g = 1{,}50/1{,}25 = 1{,}2$. Tới 9:15 xe đi
được 600 giây quãng đường; 600 giây còn lại mất $600 \times 1{,}2 = 720$ giây, nên $S = D$ = 9:27:00, dài thêm 120
giây.
% EV: X3-11 (STRETCH:1-12; forecast_weight neo ở khung hiện tại); F01; DISP:60-64; ví dụ dựng tay X3-37 (a), tính bằng hàm stretch của dự án; số học ở 07-chuong3-v3-NOTES.md mục 4

\begin{table}[htbp]
\centering
\caption[Giả mã lời hứa theo dự báo]{Giả mã lời hứa theo dự báo cho một giờ đón hay giờ trả đã công bố}
\label{bang:dubao}
\small
\begin{tabular}{@{}r@{\hspace{6pt}}p{9.6cm}@{\hspace{8pt}}p{4.6cm}@{}}
\toprule
\multicolumn{3}{@{}p{14.9cm}@{}}{\textbf{Đầu vào:} lúc công bố $t$; giờ công bố $p$; ngày $d$ và khung giờ của
kịch bản; hệ số tắc đường theo khung của các ngày trước $d$. \textbf{Đầu ra:} lời hứa hiển thị $D$}\\
\midrule
1: & \textbf{hàm LờiHứaTheoDựBáo}($t$, $p$, $d$) & \\
2: & \quad $H \gets$ các ngày cùng loại với $d$ và trước $d$, bỏ ngày 15/11 & chỉ dùng ngày đã qua\\
3: & \quad \textbf{nếu} $H$ có ít hơn 2 ngày \textbf{hoặc} $p \le t$ \textbf{thì trả về} $p$ & không dự báo\\
4: & \quad $m_i \gets$ trung bình hệ số của khung $i$ trên $H$, $i = 0, \ldots, 8$ & khung tính từ đầu khung giờ\\
5: & \quad $W \gets p - t$;\; $\tau \gets t$;\; $i_0 \gets i(t)$ & quãng còn lại theo tốc độ hiện tại\\
6: & \quad \textbf{lặp} & \\
7: & \qquad $g \gets 1$ nếu $i(\tau) = i_0$, ngược lại $g \gets \tilde m_{i(\tau)} / \tilde m_{i_0}$ &
$\tilde m_i = m_{\min(i, 8)}$\\
8: & \qquad $h \gets$ thời gian từ $\tau$ tới hết khung $i(\tau)$ & \\
9: & \qquad \textbf{nếu} $h/g \ge W$ \textbf{thì} $S \gets \tau + W g$ và thoát vòng lặp & \\
10: & \qquad $W \gets W - h/g$;\; $\tau \gets$ đầu khung kế tiếp & \\
11: & \quad \textbf{trả về} $D \gets \max(p, S)$ & $D \ge p$ ở lần công bố này\\
12: & \textbf{kết thúc hàm} & \\
\bottomrule
\end{tabular}
\end{table}
% EV (dòng 2–4): F02; X3-11 (DISP:27-57; preflight_world_v1_3.py:51, 80-82)
% EV (dòng 5–11): F01; X3-11 (STRETCH:25-39, 49-52)

Ngày nguồn là các ngày cùng loại (ngày thường hay cuối tuần) trước ngày của kịch bản, ít nhất 2 ngày, và không bao
giờ là ngày tuyết 15/11/2018 (một lựa chọn của khóa luận; hệ thống thật cần quy tắc tự động nhận ra ngày bất
thường). Vì vậy trong dữ liệu 12--18/11/2018 chỉ ba ngày 14, 15, 16/11 có dự báo: 21 trong 34 kịch bản của tập
kiểm tra chính; 13 kịch bản ngày 17, 18/11 giữ lời hứa không dự báo.
% EV: F02 (DISP:10-16, 27-52); F04; N-16; N-28

Trong các kết quả chính, $D$ được tính lại trên các lần chạy đã ghi, gọi là phát lại (replay): xe nào chở ai giữ
nguyên, chỉ giờ báo cho khách đổi, và $D$ không quay lại sổ hay phép kiểm trần. $D_0 - p_0$, với $D_0$ là lời hứa
đầu có dự báo, là giá của dự báo. Phép kiểm khả thi vẫn dùng giờ không dự báo, nên dưới luật C trần 30 giây, trong
1.471 khách được phục vụ ở 21 kịch bản có dự báo, 34 khách có giờ đón hiển thị ở lời hứa đầu muộn hơn giờ đón
muộn nhất; chỗ này cần sửa trước khi đưa dự báo vào bộ điều phối. Phép thử đưa dự báo vào bộ điều phối
(mục~\ref{muc:c4-dubao}) thì khác: nó đổi thời gian đi mà bộ điều phối nhận, còn FleetPy vẫn chạy theo hệ số đo.
Số khung lấy trung bình là $K$: các cấu hình chính bật dự báo dùng $K = 2$, gồm khung 15 phút hiện tại và khung kế tiếp.
Phép thử giả dùng $K = 1$, cho hệ số bằng 1; khi tắt dự báo không đặt $K$.
% EV: F03 (chỉ phần "phát lại trên các lần chạy đã ghi"); FR:10-11; X3-12 (FFA:84-99); F06; F07 (FR:70); F12; X3-29 (preflight_world_v1_3.py:1-26, 113-133); R26-CH3-OPEN-QUESTIONS-CODEX.md §Q9; preflight_world_v1_3.py:85-89, 99-101; run_forecast.py:67-72

% -----------------------------------------------------------------------------------------------------
\section{Sổ theo từng khách}
\label{muc:so}

Sổ do RideBound ghi trong lúc chạy: mỗi khách có bản ghi lời hứa đầu và một bản ghi sửa cho mỗi lượt mà phương án
được chọn làm $x$, $c$ hay $z$ khác 0 ở bất kỳ chiều nào. Bản ghi sửa lưu các phép chiếu ứng với $v$, $e$, $p$,
độ lớn thay đổi trên mười chiều và $L$ sau lần sửa (Bảng~\ref{bang:vidu-so}); các hiệu có dấu $x$, $c$, $z$ chỉ dùng
cho giờ. Không có ngưỡng: thay đổi 1 mili giây cũng tạo bản ghi sửa, nên số bản ghi sửa không phải số lần một ứng
dụng thật báo lại cho khách.
% EV: A04; X3-26 (DELTA:64-160); X3-25 (VALID:562-597); LEDGER:285-314; X3-13 (SPEC:35-46)

\begin{table}[htbp]
\centering
\caption[Bản ghi sổ của khách A trong ví dụ dựng tay]{Hai bản ghi giờ trả của khách A ở Hình~\ref{hinh:bamoc} (dựng tay), với giả định 9:10:00 là lời hứa đầu,
công bố lúc 8:50:00, và $B = 30$ giây. Bản ghi lời hứa đầu không có lời hứa trước nên sổ quy ước $v = e = p$. Các
độ lớn và $L$ tính bằng giây.}
\label{bang:vidu-so}
\small
\begin{tabular}{@{}lllllrrrr@{}}
\toprule
\textbf{Bản ghi} & \textbf{Lúc ghi} & $v$ & $e$ & $p$ & $|x|$ & $|c|$ & $|z|$ & $L$\\
\midrule
Lời hứa đầu & 8:50:00 & 9:10:00 & 9:10:00 & 9:10:00 & 0 & 0 & 0 & 0\\
Sửa (ghép khách M) & 9:00:00 & 9:10:00 & 9:10:40 & 9:11:10 & 40 & 30 & 70 & 30\\
\bottomrule
\end{tabular}
\end{table}
% EV: ví dụ dựng tay; v = e = p ở bản ghi đầu theo SPEC:46; L = 0 + |c| theo LEDGER:285-288

Trong một bản ghi mô phỏng của luật C trần 120 giây, một khách có $v$ = 17:43:46,044, $e$ = 17:43:46,283 và
$p$ = 17:44:01,203. Phần dời $|p - e| = 14{,}92$ giây làm phần trần còn lại bằng $120 - 14{,}92 = 105{,}08$ giây.
Cả 14 bản ghi sửa sau có $c = 0$ và $|x| < 0{,}25$ giây, nên phần trần đã dùng không tăng.
% EV: round26-trace-rider.txt, Q3 SELECTED RIDER, entries 1-15; R26-CH3-OPEN-QUESTIONS-CODEX.md §Q3; C-120-d20181117-s10-r1-w17/transcript-00.ndjson, final checkpoint ordinal 1716, request rq-c48d442603e9ce8c056b029d6d50c497f4d224b8af469945692690e92a75afb0

Sau khi chạy, độ trễ so với lời hứa đầu được chia ba phần: phần tính cho điều phối, phần được hồ sơ lịch sử giải
thích và phần còn lại,
\begin{equation}
a - p_0 \;=\; q \;+\; f \;+\; (a - D_0 - q), \qquad q = \sum_{j} c_j, \qquad f = D_0 - p_0 .
\label{eq:so}
\end{equation}
Cộng ba phần thì $q$ và $D_0$ triệt tiêu, nên đẳng thức luôn đúng. Trong đó, $a$ là giờ trả thực tế trong mô
phỏng; $q$ cộng các $c$ giữ nguyên dấu (dời muộn dương, dời sớm âm), khác với $\sum |c|$ của luật C, và dưới luật
C, $|q| \le \sum |c| \le B$; $f$ là phần dự báo đã cộng sẵn vào lời hứa đầu, bằng 0 khi không có dự báo.
Theo~\eqref{eq:dongnhat}, phần còn lại bằng $\sum_j x_j + \rho - f$ với $\rho = a - p_n$: thay đổi khi giữ nguyên
kế hoạch mà dự báo chưa tính, cộng độ lệch giữa giờ trả thực tế và lời hứa cuối $p_n$. Với khách A ở
Bảng~\ref{bang:vidu-so}, nếu $D_0$ = 9:10:30 và A tới đúng lời hứa cuối, $a$ = 9:11:10, thì độ trễ 70 giây gồm
$q = 30$, $f = 30$ và 10 giây còn lại, đúng bằng $\sum x + \rho - f = 40 + 0 - 30$.
% EV: X3-12 (FFA:84-97: q = Σ(p − e) có dấu, D0 = lời hứa hiển thị đầu); A12; A06; F09; ví dụ dựng tay X3-37 (b) (số học ở NOTES mục 4)

Chương~\ref{chuong:thucnghiem} cộng ba phần trên các khách có $a - p_0 > 60$ giây, với $p_0$ là lời hứa giờ trả
đầu không dự báo (không dùng $D_0$). Đây là phép tính theo mô hình, không kết luận ai có lỗi: ``phần tính cho điều
phối'' chỉ nói rằng mô hình gán phần này cho quyết định điều phối khi so với kế hoạch giữ nguyên. Phần được hồ sơ
lịch sử giải thích phụ thuộc dự báo đang dùng; phần còn lại là phần dự báo đó chưa giải thích được, không phải phần
``không thể biết trước''.
% EV: X3-12 (FFA:134-138: nhóm A − P0 > 60 s); F09; K04; A01 (lưu ý D03:222-224); K19

% -----------------------------------------------------------------------------------------------------
\section{Cài đặt, kiểm thử và tái lập}
\label{muc:caidat}

\subsection{Kiến trúc và trao đổi với bộ mô phỏng}
\label{muc:kientruc}

RideBound viết bằng C\# trên .NET 10, chia thành các dự án theo tầng (Bảng~\ref{bang:kientruc}). Tầng nghiệp vụ
không tham chiếu dự án hay gói thư viện ngoài nào, OR-Tools chỉ có trong cầu nối bộ giải, và một bộ kiểm thử kiến
trúc kiểm các quy tắc này; nhờ vậy phép kiểm trần, ba mốc và sổ không phụ thuộc bộ mô phỏng hay bộ giải.
% EV: X3-14 (các tệp .csproj: net10.0, ProjectReference, PackageReference; DependencyRuleTests.cs:8-31, 137-166)

\begin{table}[htbp]
\centering
\caption[Các thành phần chính của RideBound và bộ chuyển đổi FleetPy]{Các thành phần chính của RideBound và bộ chuyển đổi FleetPy}
\label{bang:kientruc}
\small
\begin{tabular}{@{}p{4.6cm}p{10.2cm}@{}}
\toprule
\textbf{Thành phần} & \textbf{Vai trò}\\
\midrule
Domain (tầng nghiệp vụ) & Yêu cầu và giai đoạn của khách, lời hứa, khóa, sổ, phép kiểm trần\\
Application (tầng ứng dụng) & Chiếu lịch theo thời gian đi, tính ba mốc giờ trả, kiểm phương án theo khóa và luật
lời hứa\\
Algorithms (thuật toán) & Bộ sinh chỉ chèn thêm, đánh giá phương án theo luật lời hứa và luật phục hồi, bộ chọn\\
Solvers.OrTools (cầu nối bộ giải) & Cầu nối tới bộ giải CP-SAT của OR-Tools 9.15\\
Contracts (định dạng bản tin) & Định dạng các bản tin trao đổi với bộ mô phỏng\\
Runner (bộ chạy) & Nhận bản tin, ghép các thành phần trên và trả quyết định\\
Benchmarking (chuẩn bị dữ liệu) & Chuẩn hóa dữ liệu Manhattan của FleetPy thành kịch bản thử\\
Bộ chuyển đổi FleetPy (Python) & Phần điều khiển đội xe của FleetPy 1.0.2: gửi sự kiện cho bộ chạy, áp quyết định
vào mô phỏng\\
\bottomrule
\end{tabular}
\end{table}
% EV: X3-14 (Google.OrTools 9.15.6755 chỉ ở RideBound.Solvers.OrTools.csproj:9); X3-22 (OrTools:3, CP-SAT); X3-15; NORM; B01

Bộ chuyển đổi viết bằng Python đóng vai phần điều khiển đội xe của FleetPy 1.0.2~\cite{engelhardt2026}, ghim theo
phiên bản; nó chạy bộ chạy RideBound như một chương trình con, hai bên nhắn tin qua đầu vào, đầu ra chuẩn, mỗi bản
tin một dòng JSON (Hình~\ref{hinh:bantin}). Bộ chạy chỉ nhận lô mới sau khi bộ mô phỏng báo đã áp dụng quyết định
trước, nên các quyết định được áp dụng đúng thứ tự. Bản tin của mỗi lần chạy được lưu, và mọi con số ở
Chương~\ref{chuong:thucnghiem} được tính từ đó.
% EV: B01; J18 (FLEET:474); X3-15 (README:1-12; runner_client.py:133-138; Program.cs:103-106; RunnerSession.cs:448-455); X3-13 (SPEC:28-37); A26

\begin{figure}[htbp]
\centering
\begin{tikzpicture}[y=1cm,font=\small,
  ben/.style={draw=xamvien,fill=xamnhat,rounded corners=2pt,align=center,inner sep=4pt,text width=4.3cm},
  tin/.style={-{Latex[length=2mm]},thick},
  nhan/.style={midway,above,font=\footnotesize,align=center}]
\node[ben] (L) at (0,0) {FleetPy 1.0.2\\và bộ chuyển đổi (Python)};
\node[ben] (R) at (9.5,0) {Bộ chạy RideBound\\(C\#, .NET 10)};
\draw[xamvien,thick] (L.south)--(0,-12.5);
\draw[xamvien,thick] (R.south)--(9.5,-12.5);
\draw[tin] (0,-1.8)--(9.5,-1.8) node[nhan]{chào hỏi; khởi tạo lần chạy\\(kịch bản, cấu hình luật)};
\draw[tin] (9.5,-2.9)--(0,-2.9) node[nhan]{đã nhận chào hỏi; đã khởi tạo};
\draw[dashed,xamvien] (-0.6,-3.35) rectangle (12.2,-9.1);
\node[anchor=north west,font=\footnotesize\itshape] at (-0.5,-3.45) {mỗi lượt quyết định};
\draw[tin] (0,-5.1)--(9.5,-5.1) node[nhan]{lô sự kiện: yêu cầu mới, xác nhận đặt xe,\\tiến độ xe, đón, trả,
thời gian đi mới};
\node[ben,text width=3.9cm,font=\footnotesize] at (9.5,-6.15) {sinh phương án, kiểm luật, chọn, ghi sổ
(Bảng~\ref{bang:luotquyetdinh})};
\draw[tin] (9.5,-7.7)--(0,-7.7) node[nhan]{quyết định: kế hoạch mới của các xe, lời hứa};
\draw[tin] (0,-8.7)--(9.5,-8.7) node[nhan]{đã áp dụng quyết định};
\draw[tin] (0,-10.15)--(9.5,-10.15) node[nhan]{yêu cầu trạng thái};
\draw[tin] (9.5,-11.15)--(0,-11.15) node[nhan]{trạng thái, gồm sổ theo từng khách};
\draw[tin] (0,-12.15)--(9.5,-12.15) node[nhan]{kết thúc: bộ chạy dừng, không trả lời};
\end{tikzpicture}
\caption[Các bản tin giữa bộ chuyển đổi FleetPy và bộ chạy RideBound]{Các bản tin chính giữa bộ chuyển đổi FleetPy và bộ chạy RideBound trong một lần chạy (đã giản lược). Mỗi
bản tin là một dòng JSON; phần trong khung nét đứt lặp lại ở mỗi lượt quyết định.}
\label{hinh:bantin}
\end{figure}
% EV: X3-15 (ProtocolPrimitives.cs:104-121; RunnerSession.cs:177-214: chỉ xử lý hello, initializeRun, eventBatch, decisionApplied, checkpoint, restore; shutdown làm bộ chạy dừng, không trả lời; finalizeRun/runSummary chỉ được khai báo; runner_client.py:158-273); X3-13 (SPEC:35-37)

\subsection{Giới hạn tính toán và thời gian chạy}
\label{muc:gioihan}

Ngoài giới hạn 2 yêu cầu mỗi lượt và 100 phương án mỗi xe, các bước sinh, kiểm và giải đều có giới hạn số bước
tính, và bộ giải dùng hạt giống ngẫu nhiên (số khởi tạo phần ngẫu nhiên) cố định bằng 7 để chạy lại cho cùng kết
quả. Vì vậy bộ sinh không xét hết mọi phương án; mục~\ref{muc:c4-hople} báo các bộ đếm của bộ sinh phương án cho biết
giới hạn có cắt bớt phương án hay không. Phép kiểm ở mỗi phương án duyệt mười chiều cho mỗi khách có lời hứa của xe và đếm đảo thứ tự
trên từng cặp điểm dừng, nên không thể coi khối lượng kiểm là tuyến tính theo độ dài kế hoạch; nó chưa được đo
riêng.
% EV: A24 (WP4:9-23); X3-17; X3-32 (GEN:54-79); N-25; X3-05 (BUDGET:76-99); X3-04 (VALID:233-239); X3-26 (DELTA:190-215)

Khóa luận chỉ báo thời gian thực của cả lần chạy: trên tập tắc đường cố định (16 lần chạy mỗi cấu hình), trung vị là
730, 667 và 641 giây với ghép tự do, điều phối không trần và luật C trần 30 giây. Thời gian này gồm mô phỏng FleetPy
và thời gian chờ RideBound trả quyết định, không phải thời gian của một lượt quyết định. Chi phí riêng của phép kiểm trần
không được báo, nên các số này không dùng để so chi phí tính toán.
% EV: H13 (EL17:2-6); X3-19; B17; R26-CH3-OPEN-QUESTIONS-CODEX.md §Q6; wp14r_scientific_protocol.py:801-803, 824-829; wp14r_supervised_process.py:297, 329-330, 979-985

\subsection{Kiểm tra kết quả}
\label{muc:kiemtra}

Kết quả được kiểm ở bốn mức.
\begin{itemize}[leftmargin=1.2cm,itemsep=2pt,topsep=3pt]
\item \textbf{Kiểm thử phần mềm:} chạy lại trên đúng mã nguồn đã dựng bộ chạy, cả 1.025 trong 1.025 kiểm thử tự động
đạt trong một lần chạy đầy đủ ngày 13/10/2026, với kiểu xuống dòng LF và một tệp chưa ghi nhận trong cây làm việc;
chi tiết ở Phụ lục~\ref{phuluc:C}.
\item \textbf{Kiểm tính nhất quán của sổ:} ở cả 382.006 bản ghi sửa của 185 lần chạy (170 lần chạy của năm cấu hình
trên tập kiểm tra chính và 15 lần chạy luật C trần 30 giây trên 15 kịch bản hiệu chỉnh), $|x|$, $|c|$, $|z|$ tính
lại từ $v$, $e$, $p$ trong sổ khớp đúng các trường độ lớn. Phép kiểm này không kiểm $e$ có được tính đúng hay không.
\item \textbf{Tính lại độc lập:} một chương trình viết riêng, độc lập về mã nhưng dùng cùng đặc tả chỉ số, tính lại
13 hạng mục số liệu từ bản ghi thô và một điều kiện luôn phải đúng, không lệch hạng mục nào; các số đếm chính được
ít nhất hai bộ đọc viết độc lập tính lại và khớp trong phạm vi đã kiểm. Hai mức kiểm này được làm bằng các chương
trình đọc riêng, không phải kiểm thủ công; chúng cho biết con số khớp bản ghi, không cho biết mô phỏng giống thực tế
đến đâu.
\item \textbf{Tái lập:} tập kiểm tra chính, tập chỉnh trần buổi sáng, các lần chạy trần 120 giây và tám lần chạy chiều
tuyết dùng cùng một bản bộ chạy đã biên dịch và lưu lại; tập tắc đường cố định dùng một bản cũ hơn. Bản dựng lại, nếu có,
phải đối chiếu với kết quả đã lưu trước khi dùng. Các lượt ưu tiên chi phí dùng bản riêng; bốn cặp kiểm
hành vi không thay thế việc kiểm đồng nhất toàn bộ hai bản chạy (Phụ lục~\ref{phuluc:D}).
\end{itemize}
% EV (kiểm thử): TESTRUN-45f8026-2026-10-12.md (mã nguồn bộ chạy; lượt 3 và 3b; điều kiện LF và tệp chưa ghi nhận); A28
% EV (sổ): A26 (T38:106; T38I:68); B10 (170 + 15); RUN38:1-8 (15 lần chạy C-30 trên 15 kịch bản hiệu chỉnh)
% EV (tính lại): C19 (T38I:79-97: 13 hạng mục + 1 bất biến); ADDENDUM-1:15-19; C04; E01; K41; A26 lưu ý; I01
% EV (tái lập): A28 (RUNW07:40-41); X3-10 (run_w07.py:8; run_w07x.py:7; run_t38c120.py:3-8); B17; NEW:13; I13; R26-CH3-OPEN-QUESTIONS-CODEX.md §Q8 (bản bộ chạy chiều tuyết)

Chương~\ref{chuong:thucnghiem} dùng các thành phần trên để đo luật C được gì và mất gì trên mô phỏng FleetPy, so
với điều phối không trần, luật hạn chót và cách đệm lời hứa, tức cộng thêm một khoảng cố định vào giờ hứa cho mọi
khách.
